% Leçon 36 — Grammaire : antonyme 新 / 旧 ; verbe modal 可以 — Chinois 1ère A — Cours signé OnBuch+
% Programme officiel MINESEC. Compiler avec : tectonic ZHA36-grammaire-antonyme-verbe-modal.tex
\newcommand{\DOCMATIERE}{Chinois}
\newcommand{\DOCNIVEAU}{1ère A}
\newcommand{\DOCMODULE}{Module 5 : Mass médias et télécommunication}
\newcommand{\DOCLECON}{ZHA36}
\newcommand{\DOCTITRE}{Leçon 36 — Grammaire : antonyme 新 / 旧 ; verbe modal 可以}
% OnBuch+ — préambule « cours signés » ENRICHI (compilé avec tectonic / XeLaTeX)
% Système visuel « Pop » d'OnBuch+ : fond crème (jamais blanc), bordures encre
% épaisses, ombres « dures » décalées, titres Archivo Black en capitales.
% Version MATHÉMATIQUES : graphes (pgfplots/tikz), boîtes pédagogiques
% (définition, propriété, méthode, attention, le savais-tu, exercice résolu,
% exercice, TP, bilan), page de garde et signature OnBuch+ ancrée en bas.
%
% Macros attendues dans main.tex AVANT \input{preamble} :
%   \DOCMATIERE  \DOCNIVEAU  \DOCMODULE  \DOCLECON (numéro)  \DOCTITRE
\documentclass[11pt]{article}

\usepackage{fontspec}
\usepackage[french]{babel} % français = langue principale ; le chinois via \zh{..} / textebox (xeCJK)
\usepackage{xeCJK}
\usepackage{pifont} % \ding{51} \ding{55} utilisés par les modèles
\usepackage[a4paper,margin=2cm,top=3.4cm,bottom=2.8cm,headheight=1.4cm,footskip=1.3cm]{geometry}
\usepackage{amsmath,amssymb,amsfonts}
\usepackage{mathtools} % \xmapsto, \coloneqq... souvent utilisés par les modèles
\usepackage{cancel} % simplifications barrées (bilans de forces)
% \abs{z} (module, valeur absolue) et \norm{u} (norme), utilisés par les modèles
\providecommand{\abs}{}\let\abs\relax\DeclarePairedDelimiter{\abs}{\lvert}{\rvert}
\providecommand{\norm}{}\let\norm\relax\DeclarePairedDelimiter{\norm}{\lVert}{\rVert}
\usepackage[table]{xcolor} % [table] : fournit aussi \rowcolors (alternance de lignes)
\usepackage{pagecolor}
\usepackage{tikz}
\usetikzlibrary{arrows.meta,positioning,calc,shapes.geometric,shapes.misc,shapes.arrows,decorations.pathmorphing,decorations.pathreplacing,decorations.markings,patterns,shadows,mindmap,trees,fit,backgrounds,matrix,intersections,angles,quotes}
\usepackage{pgfplots}
\usepackage[version=4]{mhchem} % équations nucléaires \ce{^{235}_{92}U}
\usepackage[european,siunitx]{circuitikz} % schémas électriques
\pgfplotsset{compat=1.18}
\usepgfplotslibrary{fillbetween,groupplots} % aires entre courbes (calcul intégral)
\usepackage{enumitem}
\usepackage{fancyhdr}
% [most] charge toutes les bibliothèques tcolorbox (listings, minted, raster,
% documentation...) et gonfle inutilement la mémoire TeX sur un document long
% (~300 boîtes) : on ne charge que ce qui est réellement utilisé.
\usepackage[skins,breakable]{tcolorbox}
\usepackage{graphicx}
\usepackage{colortbl}
\usepackage{booktabs}
\usepackage{tabularx}
\usepackage{diagbox} % case d'en-tête diagonale des tableaux à double entrée
% Colonnes extensibles centrée / gauche / droite (C, L, R), souvent utilisées par les modèles
\newcolumntype{C}{>{\centering\arraybackslash}X}
\newcolumntype{L}{>{\raggedright\arraybackslash}X}
\newcolumntype{R}{>{\raggedleft\arraybackslash}X}
\usepackage{array}
\usepackage{multicol}
\usepackage{multirow}
\usepackage{siunitx}
% parse-numbers=false : accepte aussi \SI{2,5 \times 10^{-3}}{...}
\sisetup{locale=FR,output-decimal-marker={,},group-minimum-digits=4,parse-numbers=false}
\DeclareSIUnit\ohm{\ensuremath{\symup{\Omega}}} % U+2126 absent de la police
\DeclareSIUnit\division{div} % réglages d'oscilloscope (ms/div, V/div)
\DeclareSIUnit\tr{tr} % tours (vitesse de rotation, ex. \SI{1500}{\tr\per\minute})
\DeclareSIUnit\molar{M} % concentration molaire (ex. \SI{3}{\molar}, \SI{1}{\micro\molar})
\DeclareSIUnit\an{an} % année (le modèle confond parfois avec \year, un compteur TeX) — ex. \SI{5}{\kilo\meter\per\an}
\DeclareSIUnit\FCFA{FCFA} % franc CFA (ex. \SI{500}{\FCFA\per\kilo\gram})
\DeclareSIUnit\degreeBrix{\textdegree Brix} % teneur en sucre d'un sirop/jus (réfractométrie)
\DeclareSIUnit\month{mois} % le modèle utilise parfois \month, un compteur TeX, comme unité de durée
\usepackage{qrcode}
% \degree : commande fréquemment utilisée par les modèles pour un angle ou une
% température en dehors de \SI{}{} (non fournie par les paquets chargés).
\providecommand{\degree}{^{\circ}}
% \qed (fin de démonstration), utilisé par les modèles sans amsthm
\providecommand{\qed}{\hfill$\square$}
% \barre : double barre des valeurs interdites dans un tableau de variations (tkz-tab)
\providecommand{\barre}{\ensuremath{\|}}
% environnement proof (amsthm non chargé) utilisé par les modèles
\ifdefined\proof\else\newenvironment{proof}[1][Démonstration]{\par\noindent\textit{#1.}\ }{\qed\par}\fi
\usepackage{lastpage}
\usepackage{hyperref}
\hypersetup{hidelinks}

% ============================================================
% Palette « Pop » OnBuch+ — mêmes hex que la classe P (plus_ui.dart)
% ============================================================
\definecolor{poporange}{HTML}{F59321}
\definecolor{popdark}{HTML}{E07A0C}
\definecolor{popcream}{HTML}{FFF6E8}
\definecolor{popink}{HTML}{1C1714}
\definecolor{poppurple}{HTML}{7A5AE0}
\definecolor{popgreen}{HTML}{2E9E6A}
\definecolor{poppink}{HTML}{E0457B}
\definecolor{popblue}{HTML}{2F7DE1}
\definecolor{popgold}{HTML}{C98A12}
\definecolor{popmuted}{HTML}{8A7F76}
\definecolor{popline}{HTML}{EADFD2}
\definecolor{popgray}{HTML}{8A7F76}
\colorlet{popgrey}{popgray}
% Alias tolérants (le modèle nomme parfois les couleurs différemment)
\colorlet{popwhite}{white}
\colorlet{popblack}{popink}
\colorlet{popred}{poppink}
\colorlet{popyellow}{popgold}
% Teintes claires (fonds de tableaux/encadrés) — le modèle nomme parfois
% les couleurs avec un suffixe L (ex. popblueL) sans les avoir définies.
\colorlet{poporangeL}{poporange!20}
\colorlet{popdarkL}{popdark!20}
\colorlet{poppurpleL}{poppurple!20}
\colorlet{popgreenL}{popgreen!20}
\colorlet{poppinkL}{poppink!20}
\colorlet{popblueL}{popblue!20}
\colorlet{popgoldL}{popgold!20}
\colorlet{popredL}{popred!20}
\colorlet{popgrayL}{popgray!20}
% Teintes claires pour fonds de tableaux / zones
\colorlet{popblueL}{popblue!10!white}
\colorlet{poporangeL}{poporange!14!white}
\colorlet{popgreenL}{popgreen!12!white}
\colorlet{poppurpleL}{poppurple!12!white}
\colorlet{poppinkL}{poppink!12!white}
\colorlet{popgoldL}{popgold!14!white}

\pagecolor{popcream}

% ============================================================
% Typographie
% ============================================================
\defaultfontfeatures{Ligatures=TeX}
\setmainfont{PlusJakartaSans-Medium.ttf}[Path=../../../fonts/,BoldFont=PlusJakartaSans-Bold.ttf]
% Symboles, API et emojis absents de la police principale : repli sur DejaVu Sans (caractères actifs)
\newfontface\symfont{DejaVuSans.ttf}[Path=../../../fonts/]
\makeatletter
\newcommand\symactive[1]{\catcode#1=13 \begingroup\lccode`\~=#1 \lowercase{\endgroup\def~}{{\symfont\char#1}}}
\newcommand\symblank[1]{\catcode#1=13 \begingroup\lccode`\~=#1 \lowercase{\endgroup\def~}{}}
\newcommand\symhyph[1]{\catcode#1=13 \begingroup\lccode`\~=#1 \lowercase{\endgroup\def~}{\mbox{-}}}
\newcount\symc
\newcommand\symrange[2]{\symc=#1 \@whilenum\symc<#2 \do{\expandafter\symactive\expandafter{\the\symc}\advance\symc by 1 }}
\newcommand\symblankrange[2]{\symc=#1 \@whilenum\symc<#2 \do{\expandafter\symblank\expandafter{\the\symc}\advance\symc by 1 }}
\makeatother
\symrange{"0250}{"02FF}\symactive{"2713}\symactive{"2714}\symactive{"2717}\symactive{"2718}\symactive{"2610}\symactive{"2611}\symactive{"2612}
\symrange{"2600}{"27BF}
\catcode"2705=13 \begingroup\lccode`\~="2705 \lowercase{\endgroup\def~}{{\symfont\char"2713}}\symblank{"26BD}\symblank{"26A1}
\symrange{"2460}{"257F}\symactive{"223C}\symactive{"2248}\symactive{"2260}
\symactive{"01F8}\symactive{"01F9}\symactive{"1E3E}\symactive{"1E3F}
\symhyph{"2011}\symhyph{"2E1A}\symblankrange{"1F300}{"1FAFF}\symblank{"FE0F}
% Caractères chinois : WenQuanYi Zen Hei (sous-ensemble GB2312 + pinyin), gras/italique simulés
\setCJKmainfont{WenQuanYiZenHei-subset.ttf}[Path=../../../fonts/,AutoFakeBold=3,AutoFakeSlant=0.2]
\xeCJKsetup{CJKspace=false}
% Pinyin : voyelles à caron (ǎ ǐ ǒ ǔ ǖ ǘ ǚ ǜ) absentes de la police principale -> caractères actifs
\newfontface\pyfont{WenQuanYiZenHei-subset.ttf}[Path=../../../fonts/]
\newcommand\pyactive[1]{\catcode#1=13 \begingroup\lccode`\~=#1 \lowercase{\endgroup\def~}{{\pyfont\char#1}}}
\pyactive{"01CD}\pyactive{"01CE}\pyactive{"01CF}\pyactive{"01D0}\pyactive{"01D1}\pyactive{"01D2}\pyactive{"01D3}\pyactive{"01D4}\pyactive{"01D5}\pyactive{"01D6}\pyactive{"01D7}\pyactive{"01D8}\pyactive{"01D9}\pyactive{"01DA}\pyactive{"01DB}\pyactive{"01DC}
% Maths en sans-serif (Fira Math) avec chiffres et lettres droites de Plus
% Jakarta, pour que les formules se fondent dans le texte.
\usepackage[math-style=ISO,bold-style=ISO]{unicode-math}
\setmathfont{FiraMath-Regular.otf}[Path=../../../fonts/]
\setmathfont{PlusJakartaSans-Medium.ttf}[Path=../../../fonts/,range={up/num,up/{Latin,latin},\mathup}]
\usepackage{icomma} % virgule décimale sans espace en mode math
% Caractères mathématiques Unicode absents des polices de texte (Plus Jakarta,
% Archivo Black) : rendus par leur équivalent mathématique, en texte comme en
% formule — sinon ils s'affichent en carré vide (ex. « Arithmétique dans ℤ »).
\usepackage{newunicodechar}
\newunicodechar{ℝ}{\ensuremath{\mathbb{R}}}
\newunicodechar{ℤ}{\ensuremath{\mathbb{Z}}}
\newunicodechar{ℂ}{\ensuremath{\mathbb{C}}}
\newunicodechar{ℕ}{\ensuremath{\mathbb{N}}}
\newunicodechar{ℚ}{\ensuremath{\mathbb{Q}}}
\newunicodechar{𝔻}{\ensuremath{\mathbb{D}}}
\newunicodechar{α}{\ensuremath{\alpha}}
\newunicodechar{β}{\ensuremath{\beta}}
\newunicodechar{γ}{\ensuremath{\gamma}}
\newunicodechar{δ}{\ensuremath{\delta}}
\newunicodechar{ε}{\ensuremath{\varepsilon}}
\newunicodechar{θ}{\ensuremath{\theta}}
\newunicodechar{λ}{\ensuremath{\lambda}}
\newunicodechar{μ}{\ensuremath{\mu}}
\newunicodechar{σ}{\ensuremath{\sigma}}
\newunicodechar{τ}{\ensuremath{\tau}}
\newunicodechar{φ}{\ensuremath{\varphi}}
\newunicodechar{ψ}{\ensuremath{\psi}}
\newunicodechar{ω}{\ensuremath{\omega}}
\newunicodechar{Σ}{\ensuremath{\Sigma}}
% majuscules produites par \MakeUppercase dans les titres (α-aminés -> Α)
\newunicodechar{Α}{\ensuremath{\alpha}}
\newunicodechar{Β}{\ensuremath{\beta}}
\newunicodechar{Γ}{\ensuremath{\Gamma}}
\newunicodechar{Δ}{\ensuremath{\Delta}}
\newunicodechar{Ε}{\ensuremath{\varepsilon}}
\newunicodechar{Θ}{\ensuremath{\Theta}}
\newunicodechar{Λ}{\ensuremath{\Lambda}}
\newunicodechar{Μ}{\ensuremath{\mu}}
\newunicodechar{Τ}{\ensuremath{\tau}}
\newunicodechar{Φ}{\ensuremath{\Phi}}
\newunicodechar{Ψ}{\ensuremath{\Psi}}
\newunicodechar{Ω}{\ensuremath{\Omega}}
\newunicodechar{↦}{\ensuremath{\mapsto}}
\newunicodechar{∈}{\ensuremath{\in}}
\newunicodechar{∉}{\ensuremath{\notin}}
\newunicodechar{∧}{\ensuremath{\wedge}}
\newunicodechar{⇔}{\ensuremath{\Leftrightarrow}}
\newunicodechar{⟺}{\ensuremath{\Longleftrightarrow}}
\newunicodechar{⇒}{\ensuremath{\Rightarrow}}
\newunicodechar{⟹}{\ensuremath{\Longrightarrow}}
\newunicodechar{∘}{\ensuremath{\circ}}
\newunicodechar{∪}{\ensuremath{\cup}}
\newunicodechar{∩}{\ensuremath{\cap}}
\newunicodechar{≡}{\ensuremath{\equiv}}
\newunicodechar{⊂}{\ensuremath{\subset}}
\newunicodechar{⊥}{\ensuremath{\perp}}
\newunicodechar{∃}{\ensuremath{\exists}}
\newunicodechar{∀}{\ensuremath{\forall}}
\newunicodechar{∛}{\ensuremath{\sqrt[3]{\,}}}
\newunicodechar{∜}{\ensuremath{\sqrt[4]{\,}}}
\newunicodechar{⃗}{\ensuremath{{}^{\to}}}
\newunicodechar{ⁿ}{\ifmmode^{n}\else\textsuperscript{\ensuremath{n}}\fi}
\newunicodechar{ˣ}{\ifmmode^{x}\else\textsuperscript{\ensuremath{x}}\fi}
\newunicodechar{ᵇ}{\ifmmode^{b}\else\textsuperscript{\ensuremath{b}}\fi}
\newunicodechar{ʳ}{\ifmmode^{r}\else\textsuperscript{\ensuremath{r}}\fi}
\newunicodechar{ᵘ}{\ifmmode^{u}\else\textsuperscript{\ensuremath{u}}\fi}
\newunicodechar{ʸ}{\ifmmode^{y}\else\textsuperscript{\ensuremath{y}}\fi}
\newunicodechar{ᵃ}{\ifmmode^{a}\else\textsuperscript{\ensuremath{a}}\fi}
\newunicodechar{ᵉ}{\ifmmode^{e}\else\textsuperscript{\ensuremath{e}}\fi}
\newunicodechar{ᵏ}{\ifmmode^{k}\else\textsuperscript{\ensuremath{k}}\fi}
\newunicodechar{ᵖ}{\ifmmode^{p}\else\textsuperscript{\ensuremath{p}}\fi}
\newunicodechar{ᴺ}{\ifmmode^{N}\else\textsuperscript{\ensuremath{N}}\fi}
\newunicodechar{ᶜ}{\ifmmode^{c}\else\textsuperscript{\ensuremath{c}}\fi}
\newunicodechar{ʰ}{\ifmmode^{h}\else\textsuperscript{\ensuremath{h}}\fi}
\newunicodechar{ᵠ}{\ifmmode^{q}\else\textsuperscript{\ensuremath{q}}\fi}
\newunicodechar{ₙ}{\ifmmode_{n}\else\textsubscript{\ensuremath{n}}\fi}
\newunicodechar{ₐ}{\ifmmode_{a}\else\textsubscript{\ensuremath{a}}\fi}
\newunicodechar{ᵢ}{\ifmmode_{i}\else\textsubscript{\ensuremath{i}}\fi}
\newunicodechar{ᵣ}{\ifmmode_{r}\else\textsubscript{\ensuremath{r}}\fi}
\newunicodechar{ₖ}{\ifmmode_{k}\else\textsubscript{\ensuremath{k}}\fi}
\newunicodechar{ᵦ}{\ifmmode_{\beta}\else\textsubscript{\ensuremath{\beta}}\fi}
\newunicodechar{ₓ}{\ifmmode_{x}\else\textsubscript{\ensuremath{x}}\fi}
\newfontfamily\popdisplay{ArchivoBlack-Regular.ttf}[Path=../../../fonts/]
\newfontfamily\poplogo{SpaceGrotesk-Bold.ttf}[Path=../../../fonts/]
\newfontfamily\popsemi{PlusJakartaSans-SemiBold.ttf}[Path=../../../fonts/]
\newfontfamily\popxbold{PlusJakartaSans-ExtraBold.ttf}[Path=../../../fonts/]
\newfontfamily\popxboldls{PlusJakartaSans-ExtraBold.ttf}[Path=../../../fonts/,LetterSpace=3.0]

\setlist[enumerate]{leftmargin=1.7em,itemsep=5pt,topsep=4pt}
\setlist[itemize]{leftmargin=1.7em,itemsep=4pt,topsep=4pt,label={\color{poporange}\textbullet}}
\setlength{\parindent}{0pt}
\setlength{\parskip}{5pt}
\renewcommand{\arraystretch}{1.25}

% pgfplots : style Pop par défaut
\pgfplotsset{
  every axis/.append style={axis line style={popink,line width=1pt},tick style={popink},
    grid style={popline,line width=0.5pt},label style={font=\small},tick label style={font=\footnotesize}},
  popcurve/.style={poporange,line width=1.8pt,no markers,smooth},
}
% Même style disponible dans un \draw TikZ simple (hors pgfplots)
\tikzset{popcurve/.style={poporange,line width=1.8pt}}
\tikzset{popgrid/.style={popline,very thin}}
\colorlet{popcurve}{poporange} % le nom du style est parfois utilisé comme couleur

% ============================================================
% Pill / squiggle
% ============================================================
\newcommand{\poppill}[3]{%
  \tikz[baseline=-0.55ex]\node[draw=popink,line width=1.2pt,rounded corners=2.6pt,%
    fill=#2,inner xsep=6.5pt,inner ysep=3pt]{%
    \popxboldls\fontsize{7}{7}\selectfont\color{#3}\MakeUppercase{#1}};%
}
\newcommand{\popsquiggle}[1]{%
  \tikz[baseline=-0.4ex]{
    \draw[#1,line width=2.6pt,line cap=round]
      plot [smooth] coordinates {(0,0.09) (0.34,0.01) (0.68,0.21) (1.02,0.01) (1.36,0.10)};
  }%
}
% Mot-clé surligné (marqueur orange sous le texte)
\newcommand{\cle}[1]{{\popxbold\color{popdark}#1}}

% ============================================================
% En-tête / pied
% ============================================================
\pagestyle{fancy}
\fancyhf{}
\renewcommand{\headrulewidth}{0pt}
\renewcommand{\footrulewidth}{0pt}
\lhead{\raisebox{-0.15ex}{\poplogo\fontsize{13}{13}\selectfont\color{popink}OnBuch\color{poporange}+}}
\rhead{\poppill{\DOCMATIERE\ \textbullet\ \DOCNIVEAU\ \textbullet\ Leçon \DOCLECON}{white}{popink}}
\lfoot{\popxbold\fontsize{7.6}{7.6}\selectfont\color{popmuted}\MakeUppercase{Cours signé OnBuch+}\ \textbullet\ {\popsemi\DOCTITRE}}
\rfoot{\tikz[baseline=-0.6ex]\node[rounded corners=8.5pt,fill=popink,minimum height=17pt,minimum width=17pt,inner xsep=5pt,inner sep=0]{\popxbold\fontsize{8}{8}\selectfont\color{popcream}\thepage\,/\,\pageref*{LastPage}};}

% ============================================================
% Titres
% ============================================================
\newcommand{\chaptitre}[2]{%
  \begin{tcolorbox}[enhanced,colback=poporange,colframe=popink,boxrule=2.6pt,arc=11pt,
      drop shadow=popink,left=15pt,right=15pt,top=12pt,bottom=11pt,before skip=3pt,after skip=18pt]
    \poppill{#2}{popink}{popcream}\par\vspace{9pt}
    {\raggedright\hyphenpenalty=10000\exhyphenpenalty=10000\popdisplay\fontsize{22}{24}\selectfont\color{popcream}\MakeUppercase{#1}\par}
  \end{tcolorbox}
}
% Section numérotée : pastille orange avec le numéro + titre Archivo
\newcounter{popsec}
\newcommand{\coursec}[1]{%
  \stepcounter{popsec}\setcounter{popsub}{0}%
  \par\vspace{16pt}\noindent
  \begin{minipage}{\linewidth}
  \tikz[baseline=-0.7ex]\node[draw=popink,line width=1.6pt,fill=poporange,rounded corners=4pt,minimum size=22pt,inner sep=0]{\popdisplay\fontsize{12}{12}\selectfont\color{popcream}\thepopsec};%
  \hspace{8pt}{\popdisplay\fontsize{15}{17}\selectfont\color{popink}\MakeUppercase{#1}}\par
  \vspace{4pt}\noindent{\color{poporange}\rule{36pt}{3.6pt}}
  \end{minipage}\par\nopagebreak\vspace{8pt}
}
\newcounter{popsub}
\newcommand{\courssub}[1]{%
  \stepcounter{popsub}%
  \par\vspace{9pt}\noindent{\popxbold\fontsize{11.5}{13}\selectfont\color{popink}{\color{poporange}\thepopsec.\thepopsub}\hspace{6pt}#1}\par\nopagebreak\vspace{4pt}
}

% ============================================================
% Boîtes pédagogiques — même recette : fond blanc, bordure épaisse colorée,
% ombre dure encre, badge plein. Titre optionnel [#1].
% ============================================================
\tcbset{popbase/.style={breakable,enhanced,colback=white,boxrule=2.3pt,arc=9pt,
  shadow={3pt}{-3pt}{0pt}{popink},left=14pt,right=14pt,top=11pt,bottom=11pt,before skip=11pt,after skip=15pt,
  fontupper=\popsemi\fontsize{10.6}{14.5}\selectfont\color{popink},
  fontlower=\popsemi\fontsize{10.6}{14.5}\selectfont\color{popink}}}
% Coche dessinée (le glyphe ✓ n'existe pas dans les polices du document)
\newcommand{\popcheck}[1]{\tikz[baseline=-0.5ex]\draw[#1,line width=1.6pt,line cap=round,line join=round](-0.13,0.0)--(-0.04,-0.09)--(0.13,0.1);}
\providecommand{\coche}{\popcheck{popgreen}}
\providecommand{\resultat}[1]{\par\smallskip\noindent\fcolorbox{popgreen}{popgreenL}{\parbox{\dimexpr\linewidth-2\fboxsep-2\fboxrule\relax}{\textbf{Résultat.} #1}}\par\smallskip}
\newcommand{\popboxhead}[3]{\poppill{#1}{#2}{white}\ifx\relax#3\relax\else\hspace{6pt}{\popxbold\fontsize{10.6}{12}\selectfont\color{popink}#3}\fi\par\vspace{7pt}}

\newtcolorbox{aretenir}[1][]{popbase,colframe=popgreen,before upper={\popboxhead{À retenir}{popgreen}{#1}}}
% Texte en chinois (document de lecture, dialogue, extrait) : cadre
% d'étude. Un titre optionnel (source, auteur) s'affiche dans l'en-tête.
\newtcolorbox{textebox}[1][]{popbase,colframe=popblue,before upper={\popboxhead{文本 · Texte}{popblue}{#1}},after upper={}}
% Marques ✓ / ✗ (forme correcte / fautive) souvent utilisées par les modèles
\providecommand{\cross}{\textcolor{poppink}{\ensuremath{\times}}}
\providecommand{\xmark}{\cross}\providecommand{\crossmark}{\cross}\providecommand{\cmark}{\ensuremath{\checkmark}}
% Ligne de réponse / justification à compléter (inventée par les modèles)
\providecommand{\justification}[1]{\par\noindent\textit{Justification~:} \dotfill\par}
% \textipa (package tipa non chargé) : la police ne couvre pas l'API -> simple passage
\providecommand{\textipa}[1]{#1}
% Soulignement (ulem non chargé) : \uline, \ul
\providecommand{\uline}[1]{\underline{#1}}\providecommand{\ul}[1]{\underline{#1}}
% \glossaire{\item ...} inventée par les modèles : liste de glossaire
\providecommand{\glossaire}[1]{\begin{itemize}#1\end{itemize}}
% \alert (beamer) inventée par les modèles : texte mis en évidence
\providecommand{\alert}[1]{\textbf{\textcolor{poppink}{#1}}}
% variantes ASCII d'ordinaux écrites en macro
\providecommand{\iere}{\textsuperscript{re}}\providecommand{\ieme}{\textsuperscript{e}}\providecommand{\ere}{\textsuperscript{re}}\providecommand{\eme}{\textsuperscript{e}}\providecommand{\er}{\textsuperscript{er}}\providecommand{\ie}{\textsuperscript{e}}
% Ordinaux français écrits en macro par les modèles : 1\ière, 3\ième
\newcommand{\ière}{\textsuperscript{re}}\newcommand{\ième}{\textsuperscript{e}}
% \exemple (inventée par les modèles) : étiquette « Exemple : »
\providecommand{\exemple}{\textbf{Exemple~:}\ }
% \square utilisable aussi hors mode math (cases à cocher)
\AtBeginDocument{\let\squareMath\square\renewcommand{\square}{\ensuremath{\squareMath}}}
% Mot ou phrase en chinois à l'intérieur d'un paragraphe français
\newcommand{\zh}[1]{\ifmmode\text{#1}\else{#1}\fi}
\let\eng\zh \let\esp\zh \let\all\zh % alias tolérés
% flèches d'intonation inventées par les modèles
\providecommand{\textsearrow}{}\renewcommand{\textsearrow}{\ensuremath{\searrow}}\providecommand{\textnearrow}{}\renewcommand{\textnearrow}{\ensuremath{\nearrow}}
% \fr{..} : mot français (inventé par les modèles)
\providecommand{\fr}[1]{\textit{#1}}
% zhbox : encadré de texte chinois inventé par les modèles
\newenvironment{zhbox}{\begin{quote}}{\end{quote}}
% \courssubsub : titre de niveau 3 inventé par les modèles
\providecommand{\courssubsub}[1]{\par\medskip\noindent\textbf{#1}\par\smallskip}
% \zhbf : texte chinois en gras inventé par les modèles
\providecommand{\zhbf}[1]{\textbf{#1}}
% \vs : « versus » inventé par les modèles
\providecommand{\vs}{\textit{vs}\ }
% \blank : case à compléter inventée par les modèles
\providecommand{\blank}[1][]{\underline{\hspace{1.6em}}}
\newtcolorbox{exemplebox}[1][]{popbase,colframe=poporange,before upper={\popboxhead{Exemple}{poporange}{#1}}}
\newtcolorbox{definition}[1][]{popbase,colframe=popblue,before upper={\popboxhead{Définition}{popblue}{#1}}}
\newtcolorbox{propriete}[1][]{popbase,colframe=poppurple,before upper={\popboxhead{Propriété}{poppurple}{#1}}}
\newtcolorbox{methode}[1][]{popbase,colframe=popgold,before upper={\popboxhead{Méthode}{popgold}{#1}}}
\newtcolorbox{attention}[1][]{popbase,colframe=poppink,before upper={\popboxhead{Attention}{poppink}{#1}}}
\newtcolorbox{experience}[1][]{popbase,colframe=popblue,colback=popblueL,before upper={\popboxhead{Expérience}{popblue}{#1}}}
% « Le savais-tu ? » : carte violette pleine (contexte, culture, Cameroun)
\newtcolorbox{savaistu}[1][]{popbase,colback=poppurple,colframe=popink,
  fontupper=\popsemi\fontsize{10.4}{14.2}\selectfont\color{white},
  before upper={\poppill{Le savais-tu\,?}{white}{poppurple}\ifx\relax#1\relax\else\hspace{6pt}{\popxbold\color{white}#1}\fi\par\vspace{7pt}}}
% Exercice résolu : énoncé en haut, \tcblower puis la solution sur fond crème
\newcounter{popexo}\newcounter{popexr}
\newtcolorbox{exoresolu}[1][]{popbase,colframe=poporange,colbacklower=popcream,
  segmentation style={popink,line width=1.2pt,dashed},
  before upper={\stepcounter{popexr}\popboxhead{Exercice résolu \thepopexr}{poporange}{#1}},
  before lower={\poppill{Solution}{popink}{popcream}\par\vspace{6pt}}}
% Exercice (non résolu en place) — la correction est dans la partie Corrigés
\newtcolorbox{exercice}[1][]{popbase,colframe=popink,
  before upper={\stepcounter{popexo}\popboxhead{Exercice \thepopexo}{popink}{#1}}}
% Corrigé
\newtcolorbox{corrige}[1][]{popbase,colframe=popgreen,colback=popgreenL,
  before upper={\popboxhead{Corrigé}{popgreen}{#1}}}
% Objectifs / prérequis / bilan
\newtcolorbox{objectifs}{popbase,colframe=popink,colback=poporangeL,
  before upper={\popboxhead{Objectifs de la leçon}{popink}{}}}
\newtcolorbox{prerequis}{popbase,colframe=popink,colback=popblueL,
  before upper={\popboxhead{Prérequis}{popblue}{}}}
\newtcolorbox{bilan}{popbase,colframe=popink,colback=white,boxrule=2.8pt,
  before upper={\popboxhead{Fiche bilan}{poporange}{L'essentiel à connaître pour l'examen}}}
% Figure encadrée avec légende
\newtcolorbox{popfigure}[1][]{enhanced,breakable=false,colback=white,colframe=popline,boxrule=1.4pt,arc=8pt,
  left=10pt,right=10pt,top=10pt,bottom=8pt,before skip=10pt,after skip=12pt,halign=center,
  lower separated=false,halign lower=center,
  fontlower=\popsemi\fontsize{9}{11.5}\selectfont\color{popmuted},
  #1}
\newcommand{\legende}[1]{\par\vspace{6pt}{\popsemi\fontsize{9}{11.5}\selectfont\color{popmuted}#1}}

% ============================================================
% Signature OnBuch+
% ============================================================
\newcommand{\poplogoinline}[1]{{\poplogo\fontsize{#1}{#1}\selectfont\color{popink}OnBuch\color{poporange}+}}
\newcommand{\signatureonbuch}{%
  \begin{tcolorbox}[enhanced,colback=popink,colframe=popink,boxrule=2.2pt,arc=10pt,
      drop shadow=poporange,left=14pt,right=14pt,top=10pt,bottom=10pt,before skip=0pt,after skip=0pt]
    \tikz[baseline=-0.7ex]\node[circle,fill=popgreen,draw=popcream,line width=1.3pt,minimum size=22pt,inner sep=0]{\popcheck{white}};%
    \hspace{9pt}\begin{minipage}[c]{0.62\linewidth}
      {\popxbold\fontsize{12}{13}\selectfont\color{popcream}Signé\ \textbullet\ {\poplogo OnBuch\color{poporange}+}}\par\vspace{2pt}
      {\raggedright\popsemi\fontsize{8}{10.5}\selectfont\color{popline}\MakeUppercase{Équipe pédagogique OnBuch+}\\\MakeUppercase{Conforme au programme officiel MINESEC}\par}
    \end{minipage}\hfill
    {\poplogo\fontsize{22}{22}\selectfont\color{popcream}OnBuch\color{poporange}+}
  \end{tcolorbox}%
}

% ============================================================
% Page de garde
% #1 titre · #2 matière · #3 niveau · #4 module · #5 n° leçon · #6 durée ·
% #7 description / objectifs (texte ou itemize)
% ============================================================
\newcommand{\pagegarde}[7]{%
  \thispagestyle{empty}
  \newgeometry{a4paper,margin=1.7cm,top=1.6cm,bottom=1.4cm}%
  \begin{tikzpicture}[remember picture,overlay]
    \fill[poporange!18] ([xshift=-3.2cm,yshift=-3.6cm]current page.north east) circle (4.6cm);
    \fill[poppurple!14] ([xshift=2.4cm,yshift=7.6cm]current page.south west) circle (3.8cm);
  \end{tikzpicture}%
  {\poplogo\fontsize{30}{30}\selectfont\color{popink}OnBuch\color{poporange}+}\hfill
  \raisebox{0.6ex}{\poppill{Cours signé \textbullet\ Premium}{popink}{popcream}}\par
  \vspace{1pt}\popsquiggle{poppurple}\par
  \vspace{8pt}
  \begin{tcolorbox}[enhanced,colback=poporange,colframe=popink,boxrule=2.8pt,arc=13pt,
      drop shadow=popink,left=18pt,right=18pt,top=12pt,bottom=13pt,after skip=10pt]
    \poppill{#2 \textbullet\ #3}{popink}{popcream}\hspace{5pt}\poppill{#4}{white}{popink}\par\vspace{8pt}
    {\popdisplay\fontsize{14}{14}\selectfont\color{popink}LEÇON #5}\par\vspace{4pt}
    {\raggedright\hyphenpenalty=10000\exhyphenpenalty=10000\popdisplay\fontsize{25}{27}\selectfont\color{popcream}\MakeUppercase{#1}\par}
    \vspace{8pt}
    {\popxbold\fontsize{9.5}{9.5}\selectfont\color{popink}\MakeUppercase{Durée conseillée : #6}}
  \end{tcolorbox}
  \begin{tcolorbox}[enhanced,colback=white,colframe=popink,boxrule=2.2pt,arc=10pt,
      drop shadow=popink,left=16pt,right=16pt,top=10pt,bottom=10pt,after skip=8pt]
    \poppill{Dans cette leçon}{poppurple}{white}\par\vspace{5pt}
    \popsemi\fontsize{9.3}{12.6}\selectfont\color{popink}#7
  \end{tcolorbox}
  \noindent\begin{minipage}[t]{0.487\linewidth}
    \begin{tcolorbox}[enhanced,colback=popgreen,colframe=popink,boxrule=2.2pt,arc=10pt,
        drop shadow=popink,left=11pt,right=11pt,top=10pt,bottom=10pt,height=3.1cm,valign=center]
      \tcbox[colback=white,colframe=popink,boxrule=1.3pt,arc=5pt,boxsep=0pt,nobeforeafter,
        left=3pt,right=3pt,top=3pt,bottom=3pt,valign=center]{\qrcode[height=1.9cm]{https://play.google.com/store/apps/details?id=com.luvvix.onbuch&hl=fr}}%
      \hspace{8pt}\begin{minipage}[c]{0.5\linewidth}
        {\popdisplay\fontsize{11}{12.5}\selectfont\color{white}\MakeUppercase{Télécharge OnBuch}}\par\vspace{4pt}
        {\raggedright\popsemi\fontsize{8.4}{11}\selectfont\color{white}Gratuit \textbullet\ Google Play\\Tuteur IA, annales, concours, résultats\par}
      \end{minipage}
    \end{tcolorbox}
  \end{minipage}\hfill
  \begin{minipage}[t]{0.487\linewidth}
    \begin{tcolorbox}[enhanced,colback=poppurple,colframe=popink,boxrule=2.2pt,arc=10pt,
        drop shadow=popink,left=11pt,right=11pt,top=10pt,bottom=10pt,height=3.1cm,valign=center]
      \tikz[baseline=-0.7ex]\node[circle,fill=white,draw=popink,line width=1.3pt,minimum size=1.6cm,inner sep=0]{\popdisplay\fontsize{24}{24}\selectfont\color{poppurple}+};%
      \hspace{8pt}\begin{minipage}[c]{0.6\linewidth}
        {\popdisplay\fontsize{11}{12.5}\selectfont\color{white}\MakeUppercase{Passe à OnBuch+}}\par\vspace{4pt}
        {\raggedright\popsemi\fontsize{8.4}{11}\selectfont\color{white}Tous les cours signés, Tuteur IA illimité, fiches PDF — onglet OnBuch+ de l'app\par}
      \end{minipage}
    \end{tcolorbox}
  \end{minipage}
  \vspace{8pt}
  \popcontenu
  \vfill
  \signatureonbuch
  \restoregeometry
  \newpage
}

% Rangée « Ce cours contient » (page de garde)
\newcommand{\poptile}[3]{% #1 couleur #2 gros texte #3 légende
  \begin{tcolorbox}[enhanced,colback=#1,colframe=popink,boxrule=2pt,arc=9pt,shadow={2.5pt}{-2.5pt}{0pt}{popink},
      left=6pt,right=6pt,top=9pt,bottom=9pt,halign=center,valign=center,height=1.75cm,width=\linewidth,nobeforeafter]
    {\popdisplay\fontsize{12.5}{13}\selectfont\color{white}\MakeUppercase{#2}}\par\vspace{3pt}
    {\centering\popsemi\fontsize{7.8}{9.5}\selectfont\color{white}#3\par}
  \end{tcolorbox}}
\newcommand{\popcontenu}{%
  \par\noindent{\popxboldls\fontsize{7.5}{7.5}\selectfont\color{popmuted}CE COURS SIGNÉ CONTIENT}\par\vspace{7pt}
  \noindent\begin{minipage}[t]{0.235\linewidth}\poptile{popblue}{Cours}{complet, illustré, pas à pas}\end{minipage}\hfill
  \begin{minipage}[t]{0.235\linewidth}\poptile{poporange}{Méthodes}{et exercices résolus}\end{minipage}\hfill
  \begin{minipage}[t]{0.235\linewidth}\poptile{poppink}{Exercices}{type examen + corrigés}\end{minipage}\hfill
  \begin{minipage}[t]{0.235\linewidth}\poptile{popgreen}{Fiche bilan}{l'essentiel à retenir}\end{minipage}\par}

% Sommaire
\newtcolorbox{sommaire}{enhanced,colback=white,colframe=popink,boxrule=2.2pt,arc=10pt,
  drop shadow=popink,left=18pt,right=18pt,top=13pt,bottom=13pt,before skip=6pt,after skip=20pt,
  before upper={\poppill{Sommaire}{poppurple}{white}\par\vspace{9pt}\popsemi\fontsize{10.6}{14.5}\selectfont\color{popink}}}

% Signature de fin de cours — poussée tout en bas de la dernière page
\newcommand{\findecours}{%
  \par\vfill\nopagebreak
  \begin{center}\popsquiggle{poporange}\end{center}\vspace{4pt}
  \signatureonbuch
}
\AtBeginDocument{\def\llbracket{\mathopen{[\mkern-3.5mu[}}\def\rrbracket{\mathclose{]\mkern-3.5mu]}}} % intervalles d'entiers (glyphe ⟦ absent de la police)
% Glyphes absents de Fira Math : redéfinis à partir de glyphes disponibles
\AtBeginDocument{%
  \def\setminus{\mathbin{\backslash}}\let\smallsetminus\setminus
  \def\vdots{\mathord{\vbox{\baselineskip3.2pt\lineskiplimit0pt\kern4pt\hbox{.}\hbox{.}\hbox{.}}}}%
  \def\ddots{\mathinner{\mkern1mu\raise6.5pt\hbox{.}\mkern2mu\raise3.6pt\hbox{.}\mkern2mu\raise0.7pt\hbox{.}\mkern1mu}}%
  \def\triangle{\mathord{\scalebox{0.9}{$\symup{\Delta}$}}}%
  \def\nabla{\mathord{\raisebox{\depth}{\scalebox{1}[-1]{$\symup{\Delta}$}}}}%
  \def\checkmark{\popcheck{popdark}}%
  \def\star{\mathbin{\ast}}\def\bigstar{\mathord{\ast}}%
  \def\diamond{\mathbin{\scalebox{0.6}{$\Diamond$}}}%
  \def\bigcup{\mathop{\mathchoice{\raisebox{-0.3em}{\scalebox{1.7}{$\cup$}}}{\scalebox{1.25}{$\cup$}}{\cup}{\cup}}\displaylimits}%
  \def\bigcap{\mathop{\mathchoice{\raisebox{-0.3em}{\scalebox{1.7}{$\cap$}}}{\scalebox{1.25}{$\cap$}}{\cap}{\cap}}\displaylimits}%
  \def\lesseqqgtr{\mathrel{\lessgtr}}\def\bigoplus{\mathop{\oplus}\displaylimits}%
}
\providecommand{\ssi}{si et seulement si} % abréviation fréquente
\AtBeginDocument{\providecommand{\wideparen}{\overparen}} % arc de cercle

%% FIGFIX-BEGIN v4 — correctifs globaux des figures (docs/onbuch-generation/tools/figfix)
\usepackage{adjustbox}\usepackage{etoolbox}
% toute figure TikZ/pgfplots plus large que la ligne est réduite à la largeur disponible
\BeforeBeginEnvironment{tikzpicture}{\begin{adjustbox}{max width=\linewidth}}
\AfterEndEnvironment{tikzpicture}{\end{adjustbox}}
% légendes pgfplots lisibles par-dessus les courbes
\pgfplotsset{every axis/.append style={legend style={fill=white,fill opacity=0.92,text opacity=1,draw=black!15,inner sep=3pt}}}
% étiquettes d'arêtes (syntaxe edge["..."]) avec fond blanc
\tikzset{every edge quotes/.append style={fill=white,inner sep=1.2pt}}
%% FIGFIX-END
\begin{document}

\pagegarde{Leçon 36 — Grammaire : antonyme 新 / 旧 ; verbe modal 可以}{Chinois}{1ère A}{Module 5 : Mass médias et télécommunication}{ZHA36}{2 h}{Cette leçon de grammaire étudie l'antonymie 新 (nouveau) / 旧 (vieux, ancien) et le verbe modal 可以 (pouvoir, permission/possibilité). L'élève apprend à construire, nier et interroger ces structures, à éviter les erreurs typiques des francophones et à les mobiliser dans des phrases et de courts textes.
\vspace{4pt}
\begin{multicols}{2}\small
\begin{itemize}
\item À la fin de cette leçon, je sais distinguer et employer correctement les antonymes 新 et 旧 devant un nom
\item je sais construire la structure 新/旧 + 的 + nom et expliquer le rôle de 的
\item je sais placer le verbe modal 可以 avant le verbe principal dans une phrase affirmative
\item je sais former la négation 不可以 / 不能 et distinguer leurs nuances
\item je sais poser une question de permission avec 可以...吗 et y répondre correctement
\item je sais comparer 可以 avec le français « pouvoir » et éviter les calques fautifs
\end{itemize}\end{multicols}}

\begin{sommaire}
\begin{multicols}{2}
\begin{enumerate}
\item Observation guidée des exemples
\item Les antonymes 新 / 旧
\item Le verbe modal 可以
\item Comparaison avec le français et pièges des francophones
\item Exercices d'application et de transformation
\item Activité d'intégration
\item Méthodes et exercices résolus
\item Exercices
\item Corrigés des exercices
\item Fiche bilan
\end{enumerate}
\end{multicols}
\end{sommaire}

\begin{objectifs}
\begin{itemize}
\item À la fin de cette leçon, je sais distinguer et employer correctement les antonymes 新 et 旧 devant un nom
\item je sais construire la structure 新/旧 + 的 + nom et expliquer le rôle de 的
\item je sais placer le verbe modal 可以 avant le verbe principal dans une phrase affirmative
\item je sais former la négation 不可以 / 不能 et distinguer leurs nuances
\item je sais poser une question de permission avec 可以...吗 et y répondre correctement
\item je sais comparer 可以 avec le français « pouvoir » et éviter les calques fautifs
\item je sais repérer et corriger les erreurs fréquentes (place de 可以, oubli de 的, double verbe)
\item je sais utiliser 新/旧 et 可以 dans un court texte sur les médias et les télécommunications
\end{itemize}
\end{objectifs}

\begin{prerequis}
\begin{itemize}
\item La structure 是 + nom et l'adjectif attributif avec 的 (mon nouveau téléphone 我的新手机)
\item Les pronoms personnels et possessifs 我/你/他 + 的
\item Le vocabulaire du Module 5 : 手机, 电脑, 电视, 报纸, 网络
\item La question avec 吗 et la négation avec 不
\item Les verbes d'action courants : 用, 看, 买, 去, 上网
\end{itemize}
\end{prerequis}

\newpage

\coursec{Observation guidée des exemples}

\courssub{Situation d'accroche : le nouveau téléphone d'Amina}

À Yaoundé, Amina vient d'acheter un nouveau téléphone pour remplacer son vieux Nokia. Elle montre fièrement ses deux appareils à ses camarades du lycée :

\zh{这是我的新手机，那是我的旧手机。} \emph{(Zhè shì wǒ de xīn shǒujī, nà shì wǒ de jiù shǒujī.)} « Voici mon nouveau téléphone, ça c'est mon vieux téléphone. »

Au marché Mokolo, son petit frère la regarde d'un air suppliant : « Est-ce que je peux utiliser ton nouveau téléphone ? »

\zh{我可以用你的新手机吗？} \emph{(Wǒ kěyǐ yòng nǐ de xīn shǒujī ma?)} « Je peux utiliser ton nouveau téléphone ? »

Amina répond avec un sourire : « Oui, mais tu ne peux pas aller sur Internet ! »

\zh{可以，但是你不可以上网。} \emph{(Kěyǐ, dànshì nǐ bù kěyǐ shàngwǎng.)} « Oui, mais tu ne peux pas aller sur Internet. »

Entre le \cle{neuf} et l'\cle{ancien}, entre le \cle{permis} et l'\cle{interdit}, le chinois exprime ces idées avec deux outils simples : l'antonymie \zh{新}/\zh{旧} et le verbe modal \zh{可以}. Comment les utiliser sans faute ? C'est toute la problématique de cette leçon : observer d'abord, deviner ensuite, vérifier enfin.

\courssub{Premier contact : le mini-texte d'Amina}

\begin{textebox}[Le nouveau téléphone d'Amina — 阿敏的新手机]
这是我的新手机，那是我的旧手机。\\
\emph{Zhè shì wǒ de xīn shǒujī, nà shì wǒ de jiù shǒujī.}\\
« Voici mon nouveau téléphone, ça c'est mon vieux téléphone. »\\[4pt]
我可以用新手机上网，不可以用旧手机上网。\\
\emph{Wǒ kěyǐ yòng xīn shǒujī shàngwǎng, bù kěyǐ yòng jiù shǒujī shàngwǎng.}\\
« Je peux aller sur Internet avec le nouveau téléphone, je ne peux pas aller sur Internet avec le vieux téléphone. »\\[4pt]
我弟弟问：« 我可以用你的新手机吗？»\\
\emph{Wǒ dìdi wèn : « Wǒ kěyǐ yòng nǐ de xīn shǒujī ma? »}\\
« Mon petit frère demande : "Je peux utiliser ton nouveau téléphone ?" »\\[4pt]
我说：« 可以，但是你不可以玩游戏！»\\
\emph{Wǒ shuō : « Kěyǐ, dànshì nǐ bù kěyǐ wán yóuxì! »}\\
« Je dis : "Oui, mais tu ne peux pas jouer aux jeux vidéo !" »
\end{textebox}

\textbf{Petit glossaire du texte}

\begin{center}
\begin{tabularx}{\linewidth}{|X|X|X|X|}
\hline
\rowcolor{popblueL} Caractères & Pinyin & Nature & Traduction \\
\hline
新 & xīn & adjectif & nouveau, neuf \\
\hline
旧 & jiù & adjectif & vieux, ancien (objet) \\
\hline
可以 & kěyǐ & verbe modal & pouvoir (permission) \\
\hline
手机 & shǒujī & nom & téléphone portable \\
\hline
上网 & shàngwǎng & verbe & aller sur Internet \\
\hline
用 & yòng & verbe & utiliser \\
\hline
弟弟 & dìdi & nom & petit frère \\
\hline
玩游戏 & wán yóuxì & locution verbale & jouer (à des jeux) \\
\hline
\end{tabularx}
\end{center}

\courssub{Les quatre phrases modèles}

Lisons maintenant les quatre phrases modèles de la leçon. Pour chacune, lis d'abord les caractères, puis le pinyin, puis la traduction, sans encore chercher la règle.

\begin{exemplebox}[Phrases modèles à observer]
\textbf{Phrase 1} : \zh{这是我的新手机。}\\
\emph{Zhè shì wǒ de xīn shǒujī.} — « Voici mon nouveau téléphone. »\\[4pt]
\textbf{Phrase 2} : \zh{那是他的旧电脑。}\\
\emph{Nà shì tā de jiù diànnǎo.} — « Ça, c'est son vieil ordinateur. »\\[4pt]
\textbf{Phrase 3} : \zh{我可以用你的手机吗？}\\
\emph{Wǒ kěyǐ yòng nǐ de shǒujī ma?} — « Je peux utiliser ton téléphone ? »\\[4pt]
\textbf{Phrase 4} : \zh{可以，但是你不可以上网。}\\
\emph{Kěyǐ, dànshì nǐ bù kěyǐ shàngwǎng.} — « Oui, mais tu ne peux pas aller sur Internet. »
\end{exemplebox}

\begin{popfigure}
\begin{center}
\begin{tabularx}{\linewidth}{|X|X|}
\hline
\rowcolor{popblueL} Phrase chinoise (caractères + pinyin) & Traduction française \\
\hline
这是我的\textbf{新}手机。\newline \emph{Zhè shì wǒ de \textbf{xīn} shǒujī.} & Voici mon \textbf{nouveau} téléphone. \\
\hline
那是他的\textbf{旧}电脑。\newline \emph{Nà shì tā de \textbf{jiù} diànnǎo.} & Ça, c'est son \textbf{vieil} ordinateur. \\
\hline
我\textbf{可以}用你的手机吗？\newline \emph{Wǒ \textbf{kěyǐ} yòng nǐ de shǒujī ma?} & Je \textbf{peux} utiliser ton téléphone ? \\
\hline
\textbf{可以}，但是你不\textbf{可以}上网。\newline \emph{\textbf{Kěyǐ}, dànshì nǐ bù \textbf{kěyǐ} shàngwǎng.} & \textbf{Oui}, mais tu ne \textbf{peux pas} aller sur Internet. \\
\hline
\end{tabularx}
\end{center}
\legende{Tableau d'observation : les mots cibles 新, 旧 et 可以 sont mis en gras.}
\end{popfigure}

\courssub{Repérage guidé : à toi de chercher}

Avant toute explication du professeur, observe les quatre phrases et réponds aux questions suivantes, par écrit, sur ton cahier d'exercices.

\begin{enumerate}
\item Souligne en rouge les mots \zh{新} et \zh{旧} dans les phrases 1 et 2. Quel mot se trouve juste après chacun d'eux ? S'agit-il d'un verbe ou d'un nom ?
\item Souligne en bleu le mot \zh{可以} dans les phrases 3 et 4. Où se place-t-il par rapport au verbe (\zh{用} « utiliser », \zh{上网} « aller sur Internet ») : avant ou après ?
\item Dans la phrase 4, quel petit mot transforme \zh{可以} en son contraire ? Que signifie alors la phrase ?
\item Compare les traductions françaises : à quels mots français correspondent \zh{新}, \zh{旧} et \zh{可以} ?
\end{enumerate}

\begin{methode}[Observer une structure grammaticale en quatre étapes]
\begin{enumerate}
\item \textbf{Lire} la phrase chinoise et sa traduction française, à voix haute, en suivant le pinyin.
\item \textbf{Repérer} le mot nouveau ou la structure nouvelle (ici \zh{新}, \zh{旧}, \zh{可以}) et le souligner.
\item \textbf{Noter sa position} dans la phrase : avant ou après le nom ? avant ou après le verbe ? accompagné de quel mot-outil (\zh{的}, \zh{不}, \zh{吗}) ?
\item \textbf{Formuler une hypothèse de règle} en une phrase simple, puis la vérifier sur un deuxième exemple avant de l'accepter.
\end{enumerate}
\end{methode}

\courssub{Tes hypothèses : confrontons-les aux faits}

\textbf{Hypothèse 1 : \zh{新} et \zh{旧} sont des adjectifs contraires.}

C'est exact. \zh{新} \emph{(xīn)} signifie « nouveau, neuf » et \zh{旧} \emph{(jiù)} signifie « vieux, ancien » — mais uniquement pour les \emph{objets} et les choses, jamais pour l'âge des personnes. Dans les deux phrases modèles, l'adjectif se place \emph{avant le nom}, exactement comme en français dans « mon \emph{nouveau} téléphone » :

\begin{center}
\begin{tabularx}{\linewidth}{|X|X|X|}
\hline
\rowcolor{popblueL} Structure & Exemple (caractères + pinyin) & Traduction \\
\hline
Possesseur + 的 + 新/旧 + nom & 我的新手机 \emph{wǒ de xīn shǒujī} & mon nouveau téléphone \\
\hline
Possesseur + 的 + 新/旧 + nom & 他的旧电脑 \emph{tā de jiù diànnǎo} & son vieil ordinateur \\
\hline
\end{tabularx}
\end{center}

\textbf{Hypothèse 2 : \zh{可以} se place avant le verbe.}

C'est exact aussi. \zh{可以} \emph{(kěyǐ)} est un \cle{verbe modal} : il exprime la \emph{permission} (« pouvoir, avoir le droit de ») et il se place toujours \emph{entre le sujet et le verbe principal}, qu'il accompagne obligatoirement :

\begin{center}
Sujet + \zh{可以} + Verbe (+ complément) + \zh{吗} ?
\end{center}

\zh{我可以用你的手机吗？} \emph{(Wǒ kěyǐ yòng nǐ de shǒujī ma?)} « Je peux utiliser ton téléphone ? » — Ici \zh{可以} précède le verbe \zh{用} « utiliser ».

\textbf{Hypothèse 3 : la négation se fait avec \zh{不}.}

Exact : \zh{不可以} \emph{(bù kěyǐ)} signifie « ne pas pouvoir, ne pas avoir le droit ». Le mot-outil \zh{不} se place \emph{avant} \zh{可以}, jamais entre \zh{可以} et le verbe : on dit \zh{你不可以上网} « tu ne peux pas aller sur Internet », jamais \zh{你可以不上网} dans ce sens.

\begin{attention}[Ne traduis pas mot à mot !]
\zh{可以} n'est \textbf{pas un verbe complet} : il ne peut jamais terminer une phrase à lui tout seul (sauf comme réponse courte : « — 我可以去吗？— 可以。 » « — Je peux y aller ? — Oui. »). Il \emph{accompagne toujours un autre verbe}. Le francophone qui traduit « Je peux utiliser » par \zh{我可以} oublie le verbe \zh{用} : la phrase correcte est \zh{我可以用}. De même, ne confonds pas \zh{可以} (permission : « avoir le droit ») avec \zh{会} (capacité apprise : « savoir ») : « Je sais nager » se dit \zh{我会游泳}, pas \zh{我可以游泳}.
\end{attention}

\courssub{Carte mentale de première observation}

\begin{popfigure}
\begin{center}
\begin{tikzpicture}[mindmap, grow cyclic,
  every node/.style={concept, align=center, font=\small},
  concept color=popblue,
  level 1/.append style={level distance=3.2cm, sibling angle=120},
  level 2/.append style={level distance=2.4cm, sibling angle=60}]
\node[concept color=popblue, text=white, align=center]{Leçon 36\\ deux outils}
  child[concept color=popgreen] { node[align=center]{新 xīn\\ « nouveau »}
    child { node[concept color=popgreenL, text=popdark, align=center]{avant le nom\\ 新手机}
    }
    child { node[concept color=popgreenL, text=popdark, align=center]{objets\\ seulement}
    }
  }
  child[concept color=poporange] { node[align=center]{旧 jiù\\ « vieux »}
    child { node[concept color=poporangeL, text=popdark, align=center]{avant le nom\\ 旧电脑}
    }
    child { node[concept color=poporangeL, text=popdark, align=center]{antonyme\\ de 新}
    }
  }
  child[concept color=poppurple] { node[align=center]{可以 kěyǐ\\ « pouvoir »}
    child { node[concept color=poppurpleL, text=popdark, align=center]{avant le verbe\\ 可以用}
    }
    child { node[concept color=poppurpleL, text=popdark, align=center]{négation\\ 不可以}
    }
  };
\end{tikzpicture}
\end{center}
\legende{Carte mentale : ce que l'observation des quatre phrases nous a déjà appris.}
\end{popfigure}

\courssub{Vérification : un exercice résolu}

\begin{exoresolu}[Repérage dans un mini-dialogue]
Énoncé — Dans le dialogue suivant, souligne tous les emplois de \zh{新}, \zh{旧} et \zh{可以}, puis indique pour chacun sa fonction (adjectif devant nom / modal devant verbe / réponse courte).

\zh{A：这是我的新自行车。}\\
\zh{B：真漂亮！我可以用一下吗？}\\
\zh{A：可以，但是你不可以骑我的旧自行车，它坏了。}

\emph{A : Zhè shì wǒ de xīn zìxíngchē. B : Zhēn piàoliang! Wǒ kěyǐ yòng yíxià ma? A : Kěyǐ, dànshì nǐ bù kěyǐ qí wǒ de jiù zìxíngchē, tā huài le.}

« A : Voici mon nouveau vélo. B : Comme il est beau ! Je peux l'essayer un peu ? A : Oui, mais tu ne peux pas utiliser mon vieux vélo, il est en panne. »

\tcblower

Solution détaillée —
\begin{enumerate}
\item \zh{新} dans \zh{新自行车} : adjectif « nouveau » placé \emph{avant le nom} \zh{自行车} « vélo ».
\item \zh{可以} dans \zh{我可以用一下吗} : verbe modal de permission placé \emph{avant le verbe} \zh{用} « utiliser » ; \zh{一下} adoucit la demande (« un peu ») et la particule \zh{吗} marque la question.
\item \zh{可以} seul au début de la réponse d'A : \emph{réponse courte} signifiant « oui, tu peux » — c'est le seul cas où \zh{可以} s'emploie sans verbe derrière.
\item \zh{不可以} dans \zh{你不可以骑} : négation du modal (\zh{不} + \zh{可以}) placée \emph{avant le verbe} \zh{骑} « monter (à vélo, à cheval) ».
\item \zh{旧} dans \zh{旧自行车} : adjectif « vieux » placé \emph{avant le nom}, antonyme de \zh{新}.
\end{enumerate}
\end{exoresolu}

\begin{exercice}[À ton tour : premières hypothèses]
Observe les trois phrases suivantes et, sans consulter la suite de la leçon, formule par écrit une hypothèse de règle pour chaque mot en gras : \zh{新}, \zh{旧}, \zh{可以}.
\begin{enumerate}
\item \zh{我们学校的新图书馆很大。} \emph{(Wǒmen xuéxiào de xīn túshūguǎn hěn dà.)} « La nouvelle bibliothèque de notre école est très grande. »
\item \zh{他的旧手机不能上网。} \emph{(Tā de jiù shǒujī bù néng shàngwǎng.)} « Son vieux téléphone ne peut pas aller sur Internet. »
\item \zh{老师，我可以问一个问题吗？} \emph{(Lǎoshī, wǒ kěyǐ wèn yí ge wèntí ma?)} « Monsieur le professeur, je peux poser une question ? »
\end{enumerate}
\end{exercice}

\begin{corrige}[Exercice « À ton tour »]
\begin{enumerate}
\item \zh{新} est un adjectif placé entre le mot-outil \zh{的} et le nom \zh{图书馆} « bibliothèque » : structure « possesseur + 的 + adjectif + nom ».
\item \zh{旧} occupe exactement la même position : adjectif avant le nom \zh{手机}. On confirme que \zh{新} et \zh{旧} fonctionnent comme une paire d'antonymes devant un nom d'objet.
\item \zh{可以} est placé entre le sujet \zh{我} et le verbe \zh{问} « demander, poser » : c'est un verbe modal de permission. La question se termine par \zh{吗}.
\end{enumerate}
Hypothèse générale validée : \emph{adjectif avant le nom, modal avant le verbe}. Nous allons maintenant étudier chaque outil en détail.
\end{corrige}

\begin{aretenir}[Ce que l'observation nous a déjà appris]
\begin{itemize}
\item \zh{新} \emph{(xīn)} « nouveau » et \zh{旧} \emph{(jiù)} « vieux » sont des \textbf{adjectifs antonymes} qui se placent \emph{avant le nom} qu'ils qualifient : \zh{新手机} « nouveau téléphone », \zh{旧电脑} « vieil ordinateur ».
\item \zh{可以} \emph{(kěyǐ)} est un \textbf{verbe modal de permission} qui se place \emph{avant le verbe principal} : \zh{我可以用} « je peux utiliser ».
\item La négation est \zh{不可以} \emph{(bù kěyǐ)} « ne pas avoir le droit » : \zh{你不可以上网} « tu ne peux pas aller sur Internet ».
\item En réponse courte, \zh{可以。} seul signifie « oui, tu peux / c'est d'accord ».
\end{itemize}
\end{aretenir}

\coursec{Les antonymes 新 et 旧}

\courssub{Deux adjectifs, une seule paire}

\begin{definition}[新 et 旧 : adjectifs antonymes]
\zh{新} \emph{(xīn)} signifie « nouveau, neuf, récent » ; \zh{旧} \emph{(jiù)} signifie « vieux, ancien, usagé ». Ce sont des \textbf{antonymes} : ils désignent les deux extrémités d'un même axe, celui de l'usure des \emph{choses}. Comme tous les adjectifs chinois, ils sont \emph{invariables} : pas de féminin, pas de pluriel.
\end{definition}

\begin{propriete}[Place de 新 / 旧 devant le nom]
Structure : \textbf{(Possesseur + 的 +) 新/旧 + Nom}. L'adjectif se colle directement devant le nom, sans mot-outil entre les deux. S'il y a un possesseur, le mot-outil \zh{的} se place entre le possesseur et l'adjectif.
\end{propriete}

\begin{exemplebox}[新 et 旧 devant le nom]
\zh{新手机} \emph{xīn shǒujī} — « nouveau téléphone »\\
\zh{旧衣服} \emph{jiù yīfu} — « vieux vêtements »\\
\zh{我的新书} \emph{wǒ de xīn shū} — « mon nouveau livre »\\
\zh{老师的旧电脑} \emph{lǎoshī de jiù diànnǎo} — « le vieil ordinateur du professeur »\\
\zh{我们学校的新图书馆} \emph{wǒmen xuéxiào de xīn túshūguǎn} — « la nouvelle bibliothèque de notre école »
\end{exemplebox}

\courssub{新 et 旧 en position attributive et en position épithète de phrase}

Comme les autres adjectifs chinois, \zh{新} et \zh{旧} peuvent aussi fonctionner \emph{après le sujet}, avec \zh{很} (ou un autre adverbe) et \textbf{sans} \zh{是} :

\begin{center}
\begin{tabularx}{\linewidth}{|X|X|X|}
\hline
\rowcolor{popblueL} Structure & Exemple (caractères + pinyin) & Traduction \\
\hline
Nom + 很 + 新/旧 & 这部手机很新。\emph{Zhè bù shǒujī hěn xīn.} & Ce téléphone est très neuf. \\
\hline
Nom + 很 + 新/旧 & 他的电脑很旧。\emph{Tā de diànnǎo hěn jiù.} & Son ordinateur est très vieux. \\
\hline
Nom + 不 + 新/旧 & 这本书不旧。\emph{Zhè běn shū bù jiù.} & Ce livre n'est pas vieux. \\
\hline
\end{tabularx}
\end{center}

\begin{attention}[Jamais 是 devant un adjectif !]
Le francophone traduit « le téléphone est neuf » par \zh{手机是新} : c'est une \textbf{faute}. En chinois, l'adjectif fait office de verbe : on dit \zh{手机很新} \emph{(shǒujī hěn xīn)}. Le verbe \zh{是} ne s'emploie qu'entre deux noms : \zh{这是我的新手机} « ceci est mon nouveau téléphone » (ici \zh{是} relie \zh{这} au groupe nominal \zh{我的新手机}).
\end{attention}

\courssub{Exception capitale : 旧 ne s'applique pas aux personnes}

\begin{propriete}[旧 = vieux pour les choses ; 老 = vieux pour les personnes]
\zh{旧} qualifie uniquement les \textbf{objets, bâtiments, idées} (téléphone, voiture, école, méthode). Pour dire qu'une \textbf{personne} est âgée, on utilise \zh{老} \emph{(lǎo)} : \zh{我的爷爷很老} « mon grand-père est très vieux ». Dire \zh{旧老师} pour « le vieux professeur » est une faute ; on dit \zh{老老师} n'existe pas non plus : on dit simplement \zh{年纪大的老师} ou, avec respect, on garde \zh{老} dans les appellations comme \zh{老王} « le vieux Wang ».
\end{propriete}

\begin{center}
\begin{tabularx}{\linewidth}{|X|X|X|}
\hline
\rowcolor{popblueL} Mot & S'applique à & Exemple \\
\hline
新 xīn & choses, personnes (nouveau venu) & 新同学 \emph{xīn tóngxué} « nouveau camarade » \\
\hline
旧 jiù & choses uniquement & 旧汽车 \emph{jiù qìchē} « vieille voiture » \\
\hline
老 lǎo & personnes (âge), ancienneté & 老人 \emph{lǎorén} « personne âgée » \\
\hline
\end{tabularx}
\end{center}

Note l'asymétrie : \zh{新} peut qualifier une personne au sens de « nouvellement arrivé » (\zh{新同学} « le nouveau camarade de classe », \zh{新老师} « le nouveau professeur »), mais son antonyme pour les personnes n'est pas \zh{旧}.

\courssub{Écriture des caractères 新 et 旧}

\textbf{新} \emph{(xīn)}, 13 traits. Caractère composé : à gauche \zh{亲} (qui donne le son), à droite la clé \zh{斤} « hache » (qui donne le sens : couper du bois frais, donc « nouveau »). Ordre d'écriture : on écrit d'abord la partie gauche de haut en bas (le \zh{立} puis le \zh{木}), puis la partie droite \zh{斤} : trait horizontal tombant, puis le grand trait vertical-courbe, puis l'horizontal, puis le vertical final. Règle générale respectée : \emph{gauche avant droite, haut avant bas}.

\textbf{旧} \emph{(jiù)}, 5 traits. À gauche un trait vertical (ressemblant à \zh{丨}), à droite la clé \zh{日} « soleil, jour » : les « jours passés » donnent l'idée d'ancien. Ordre : le vertical gauche d'abord, puis \zh{日} (vertical gauche, horizontal-vertical fermant, horizontal du milieu, horizontal de fermeture en bas).

\begin{savaistu}[D'où viennent ces caractères ?]
La forme ancienne de \zh{旧} contenait un oiseau (\zh{臼} évoquant un nid) : un « vieil » oiseau retourne à son nid. La simplification moderne n'a gardé que \zh{日}. Quant à \zh{新}, la « hache » \zh{斤} rappelle que le bois fraîchement coupé est « neuf ». Retenir l'histoire d'un caractère aide à ne plus l'oublier !
\end{savaistu}

\coursec{Le verbe modal 可以}

\courssub{Définition et structure}

\begin{definition}[可以, verbe modal de permission]
\zh{可以} \emph{(kěyǐ)} est un \textbf{verbe modal} (ou auxiliaire) qui exprime la \textbf{permission} : « pouvoir, avoir le droit de, il est permis de ». Comme tous les modaux chinois (\zh{会}, \zh{能}, \zh{想}, \zh{要}), il se place \emph{avant le verbe principal} et ne se conjugue jamais.
\end{definition}

\begin{propriete}[Les quatre schémas de 可以]
\begin{itemize}
\item \textbf{Affirmation} : Sujet + 可以 + Verbe (+ complément).\\ \zh{我可以用你的手机。} \emph{Wǒ kěyǐ yòng nǐ de shǒujī.} « Je peux utiliser ton téléphone. »
\item \textbf{Négation} : Sujet + 不 + 可以 + Verbe.\\ \zh{你不可以玩游戏。} \emph{Nǐ bù kěyǐ wán yóuxì.} « Tu ne peux pas jouer aux jeux. »
\item \textbf{Question} : Sujet + 可以 + Verbe + 吗 ?\\ \zh{我可以用一下吗？} \emph{Wǒ kěyǐ yòng yíxià ma?} « Je peux l'utiliser un peu ? »
\item \textbf{Réponse courte} : 可以。/ 不可以。\\ « Oui (tu peux). / Non (tu ne peux pas). »
\end{itemize}
\end{propriete}

\begin{exemplebox}[可以 dans la vie du lycée]
\zh{老师，我可以问一个问题吗？} \emph{Lǎoshī, wǒ kěyǐ wèn yí ge wèntí ma?} — « Professeur, je peux poser une question ? »\\
\zh{下课以后，我们可以去图书馆。} \emph{Xiàkè yǐhòu, wǒmen kěyǐ qù túshūguǎn.} — « Après les cours, nous pouvons aller à la bibliothèque. »\\
\zh{上课的时候不可以打电话。} \emph{Shàngkè de shíhou bù kěyǐ dǎ diànhuà.} — « Pendant les cours, on ne peut pas téléphoner. »\\
\zh{在杜阿拉，你可以坐船。} \emph{Zài Dù'ālā, nǐ kěyǐ zuò chuán.} — « À Douala, tu peux prendre le bateau. »
\end{exemplebox}

\courssub{可以, 能, 会 : trois « pouvoir » à ne pas confondre}

Le français n'a qu'un mot « pouvoir » ; le chinois en a trois, et l'examen de Première les oppose souvent :

\begin{center}
\begin{tabularx}{\linewidth}{|X|X|X|}
\hline
\rowcolor{popblueL} Modal & Sens & Exemple \\
\hline
可以 kěyǐ & permission (on m'autorise) & 我可以出去吗？\emph{Wǒ kěyǐ chūqù ma?} « Je peux sortir ? » \\
\hline
能 néng & possibilité matérielle, capacité physique & 他的旧手机不能上网。\emph{Tā de jiù shǒujī bù néng shàngwǎng.} « Son vieux téléphone ne peut pas (techniquement) aller sur Internet. » \\
\hline
会 huì & capacité apprise (savoir-faire) & 我会说汉语。\emph{Wǒ huì shuō Hànyǔ.} « Je sais parler chinois. » \\
\hline
\end{tabularx}
\end{center}

\begin{attention}[Permission ou possibilité ?]
Compare : \zh{你不可以上网} « tu n'as \emph{pas le droit} d'aller sur Internet » (interdiction d'Amina) et \zh{旧手机不能上网} « le vieux téléphone \emph{ne peut pas} aller sur Internet » (impossibilité technique). Traduire l'interdiction par \zh{不能} ou l'impossibilité technique par \zh{不可以} est l'erreur classique du francophone. Astuce : si le sujet est une \emph{chose}, c'est presque toujours \zh{能} ; si c'est une \emph{personne} qu'on autorise ou qu'on interdit, c'est \zh{可以}.
\end{attention}

\courssub{La réponse courte : un usage à part}

À la question \zh{我可以……吗？}, on répond :

\begin{itemize}
\item \zh{可以。} \emph{Kěyǐ.} « Oui, tu peux. » (ou \zh{可以，可以！} pour insister avec chaleur)
\item \zh{不可以。} \emph{Bù kěyǐ.} « Non, tu ne peux pas. »
\item Plus poli pour refuser : \zh{对不起，不可以。} \emph{Duìbuqǐ, bù kěyǐ.} « Désolé, ce n'est pas possible. »
\end{itemize}

C'est le \textbf{seul} contexte où \zh{可以} apparaît sans verbe derrière lui : le verbe est sous-entendu car il vient d'être prononcé dans la question.

\coursec{Comparaison avec le français et pièges des francophones}

\courssub{Ce qui se ressemble, ce qui diffère}

\begin{center}
\begin{tabularx}{\linewidth}{|X|X|}
\hline
\rowcolor{popblueL} Français & Chinois \\
\hline
L'adjectif s'accorde (nouveau, nouvelle, nouveaux) & 新 et 旧 sont invariables \\
\hline
L'adjectif peut se placer après le nom (un téléphone neuf) & 新/旧 se placent toujours avant le nom \\
\hline
« est neuf » avec le verbe être & 很新 sans 是 : 手机很新 \\
\hline
« pouvoir » se conjugue (je peux, tu peux) & 可以 est invariable, suivi directement du verbe \\
\hline
« ne pas pouvoir » : la négation encadre le verbe & 不可以 : 不 se place avant 可以, un seul bloc \\
\hline
\end{tabularx}
\end{center}

\courssub{Les cinq erreurs à éliminer}

\begin{attention}[Erreurs fréquentes des francophones]
\begin{enumerate}
\item \textbf{Oublier le verbe après 可以} : « Je peux utiliser » n'est pas \zh{我可以} mais \zh{我可以用}. Le modal ne vaut rien sans son verbe.
\item \textbf{Mettre 是 devant l'adjectif} : « le téléphone est neuf » n'est pas \zh{手机是新} mais \zh{手机很新}.
\item \textbf{Placer 不 au mauvais endroit} : « tu ne peux pas jouer » est \zh{你不可以玩}, jamais \zh{你可以不玩} (qui signifierait « tu peux ne pas jouer », c'est-à-dire « ce n'est pas obligatoire » — tout autre sens !).
\item \textbf{Utiliser 旧 pour une personne} : « mon vieux grand-père » n'est pas \zh{旧爷爷} ; pour les personnes, on utilise \zh{老}.
\item \textbf{Confondre 可以 et 会} : « je sais nager » est \zh{我会游泳} (capacité apprise), pas \zh{我可以游泳} (qui voudrait dire « j'ai le droit de nager »).
\end{enumerate}
\end{attention}

\courssub{Schéma récapitulatif des transformations}

\begin{popfigure}
\begin{center}
\begin{tikzpicture}[node distance=0.6cm,
  box/.style={draw=popblue, fill=popblueL, rounded corners, align=center, font=\small, inner sep=5pt},
  arr/.style={-{Stealth[length=3mm]}, thick, popdark}]
\node[box, align=center] (aff){\textbf{Affirmation}\\ 我可以用。\\\emph{Wǒ kěyǐ yòng.}};
\node[box, right=2.2cm of aff, align=center] (neg){\textbf{Négation}\\ 我不可以用。\\\emph{Wǒ bù kěyǐ yòng.}};
\node[box, right=2.2cm of neg, align=center] (que){\textbf{Question}\\ 我可以用吗？\\\emph{Wǒ kěyǐ yòng ma?}};
\draw[arr] (aff) -- node[above, font=\footnotesize]{+ 不} (neg);
\draw[arr] (neg) -- node[above, font=\footnotesize]{+ 吗} (que);
\end{tikzpicture}
\end{center}
\legende{Transformations autour de 可以 : la négation ajoute 不 devant le modal ; la question ajoute 吗 en fin de phrase.}
\end{popfigure}

\coursec{Exercices d'application et de transformation}

\courssub{Exercice résolu}

\begin{exoresolu}[Transformation : affirmation, négation, question]
Énoncé — À partir de la phrase \zh{他可以用我的新电脑。} \emph{(Tā kěyǐ yòng wǒ de xīn diànnǎo.)} « Il peut utiliser mon nouvel ordinateur », forme : a) la négation ; b) la question avec \zh{吗} ; c) la réponse courte négative.

\tcblower

Solution détaillée —
\begin{enumerate}
\item[a)] Négation : on insère \zh{不} \emph{devant} \zh{可以} : \zh{他不可以用我的新电脑。} \emph{(Tā bù kěyǐ yòng wǒ de xīn diànnǎo.)} « Il ne peut pas utiliser mon nouvel ordinateur. »
\item[b)] Question : on ajoute \zh{吗} en fin de phrase et on remplace le point par le point d'interrogation : \zh{他可以用我的新电脑吗？} \emph{(Tā kěyǐ yòng wǒ de xīn diànnǎo ma?)} « Il peut utiliser mon nouvel ordinateur ? »
\item[c)] Réponse courte négative : \zh{不可以。} \emph{(Bù kěyǐ.)} « Non, il ne peut pas. »
\end{enumerate}
\end{exoresolu}

\courssub{Exercices d'entraînement}

\begin{exercice}[Exercice 1 : complète avec 新 ou 旧]
\begin{enumerate}
\item \zh{这是我的\_\_\_手机，那是我的\_\_\_手机。} (Voici mon nouveau téléphone, ça c'est mon vieux téléphone.)
\item \zh{我们学校有一个\_\_\_图书馆。} (Notre école a une nouvelle bibliothèque.)
\item \zh{他的\_\_\_自行车坏了。} (Son vieux vélo est en panne.)
\item \zh{我们班有两个\_\_\_同学。} (Notre classe a deux nouveaux camarades.)
\end{enumerate}
\end{exercice}

\begin{exercice}[Exercice 2 : complète avec 可以, 能 ou 会]
\begin{enumerate}
\item \zh{老师，我\_\_\_问一个问题吗？} (Professeur, je peux poser une question ?)
\item \zh{这个旧手机不\_\_\_上网。} (Ce vieux téléphone ne peut pas aller sur Internet.)
\item \zh{我\_\_\_说一点儿汉语。} (Je sais parler un peu chinois.)
\item \zh{上课的时候不\_\_\_打电话。} (Pendant les cours, on ne peut pas téléphoner.)
\end{enumerate}
\end{exercice}

\begin{exercice}[Exercice 3 : traduis en chinois (thème)]
\begin{enumerate}
\item Voici mon nouveau téléphone, ça c'est mon vieil ordinateur.
\item Je peux utiliser ton téléphone ?
\item Oui, mais tu ne peux pas jouer aux jeux vidéo.
\item La nouvelle bibliothèque de notre école est très grande.
\end{enumerate}
\end{exercice}

\begin{exercice}[Exercice 4 : traduis en français (version)]
\begin{enumerate}
\item \zh{我可以用你的新自行车吗？}
\item \zh{不可以，我的新自行车坏了。}
\item \zh{下课以后，我们可以去图书馆看书。}
\end{enumerate}
\end{exercice}

\begin{corrige}[Exercice 1]
\begin{enumerate}
\item \zh{新} ; \zh{旧} — \zh{这是我的新手机，那是我的旧手机。}
\item \zh{新} — \zh{我们学校有一个新图书馆。} (classificateur \zh{个} devant \zh{新图书馆})
\item \zh{旧} — \zh{他的旧自行车坏了。}
\item \zh{新} — \zh{我们班有两个新同学。} (\zh{新} peut qualifier une personne au sens de « nouvellement arrivé »)
\end{enumerate}
\end{corrige}

\begin{corrige}[Exercice 2]
\begin{enumerate}
\item \zh{可以} — permission demandée au professeur.
\item \zh{能} — impossibilité technique (le sujet est une chose).
\item \zh{会} — capacité apprise (savoir parler).
\item \zh{可以} — interdiction (règle de vie scolaire).
\end{enumerate}
\end{corrige}

\begin{corrige}[Exercice 3]
\begin{enumerate}
\item \zh{这是我的新手机，那是我的旧电脑。} \emph{(Zhè shì wǒ de xīn shǒujī, nà shì wǒ de jiù diànnǎo.)}
\item \zh{我可以用你的手机吗？} \emph{(Wǒ kěyǐ yòng nǐ de shǒujī ma?)}
\item \zh{可以，但是你不可以玩游戏。} \emph{(Kěyǐ, dànshì nǐ bù kěyǐ wán yóuxì.)}
\item \zh{我们学校的新图书馆很大。} \emph{(Wǒmen xuéxiào de xīn túshūguǎn hěn dà.)}
\end{enumerate}
\end{corrige}

\begin{corrige}[Exercice 4]
\begin{enumerate}
\item \zh{我可以用你的新自行车吗？} \emph{(Wǒ kěyǐ yòng nǐ de xīn zìxíngchē ma?)} — « Je peux utiliser ton nouveau vélo ? »
\item \zh{不可以，我的新自行车坏了。} \emph{(Bù kěyǐ, wǒ de xīn zìxíngchē huài le.)} — « Non, mon nouveau vélo est en panne. »
\item \zh{下课以后，我们可以去图书馆看书。} \emph{(Xiàkè yǐhòu, wǒmen kěyǐ qù túshūguǎn kànshū.)} — « Après les cours, nous pouvons aller lire à la bibliothèque. »
\end{enumerate}
\end{corrige}

\courssub{Activité orale : à toi de jouer}

\begin{experience}[Jeu de rôle : le prêt du téléphone]
Par groupes de deux, jouez la scène d'Amina et de son petit frère, en la transposant : un élève possède un objet neuf (téléphone, vélo, livre, ballon de football) et un objet vieux ; l'autre demande la permission d'utiliser l'objet neuf.
\begin{enumerate}
\item L'élève A présente ses deux objets avec \zh{新} et \zh{旧} : \zh{这是我的新……，那是我的旧……。}
\item L'élève B demande la permission : \zh{我可以用你的新……吗？}
\item L'élève A accorde ou refuse, avec une condition : \zh{可以，但是你不可以……。} ou \zh{对不起，不可以。}
\item Inversez les rôles, puis présentez votre dialogue devant la classe.
\end{enumerate}
Critères de réussite : modal placé avant le verbe, négation \zh{不可以} correcte, tons de \zh{xīn} (1\ier{} ton) et \zh{jiù} (4\ieme{} ton) bien distingués.
\end{experience}

\begin{savaistu}[Le téléphone en Chine et au Cameroun]
En Chine comme au Cameroun, le téléphone portable (\zh{手机}, littéralement « machine à main ») est l'objet indispensable du quotidien. En Chine, on paie presque tout avec le téléphone, même au marché ; au Cameroun, le paiement mobile se développe aussi rapidement à Douala et à Yaoundé. Demander \zh{我可以用一下你的手机吗？} « je peux utiliser ton téléphone un moment ? » est une phrase que tu pourras vraiment utiliser !
\end{savaistu}

\begin{aretenir}[Bilan de la leçon 36]
\begin{itemize}
\item \textbf{新 / 旧} : adjectifs antonymes, invariables, placés \emph{avant le nom} (\zh{新手机}, \zh{旧电脑}) ; en position après le sujet, avec \zh{很} et sans \zh{是} (\zh{手机很新}). \zh{旧} ne qualifie que les choses ; pour les personnes âgées, on utilise \zh{老}.
\item \textbf{可以} : verbe modal de \emph{permission}, placé entre le sujet et le verbe : Sujet + 可以 + Verbe. Négation : \zh{不可以}. Question : \zh{……吗？}. Réponse courte : \zh{可以。/ 不可以。}
\item \textbf{Distinction des modaux} : \zh{可以} = permission, \zh{能} = possibilité matérielle, \zh{会} = capacité apprise.
\item \textbf{Pièges} : ne jamais oublier le verbe après \zh{可以} ; ne jamais mettre \zh{是} devant un adjectif ; \zh{不} se place devant \zh{可以}, pas entre \zh{可以} et le verbe.
\end{itemize}
\end{aretenir}

\coursec{Leçon 36 — Grammaire : antonymes 新 / 旧 ; verbe modal 可以}

\begin{textebox}[Dialogue d'observation : Au magasin de téléphonie à Yaoundé]
\zh{— 老板，这个手机是新的吗？} \\
\emph{Lǎobǎn, zhège shǒujī shì xīn de ma?} \\
« Patron, ce téléphone est-il neuf ? » \\[4pt]
\zh{— 当然是新的。旧的我都不卖。您可以试一试。} \\
\emph{Dāngrán shì xīn de. Jiù de wǒ dōu bù mài. Nín kěyǐ shì yi shì.} \\
« Bien sûr qu'il est neuf. Les vieux, je ne les vends pas. Vous pouvez l'essayer. » \\[4pt]
\zh{— 好的。那旧手机怎么处理？} \\
\emph{Hǎo de. Nà jiù shǒujī zěnme chǔlǐ?} \\
« D'accord. Et les vieux téléphones, comment les gérez-vous ? » \\[4pt]
\zh{— 可以回收，也可以换新的。} \\
\emph{Kěyǐ huíshōu, yě kěyǐ huàn xīn de.} \\
« On peut les recycler, ou les échanger contre du neuf. »
\end{textebox}

\courssub{Observation guidée}
Dans ce dialogue authentique, repérez deux notions grammaticales essentielles :
\begin{enumerate}
\item Les \textbf{antonymes} \zh{新} (xīn, « neuf ») et \zh{旧} (jiù, « vieux ») qui qualifient le nom \zh{手机} (téléphone). Observez leur place : devant le nom (\zh{新手机}, \zh{旧手机}) ou après le verbe \zh{是} (\zh{是新的}).
\item Le \textbf{verbe modal} \zh{可以} (kěyǐ, « pouvoir, avoir la permission/possibilité ») qui exprime l'autorisation ou la possibilité objective : \zh{您可以试一试} (« Vous pouvez essayer »), \zh{可以回收} (« On peut recycler »).
\end{enumerate}
Nous allons les étudier l'une après l'autre.

\coursec{Les antonymes 新 (xīn) / 旧 (jiù)}

\courssub{Sens, nature et emploi de base}

\begin{definition}[Deux adjectifs antonymes invariables]
\zh{新} (xīn) signifie « nouveau, neuf, récent » ; \zh{旧} (jiù) signifie « vieux, ancien, usagé, d'occasion ». Ce sont des \textbf{adjectifs qualificatifs} qui s'appliquent aux \textbf{choses} (objets, lieux, vêtements, appareils, numéros, livres) et \textbf{jamais aux personnes} pour exprimer l'âge (voir encadré « Attention » p.~4). En chinois, les adjectifs sont \textbf{invariables} : pas de genre, pas de nombre.
\end{definition}

\begin{exemplebox}[Exemples contextuels (Module : Mass médias et télécommunication)]
\zh{这是一部新手机。} \emph{Zhè shì yī bù xīn shǒujī.} « C'est un nouveau téléphone portable. »\\
\zh{那是一台旧电脑。} \emph{Nà shì yī tái jiù diànnǎo.} « C'est un vieil ordinateur. »\\
\zh{我的新号码很好记。} \emph{Wǒ de xīn hàomǎ hěn hǎo jì.} « Mon nouveau numéro est facile à retenir. »\\
\zh{这个旧电视还能看。} \emph{Zhège jiù diànshì hái néng kàn.} « Cette vieille télévision marche encore. »
\end{exemplebox}

\courssub{Écriture et mémorisation des caractères}

\begin{aretenir}[新 xīn — 13 traits — Clé : 斤 (jīn, « hache »)]
Structure : \textbf{gauche-droite} (\zh{亲} + \zh{斤}).
Ordre des traits (gauche → droite) :
\begin{enumerate}
\item Partie gauche \zh{亲} (qīn, 9 traits) : haut → bas. Écrire le haut (variante de \zh{立} : 丿 一 丨 丿 ㇏), puis \zh{木} (一 丨 一 一).
\item Partie droite \zh{斤} (jīn, 4 traits) : 丿 (trait tombant), 一 (horizontal), 一 (horizontal), 丨 (vertical).
\end{enumerate}
\emph{Astuce :} Du bois (\zh{木}) fraîchement coupé à la hache (\zh{斤}) est \textbf{neuf}.
\end{aretenir}

\begin{aretenir}[旧 jiù — 5 traits — Clé : 日 (rì, « soleil/jour »)]
Structure : \textbf{gauche-droite} (trait vertical + \zh{日}).
Ordre des traits :
\begin{enumerate}
\item Trait vertical \zh{丨} à gauche (1 trait).
\item Clé \zh{日} (4 traits) : vertical gauche \zh{丨}, horizontal-crochet \zh{(trait brisé)}, horizontal du milieu \zh{一}, horizontal de fermeture \zh{一}.
\end{enumerate}
\emph{Attention :} Ne pas confondre avec \zh{日} (rì, 4 traits) ni avec \zh{白} (bái, « blanc », 5 traits, trait horizontal initial).
\end{aretenir}

\courssub{Structure 1 : Attribut (épithète) — 新/旧 + 的 + Nom}

\begin{propriete}[Règle de l'attribut adjectival]
Quand \zh{新} ou \zh{旧} qualifie un nom placé \textbf{après eux}, la structure standard est :\\[4pt]
\centering\textbf{新 / 旧 + 的 + Nom}\\[4pt]
\raggedright
Le mot-outil \zh{的} (de) marque la relation de modification. Comme \zh{新} et \zh{旧} sont monosyllabiques, \zh{的} est \textbf{souvent omis} dans l'usage courant pour former des mots composés figés, mais il \textbf{redevient obligatoire} si l'adjectif est modifié (ex: \zh{很新的}) ou pour insister sur la qualité.
\end{propriete}

\begin{center}
\begin{tabularx}{\linewidth}{|X|X|X|X|}
\hline
\rowcolor{popblueL} Structure & Caractères + Pinyin & Traduction & Note \\
\hline
Adj + 的 + N (standard) & \zh{新的电脑} xīn de diànnǎo & le nouvel ordinateur & Insiste sur la qualité « neuf » \\
\hline
Adj + N (composé, sans 的) & \zh{新电脑} xīn diànnǎo & un nouvel ordi / ordi neuf & Très courant, sens générique \\
\hline
Possessif + Adj + N (sans 的) & \zh{我的新手机} wǒ de xīn shǒujī & mon nouveau téléphone & \textbf{Jamais} \zh{我新手机} ni \zh{我的新的手机} (lourd) \\
\hline
Démonstratif + Adj + N & \zh{这台旧电视} zhè tái jiù diànshì & cette vieille télé & Classificateur \zh{台} requis après démonstratif \\
\hline
\end{tabularx}
\end{center}

\begin{attention}[Piège francophone : la place du possessif]
En français : « \emph{mon} nouveau téléphone » (Possessif + Adj + N).
En chinois : \zh{我 的 新 手机} (Possessif + \zh{的} + Adj + N).
\cle{Ne jamais placer l'adjectif entre le possessif et le nom sans 的}. \zh{我新手机} est faux. \zh{我的新的手机} est possible mais lourd (deux \zh{的}).
\end{attention}

\begin{popfigure}
\begin{tikzpicture}[
  node/.style={rectangle, rounded corners=3pt, draw=popblue, fill=popblueL, align=center, minimum height=0.85cm, inner xsep=6pt, font=\small},
  arrow/.style={-{Stealth[length=2.5mm]}, thick, popdark},
  emph/.style={rectangle, rounded corners=3pt, draw=poporange, fill=poporangeL, align=center, minimum height=0.85cm, inner xsep=6pt, font=\small}
]
\node[node, align=center] (adj) at (0,0){\zh{新 / 旧}\\ xīn / jiù};
\node[node, align=center] (de) at (3,0){\zh{的}\\ de};
\node[node, align=center] (n) at (6,0){\zh{名词}\\ 手机 shǒujī};
\draw[arrow] (adj) -- (de);
\draw[arrow] (de) -- (n);
\node[emph, align=center] at (3,-1.3){\zh{我的新手机}\\ wǒ de xīn shǒujī\\ (Poss + 的 + Adj + N)};
\node[emph, fill=popgreenL, draw=popgreen, align=center] at (7.5,0){\zh{这台旧电视}\\ zhè tái jiù diànshì\\ (Dem + Class + Adj + N)};
\end{tikzpicture}
\legende{Schéma de la structure attributive : l'adjectif précède le nom, \zh{的} facultatif selon le contexte.}
\end{popfigure}

\courssub{Structure 2 : Prédicat — Sujet + 很 + 新/旧}

\begin{propriete}[Adjectif prédicatif : pas de verbe « être »]
Quand \zh{新} ou \zh{旧} exprime une qualité du sujet (fonction de \textbf{prédicat}), il se place \textbf{après le sujet}, sans verbe \zh{是}. On utilise presque toujours l'adverbe de degré \zh{很} (hěn) comme lien (souvent affaibli, ne signifie pas toujours « très »).\\[4pt]
\centering\textbf{Sujet + 很 + 新 / 旧}\\[4pt]
\raggedright
\zh{这部手机很新。} \emph{Zhè bù shǒujī hěn xīn.} « Ce téléphone est neuf. »\\
\zh{那台电脑很旧。} \emph{Nà tái diànnǎo hěn jiù.} « Cet ordinateur est vieux. »
\end{propriete}

\begin{center}
\begin{tabularx}{\linewidth}{|X|X|X|}
\hline
\rowcolor{popblueL} Forme & Exemple (Caractères + Pinyin) & Traduction \\
\hline
Affirmatif & \zh{书很新。} shū hěn xīn & Le livre est neuf. \\
\hline
Négatif (\zh{不} remplace \zh{很}) & \zh{书不旧。} shū bù jiù & Le livre n'est pas vieux. \\
\hline
Insistance (\zh{非常} fēicháng) & \zh{手机非常新。} shǒujī fēicháng xīn & Le téléphone est \emph{très} neuf. \\
\hline
Question & \zh{电脑新吗？} diànnǎo xīn ma? & L'ordi est-il neuf ? \\
\hline
\end{tabularx}
\end{center}

\begin{attention}[Erreur majeure : \zh{是} + Adjectif]
\cle{Interdiction formelle :} \zh{手机是新} \textbf{(FAUX)}.
En chinois, l'adjectif \emph{est} le prédicat (verbe d'état). \zh{是} (shì) ne s'emploie que devant un \textbf{nom} (prédicat nominal) : \zh{这是新手机} (C'est un nouveau téléphone).
\emph{Traduction mentale :} « Le téléphone \textbf{est-neuf} » → \zh{手机 很新}。
\end{attention}

\courssub{Récapitulatif visuel : les deux emplois}

\begin{popfigure}
\begin{tikzpicture}[
  mindmap, concept color=popblueL,
  every node/.style={concept, align=center, font=\small, text=popdark},
  level 1 concept/.append style={sibling angle=120, level distance=3.5cm},
  level 2 concept/.append style={sibling angle=60, level distance=3cm}
]
\node[concept color=poporangeL, align=center]{\textbf{新 / 旧}\\ xīn / jiù}
  child[concept color=popblueL] { node {Attribut (devant N)}
    child { node[concept color=popblue!20, align=center]{\zh{新的书}\\ xīn de shū} }
    child { node[concept color=popblue!20, align=center]{\zh{旧手机}\\ jiù shǒujī} }
    child { node[concept color=popblue!20, align=center]{\zh{我的新号码}\\ wǒ de xīn hàomǎ} }
  }
  child[concept color=popgreenL] { node {Prédicat (après Sujet)}
    child { node[concept color=popgreen!20, align=center]{\zh{书很新。}\\ shū hěn xīn} }
    child { node[concept color=popgreen!20, align=center]{\zh{手机不旧。}\\ shǒujī bù jiù} }
    child { node[concept color=popgreen!20, align=center]{\zh{电视非常旧。}\\ diànshì fēicháng jiù} }
  };
\end{tikzpicture}
\legende{Carte mentale : Attribut vs Prédicat.}
\end{popfigure}

\courssub{Limite d'emploi : 旧 ne qualifie pas les personnes}

\begin{attention}[Erreur fréquente : « 旧人 » = personne âgée ?]
\zh{旧} (jiù) s'applique \textbf{uniquement aux objets/choses}.
\begin{itemize}
\item \textbf{Faux :} \zh{他是一个旧人} (pour « C'est une personne âgée »).
\item \textbf{Correct (âge) :} \zh{老} (lǎo) ↔ \zh{年轻} (niánqīng). Ex: \zh{老人} (personne âgée), \zh{年轻人} (jeune).
\item \textbf{Correct (amitié longue) :} \zh{老朋友} (lǎo péngyou, « vieux ami » = ami de longue date), \textbf{pas} \zh{旧朋友}.
\end{itemize}
\zh{旧人} (jiùrén) existe mais signifie « ancien ami / connaissance d'autrefois / ex-partenaire », avec une nuance souvent mélancolique ou péjorative.
\end{attention}

\begin{savaistu}[Neuf et chance : regards croisés Cameroun / Chine]
En Chine, le \zh{新} (neuf) porte chance. Au Nouvel An (\zh{春节} Chūnjié), on porte des \zh{新衣服} (xīn yīfu), on change les vieilles choses pour « chasser le mauvais sort » (\zh{辞旧迎新} cí jiù yíng xīn : « dire adieu au vieux, accueillir le neuf »). Offrir un objet \zh{旧} (d'occasion) peut être mal vu (manque de respect).
Au Cameroun, offrir un cadeau \textbf{neuf} (emballé, étiqueté) est aussi la norme de respect et de considération. Dans les deux cultures, le « neuf » symbolise l'estime, le renouveau et la prospérité — d'où l'omniprésence de \zh{新} dans la pub tech (nouveau modèle, nouvelle version, nouveau numéro).
\end{savaistu}

\coursec{Le verbe modal 可以 (kěyǐ) — Permission, possibilité, proposition}

\courssub{Sens et emplois fondamentaux}

\begin{definition}[可以 kěyǐ — Verbe modal (auxiliaire modal)]
\zh{可以} exprime :\\[2pt]
1. \textbf{La permission / l'autorisation} (souvent donnée par l'interlocuteur ou la règle) : « pouvoir, avoir le droit de ».\\
2. \textbf{La possibilité objective / la faisabilité} : « il est possible de, on peut ».\\
3. \textbf{La proposition / la suggestion polie} : « veux-tu / pouvons-nous… ? ».
\end{definition}

\begin{exemplebox}[Dialogue : Demander la permission / Proposer]
\zh{— 老师，我可以用手机查单词吗？} \\
\emph{Lǎoshī, wǒ kěyǐ yòng shǒujī chá dāncí ma?} \\
« Monsieur, puis-je utiliser mon téléphone pour chercher le mot ? » (Permission) \\[4pt]
\zh{— 可以，但不可以玩游戏。} \\
\emph{Kěyǐ, dàn bù kěyǐ wán yóuxì.} \\
« Oui (tu peux), mais tu n'as pas le droit de jouer. » (Autorisation / Interdiction) \\[4pt]
\zh{— 这个旧手机可以换新的吗？} \\
\emph{Zhège jiù shǒujī kěyǐ huàn xīn de ma?} \\
« Ce vieux téléphone, peut-on l'échanger contre un neuf ? » (Possibilité objective) \\[4pt]
\zh{— 可以，我们可以去那家店。} \\
\emph{Kěyǐ, wǒmen kěyǐ qù nà jiā diàn.} \\
« Oui, on peut aller dans ce magasin-là. » (Proposition / Faisabilité)
\end{exemplebox}

\courssub{Structures grammaticales de 可以}

\begin{propriete}[Schéma de base : Sujet + 可以 + Verbe (+ Complément)]
\zh{可以} se place \textbf{avant le verbe d'action} (ou l'adjectif dans certains cas). Il ne se conjugue pas.
\end{propriete}

\begin{center}
\begin{tabularx}{\linewidth}{|X|X|X|}
\hline
\rowcolor{popblueL} Type de phrase & Structure & Exemple (Caractères + Pinyin + Traduction) \\
\hline
Affirmative & Suj + \zh{可以} + V (+ Obj) & \zh{你可以去。} nǐ kěyǐ qù. « Tu peux y aller. » \\
\hline
Négative & Suj + \zh{不可以} + V & \zh{你不可以去。} nǐ bù kěyǐ qù. « Tu ne peux pas / Tu n'as pas le droit d'y aller. » \\
\hline
Interrogative (oui/non) & Suj + \zh{可以} + V + \zh{吗}? & \zh{我可以进来吗？} wǒ kěyǐ jìnlái ma? « Puis-je entrer ? » \\
\hline
Interrogative (choix) & Suj + \zh{可以} + V + \zh{不可以}? & \zh{我可以进来不可以？} wǒ kěyǐ jìnlái bù kěyǐ? « Puis-je entrer ou non ? » \\
\hline
Réponse courte & \zh{可以} / \zh{不可以} & \zh{— 可以。} / \zh{— 不可以。} \\
\hline
\end{tabularx}
\end{center}

\begin{attention}[Ne pas confondre : \zh{不可以} vs \zh{不能} vs \zh{不用}]
\begin{itemize}
\item \zh{不可以} (bù kěyǐ) = \textbf{Interdiction / Permission refusée} (règle, autorité, politesse). « Tu n'as \textbf{pas le droit} de… »
\item \zh{不能} (bù néng) = \textbf{Impossibilité physique / technique / capacité}. « Tu \textbf{ne sais pas / ne peux pas physiquement} faire… »
\item \zh{不用} (bù yòng) = \textbf{Absence de nécessité}. « Tu \textbf{n'as pas besoin} de faire… »
\end{itemize}
\emph{Exemple :} \zh{这里不可以抽烟} (Interdiction : « Défense de fumer ici »). \zh{我不能抽烟} (Incapacité : « Je ne sais pas fumer / Je ne peux pas fumer pour raison de santé »). \zh{你不用抽烟} (Inutilité : « Tu n'es pas obligé de fumer »).
\end{attention}

\courssub{可以 vs 想 vs 要 vs 能 — Nuances modales (HSK 3)}

\begin{center}
\begin{tabularx}{\linewidth}{|X|X|X|X|}
\hline
\rowcolor{popblueL} Verbe modal & Sens principal & Nuance & Exemple \\
\hline
\zh{可以} kěyǐ & Permission / Possibilité objective & « C'est permis / C'est possible / Je te le permets » & \zh{你可以坐这儿。} (Tu peux t'asseoir ici.) \\
\hline
\zh{想} xiǎng & Désir / Volonté subjective & « Je veux / J'aimerais » (intention interne) & \zh{我想喝水。} (J'ai envie de boire de l'eau.) \\
\hline
\zh{要} yào & Volonté forte / Projet / Futur proche & « Je vais / Je dois / Je veux fermement » & \zh{我要买新手机。} (Je vais acheter un nouveau tel.) \\
\hline
\zh{能} néng & Capacité / Possibilité matérielle / Savoir-faire & « Je sais faire / C'est matériellement possible » & \zh{我能说中文。} (Je sais parler chinois.) \\
\hline
\end{tabularx}
\end{center}

\begin{methode}[Choisir le bon modal : 3 questions]
Pour traduire « pouvoir / vouloir » en chinois, posez-vous :
\begin{enumerate}
\item \textbf{Est-ce une permission / une règle / une proposition polie ?} → \zh{可以}.
\item \textbf{Est-ce un désir personnel, une envie ?} → \zh{想}.
\item \textbf{Est-ce une capacité, un savoir-faire, une possibilité matérielle (temps, place, argent) ?} → \zh{能}.
\item \textbf{Est-ce une décision ferme, un projet, une obligation ?} → \zh{要}.
\end{enumerate}
\end{methode}

\courssub{可以 dans la politesse et les formules types}

\begin{exemplebox}[Formules de politesse indispensables (niveau HSK 2-3)]
\zh{可以吗？} \emph{Kěyǐ ma?} « Est-ce possible / Ça va ? / C'est d'accord ? » (Pour valider une proposition).\\
\zh{这可以。} \emph{Zhè kěyǐ.} « Ça marche / C'est bon / C'est acceptable. »\\
\zh{不可以。} \emph{Bù kěyǐ.} « Non / Pas question / Interdit. » (Ferme).\\
\zh{请问，可以借用一下吗？} \emph{Qǐngwèn, kěyǐ jièyòng yíxià ma?} « Excusez-moi, puis-je emprunter (utiliser) un instant ? »\\
\zh{你可以试试看。} \emph{Nǐ kěyǐ shì shi kàn.} « Tu peux essayer voir / Essaie donc. » (\zh{试试看} = essayer pour voir le résultat).
\end{exemplebox}

\coursec{Exercices d'application et de transformation}

\begin{exoresolu}[Transformation : Attribut $\leftrightarrow$ Prédicat (新/旧) + Modal 可以]
Énoncé : Transformez les phrases selon le modèle indiqué.
\begin{enumerate}
\item (Attribut $\to$ Prédicat) \zh{这是一本新书。} $\to$ \zh{这本书\underline{\hspace{2cm}}。}
\item (Prédicat $\to$ Attribut avec possessif) \zh{他的电脑很旧。} $\to$ \zh{\underline{\hspace{2cm}}电脑很旧。}
\item (Ajouter 可以 : permission) \zh{我用你的手机。} $\to$ \zh{\underline{\hspace{2cm}}你的手机？}
\item (Négation interdiction) \zh{学生可以带手机上课。} $\to$ \zh{学生\underline{\hspace{2cm}}带手机上课。}
\end{enumerate}
\tcblower
Solution détaillée :
\begin{enumerate}
\item \zh{这本书很新。} (Prédicat : Sujet + 很 + Adj. Pas de \zh{是}).
\item \zh{他的旧电脑很旧。} ou plus naturel : \zh{他的那台旧电脑很旧。} (Attribut : Possessif + Adj + N, sans \zh{的} entre Adj et N). \emph{Note :} L'exercice demande « Attribut avec possessif », la phrase source a déjà le possessif. La transformation attendue est probablement : \zh{那是他的旧电脑。} (Prédicat nominal avec \zh{是}) ou garder \zh{他的旧电脑很旧。} (Attribut + Prédicat). \textbf{Correction pédagogique :} Transformons en phrase nominale : \zh{那是他的旧电脑。} (C'est son vieil ordi).
\item \zh{我可以用你的手机吗？} (Permission : Suj + 可以 + V + Obj + \zh{吗}).
\item \zh{学生不可以带手机上课。} (Interdiction : \zh{不可以}).
\end{enumerate}
\end{exoresolu}

\begin{exercice}[Entraînement complet — À rendre / Travail en classe]
\textbf{A. Traduisez en chinois (caractères + pinyin).}
1. Mon vieux téléphone est cassé, je veux en acheter un neuf.\\
2. Ce nouveau numéro ne marche pas. Puis-je utiliser l'ancien ?\\
3. Les élèves n'ont pas le droit d'utiliser l'ordinateur du prof.\\
4. Ce n'est pas un vieil ordinateur, c'est un nouvel ordinateur.\\
5. Peut-on recycler les vieux téléphones ici ? Oui, c'est possible.

\textbf{B. Corrigez les erreurs (une par phrase).}
1. \zh{我的旧的手机坏了。} \\
2. \zh{这个手机是很新。} \\
3. \zh{我不可以会说中文。} (Contexte : capacité) \\
4. \zh{你要不可以进来？} (Contexte : permission polie) \\
5. \zh{这是一台旧电脑，那台是新的电脑。} (Style)

\textbf{C. Expression écrite (5-6 phrases).}
Rédigez un court dialogue (4-6 répliques) entre un client et un vendeur dans une boutique de téléphonie à Douala. Utilisez : \zh{新}, \zh{旧}, \zh{可以}, \zh{换} (échanger), \zh{修} (réparer), \zh{贵} (cher), \zh{便宜} (pas cher).
\end{exercice}

\begin{corrige}[Exercice « Entraînement complet »]
\textbf{A. Traduction}
1. \zh{我的旧手机坏了，我想买一个新的。} \emph{Wǒ de jiù shǒujī huài le, wǒ xiǎng mǎi yī ge xīn de.} (\zh{新的} = nouveau (téléphone), pronominalisation par \zh{的}).\\
2. \zh{这个新号码不行。我可以用旧号码吗？} \emph{Zhège xīn hàomǎ bù xíng. Wǒ kěyǐ yòng jiù hàomǎ ma?} (\zh{不行} = ne marche pas / ne va pas).\\
3. \zh{学生不可以用老师的电脑。} \emph{Xuésheng bù kěyǐ yòng lǎoshī de diànnǎo.} (\zh{不可以} = interdiction).\\
4. \zh{这不是旧电脑，这是新电脑。} \emph{Zhè bù shì jiù diànnǎo, zhè shì xīn diànnǎo.} (Prédicat nominal $\to$ \zh{是}).\\
5. \zh{这里旧手机可以回收吗？} \emph{Zhèlǐ jiù shǒujī kěyǐ huíshōu ma?} — \zh{可以，可以回收。} \emph{Kěyǐ, kěyǐ huíshōu.}

\textbf{B. Corrections}
1. \zh{我的旧手机坏了。} (Possessif + Adj + N sans \zh{的}).
2. \zh{这个手机很新。} (Prédicat adj $\to$ \zh{很}, pas \zh{是}).
3. \zh{我不会说中文。} ou \zh{我不能说中文。} (\zh{可以} = permission ; capacité $\to$ \zh{能} ou \zh{会}). \emph{Correction attendue :} \zh{我不能说中文。} (impossibilité actuelle) ou \zh{我不会说中文。} (ne pas savoir).
4. \zh{请问，我可以进来吗？} (\zh{要不可以} n'existe pas ; question permission $\to$ \zh{可以...吗?}).
5. \zh{这是一台旧电脑，那台是新的。} (Éviter répétition \zh{电脑} ; \zh{新的} pronominalisé).

\textbf{C. Proposition de dialogue}
\zh{— 老板，这部新手机多少钱？} \\
\zh{— 三千元。您的旧手机可以换五百元。} \\
\zh{— 太贵了。我的旧手机只是屏幕坏了，可以修吗？} \\
\zh{— 可以修，但很贵。买新的更划算。} \\
\zh{— 好吧，给我这部新的吧。}
\emph{Traduction : — Patron, combien ce nouveau téléphone ? — 3000 yuans. Votre vieux téléphone peut être repris 500. — Trop cher. Mon vieux tel a juste l'écran cassé, peut-on le réparer ? — On peut réparer, mais c'est cher. Acheter du neuf est plus rentable. — Bon, je prends ce neuf.}
\end{corrige}

\begin{experience}[Activité orale : « Marché de l'occasion » (Jeu de rôle)]
\textbf{Contexte :} Marché Mokolo (Yaoundé) ou Marché Central (Douala). Stand de téléphones d'occasion.
\textbf{Rôles :} A = Vendeur (vanne les mérites du \zh{新} ou du \zh{旧} reconditionné). B = Client (négocie, demande permission d'essayer, vérifie l'état).
\textbf{Contraintes linguistiques :} Utiliser au moins 3 fois \zh{新/旧} (attribut + prédicat), 2 fois \zh{可以} (permission/possibilité), 1 fois \zh{不可以} (refus/limite), 1 comparatif \zh{比} (révision).
\textbf{Exemple de début :}
\zh{A: 同学，这部手机很新，屏幕没坏。你可以试试。} \\
\zh{B: 真的很新吗？我看看…… 这个旧电池可以换吗？} \\
\zh{A: 电池不可以换，但充电很快。} \\
\zh{B: 可以便宜一点吗？}
\end{experience}

\coursec{Le verbe modal 可以}

Dans la section précédente, nous avons appris à opposer \zh{新} (xīn, « nouveau ») et \zh{旧} (jiù, « vieux, ancien ») pour décrire les objets : un nouveau téléphone \zh{新手机}, un vieil ordinateur \zh{旧电脑}. Mais posséder un objet ne suffit pas : encore faut-il avoir le \emph{droit} de l'utiliser ! Au lycée, peut-on utiliser son téléphone portable en classe ? À la maison, peut-on regarder la télévision tard le soir ? Pour poser ces questions et y répondre, le chinois utilise un petit mot essentiel : le verbe modal \zh{可以} (kěyǐ), « pouvoir ». C'est l'objet de cette section.

\courssub{Observation guidée}

\begin{textebox}[À la maison, le soir]
小明：妈妈，我可以用电脑吗？\\
Xiǎo Míng : Māma, wǒ kěyǐ yòng diànnǎo ma ?\\
« Maman, je peux utiliser l'ordinateur ? »\\
妈妈：可以，但是你不可以玩太久。\\
Māma : Kěyǐ, dànshì nǐ bù kěyǐ wán tài jiǔ.\\
« Oui, mais tu ne peux pas jouer trop longtemps. »\\
小明：那我可以看电视吗？\\
Xiǎo Míng : Nà wǒ kěyǐ kàn diànshì ma ?\\
« Alors, je peux regarder la télévision ? »\\
妈妈：不可以！你要先做作业。\\
Māma : Bù kěyǐ ! Nǐ yào xiān zuò zuòyè.\\
« Non ! Tu dois d'abord faire tes devoirs. »
\end{textebox}

Observons attentivement ce dialogue :

\begin{itemize}
\item \zh{可以} se trouve toujours \textbf{entre le sujet et le verbe} : \zh{我} + \zh{可以} + \zh{用} (je + pouvoir + utiliser).
\item Pour dire « non », on ajoute simplement \zh{不} devant : \zh{不可以}.
\item Pour poser une question, on ajoute \zh{吗} à la fin : \zh{我可以用电脑吗？}
\item Pour répondre, on reprend seulement \zh{可以} ou \zh{不可以}, sans répéter le verbe.
\end{itemize}

\courssub{Définition et nature de 可以}

\begin{definition}[可以 kěyǐ]
\zh{可以} (kěyǐ) est un \cle{verbe modal} qui signifie « pouvoir » au sens de la \textbf{permission} (avoir le droit de) ou de la \textbf{possibilité} (il est possible de). Comme tous les verbes modaux chinois (\zh{会}, \zh{想}, \zh{要}, \zh{能}), il ne s'emploie \textbf{jamais seul} dans une phrase complète : il est toujours suivi d'un verbe principal (sauf dans la réponse brève, où le verbe est sous-entendu). Il est invariable : pas de conjugaison, pas de marque de temps.
\end{definition}

\begin{propriete}[Place de 可以 dans la phrase]
\zh{可以} se place \textbf{avant le verbe principal}, juste après le sujet (ou après le complément de temps s'il y en a un). Il ne change jamais de forme, quelle que soit la personne.

Schéma affirmatif : \textbf{Sujet + 可以 + Verbe + (complément)}

Exemple : \zh{他可以上网。} (Tā kěyǐ shàngwǎng.) « Il peut aller sur Internet. »

En français, « pouvoir » se conjugue (je peux, tu peux, il peut...) ; en chinois, \zh{可以} reste identique pour toutes les personnes : \zh{我可以}, \zh{你可以}, \zh{他可以}, \zh{我们可以}.
\end{propriete}

\begin{exemplebox}[Forme affirmative]
\zh{我可以用电脑。} (Wǒ kěyǐ yòng diànnǎo.) « Je peux utiliser l'ordinateur. »\\
\zh{你可以坐这里。} (Nǐ kěyǐ zuò zhèlǐ.) « Tu peux t'asseoir ici. »\\
\zh{我们可以去图书馆。} (Wǒmen kěyǐ qù túshūguǎn.) « Nous pouvons aller à la bibliothèque. »\\
\zh{学生可以在网上学习汉语。} (Xuésheng kěyǐ zài wǎngshang xuéxí Hànyǔ.) « Les élèves peuvent apprendre le chinois sur Internet. »\\
\zh{明天你可以晚一点儿来。} (Míngtiān nǐ kěyǐ wǎn yìdiǎnr lái.) « Demain, tu peux venir un peu plus tard. »
\end{exemplebox}

\courssub{La négation : 不可以}

\begin{propriete}[Forme négative]
La négation se construit avec \textbf{不 + 可以 + Verbe} : \zh{不可以} signifie « ne pas pouvoir » au sens de « ne pas avoir le droit », « il est interdit de ».

Schéma : \textbf{Sujet + 不可以 + Verbe + (complément)}

Exemple : \zh{你不可以看电视。} (Nǐ bù kěyǐ kàn diànshì.) « Tu ne peux pas regarder la télévision. »
\end{propriete}

\begin{exemplebox}[Forme négative]
\zh{上课的时候不可以打电话。} (Shàngkè de shíhou bù kěyǐ dǎ diànhuà.) « Pendant le cours, on ne peut pas téléphoner. »\\
\zh{这里不可以停车。} (Zhèlǐ bù kěyǐ tíngchē.) « Ici, il est interdit de se garer. »\\
\zh{你不可以太晚回家。} (Nǐ bù kěyǐ tài wǎn huí jiā.) « Tu ne peux pas rentrer trop tard à la maison. »\\
\zh{图书馆里不可以大声说话。} (Túshūguǎn lǐ bù kěyǐ dàshēng shuōhuà.) « Dans la bibliothèque, on ne peut pas parler fort. »
\end{exemplebox}

\begin{aretenir}[不可以 ou 不能 ?]
Pour l'examen, il faut distinguer deux négations proches :
\begin{itemize}
\item \zh{不可以} = \textbf{interdiction}, « tu n'as pas le droit » : \zh{你不可以抽烟。} (Nǐ bù kěyǐ chōuyān.) « Tu n'as pas le droit de fumer. »
\item \zh{不能} = \textbf{impossibilité} (physique, matérielle), « ce n'est pas possible » : \zh{他病了，不能来上课。} (Tā bìng le, bù néng lái shàngkè.) « Il est malade, il ne peut pas venir en cours. »
\end{itemize}
En français, les deux se traduisent par « ne pas pouvoir » ; en chinois, le choix dépend du sens : droit refusé (\zh{不可以}) ou impossibilité (\zh{不能}).
\end{aretenir}

\courssub{La question avec 吗 et les réponses brèves}

\begin{propriete}[Forme interrogative]
Pour demander la permission, on ajoute la particule interrogative \zh{吗} à la fin de la phrase affirmative, sans changer l'ordre des mots :

Schéma : \textbf{Sujet + 可以 + Verbe + 吗 ?}

Exemple : \zh{我可以进来吗？} (Wǒ kěyǐ jìnlái ma ?) « Je peux entrer ? »

Pour répondre, on utilise des \textbf{réponses brèves} sans répéter le verbe principal :
\begin{itemize}
\item Oui : \zh{可以。} (Kěyǐ.) « Oui. »
\item Non : \zh{不可以。} (Bù kěyǐ.) « Non. »
\end{itemize}
\end{propriete}

\begin{exemplebox}[Questions et réponses]
— \zh{学生可以用手机吗？} (Xuésheng kěyǐ yòng shǒujī ma ?) « Les élèves peuvent-ils utiliser le téléphone portable ? »\\
— \zh{不可以。} (Bù kěyǐ.) « Non. »

— \zh{老师，我可以问一个问题吗？} (Lǎoshī, wǒ kěyǐ wèn yí ge wèntí ma ?) « Monsieur le professeur, je peux poser une question ? »\\
— \zh{可以，请问。} (Kěyǐ, qǐng wèn.) « Oui, je t'en prie. »

— \zh{我们可以现在回家吗？} (Wǒmen kěyǐ xiànzài huí jiā ma ?) « Nous pouvons rentrer à la maison maintenant ? »\\
— \zh{可以。} (Kěyǐ.) « Oui. »
\end{exemplebox}

\begin{savaistu}[La question affirmative-négative]
À un niveau plus avancé, il existe une autre façon de poser la question, sans \zh{吗} : on répète le modal en forme affirmative-négative : \zh{你可以不可以来？} (Nǐ kěyǐ bù kěyǐ lái ?) « Peux-tu venir, oui ou non ? » C'est la forme « V-不-V » appliquée au verbe modal. Retenez-la comme complément culturel ; à l'examen de Première, la forme avec \zh{吗} suffit.
\end{savaistu}

\begin{popfigure}
\begin{tikzpicture}[font=\small, node distance=0.5cm]
\node[draw=popblue, fill=popblueL, rounded corners, minimum width=2.2cm, minimum height=0.9cm, align=center] (s) {Sujet\\我 / 你 / 他};
\node[draw=poporange, fill=poporangeL, rounded corners, minimum width=2.2cm, minimum height=0.9cm, align=center, right=of s] (m) {可以\\(不 + 可以)};
\node[draw=popgreen, fill=popgreenL, rounded corners, minimum width=2.2cm, minimum height=0.9cm, align=center, right=of m] (v) {Verbe\\用 / 看 / 去};
\node[draw=poppurple, fill=poppurpleL, rounded corners, minimum width=2.2cm, minimum height=0.9cm, align=center, right=of v] (q) {(吗 ?)\\。};
\draw[-{Stealth}, thick, popdark] (s) -- (m);
\draw[-{Stealth}, thick, popdark] (m) -- (v);
\draw[-{Stealth}, thick, popdark] (v) -- (q);
\node[below=0.35cm of m, align=center, text=popmuted] {invariable,\\jamais seul};
\end{tikzpicture}
\legende{Schéma de la structure : Sujet + 可以 + Verbe (+ 吗 ?)}
\end{popfigure}

\begin{popfigure}
\begin{tikzpicture}[font=\small, node distance=0.6cm and 1.6cm]
\node[draw=popblue, fill=popblueL, rounded corners, align=center, minimum width=4.5cm, minimum height=0.9cm] (q) {Question : 可以 + Verbe + 吗 ?\\我可以进来吗？};
\node[draw=popgreen, fill=popgreenL, rounded corners, align=center, minimum width=3cm, minimum height=0.9cm, below left=of q] (o) {Oui\\可以。};
\node[draw=poppink, fill=poppinkL, rounded corners, align=center, minimum width=3cm, minimum height=0.9cm, below right=of q] (n) {Non\\不可以。};
\draw[-{Stealth}, thick, popdark] (q) -- (o);
\draw[-{Stealth}, thick, popdark] (q) -- (n);
\end{tikzpicture}
\legende{Organigramme de réponse à une question en 可以...吗}
\end{popfigure}

\courssub{Tableaux récapitulatifs}

\begin{center}
\begin{tabularx}{\linewidth}{|X|X|X|}
\hline
\rowcolor{popblueL} Forme & Exemple (caractères + pinyin) & Traduction \\
\hline
Affirmative & 我可以用电脑。Wǒ kěyǐ yòng diànnǎo. & Je peux utiliser l'ordinateur. \\
\hline
Négative & 你不可以看电视。Nǐ bù kěyǐ kàn diànshì. & Tu ne peux pas regarder la télé. \\
\hline
Interrogative & 我可以进来吗？Wǒ kěyǐ jìnlái ma ? & Je peux entrer ? \\
\hline
Réponse positive & 可以。Kěyǐ. & Oui. \\
\hline
Réponse négative & 不可以。Bù kěyǐ. & Non. \\
\hline
Affirmative-négative & 你可以不可以来？Nǐ kěyǐ bù kěyǐ lái ? & Peux-tu venir ? \\
\hline
\end{tabularx}
\end{center}

\begin{center}
\begin{tabularx}{\linewidth}{|X|X|X|X|}
\hline
\rowcolor{popblueL} Caractères & Pinyin & Nature & Traduction \\
\hline
可以 & kěyǐ & verbe modal & pouvoir (permission) \\
\hline
用 & yòng & verbe & utiliser \\
\hline
上网 & shàngwǎng & verbe & aller sur Internet \\
\hline
进来 & jìnlái & verbe & entrer \\
\hline
手机 & shǒujī & nom & téléphone portable \\
\hline
电脑 & diànnǎo & nom & ordinateur \\
\hline
问题 & wèntí & nom & question, problème \\
\hline
停车 & tíngchē & verbe & se garer \\
\hline
\end{tabularx}
\end{center}

\courssub{Comparaison avec le français et pièges des francophones}

En français, « pouvoir » se conjugue et la question peut se faire par inversion (« Puis-je... ? ») ; en chinois, la structure est beaucoup plus régulière, mais elle impose une discipline stricte :

\begin{itemize}
\item \textbf{Un seul verbe principal} après \zh{可以} (suivi de ses compléments).
\item \textbf{Pas de conjugaison} : \zh{可以} ne change jamais.
\item \textbf{Ordre fixe} : Sujet → \zh{可以} → Verbe → \zh{吗}.
\end{itemize}

\begin{attention}[Erreurs fréquentes des francophones]
\begin{itemize}
\item Ne jamais placer \zh{吗} au milieu de la phrase : il se met \textbf{à la toute fin}. Correct : \zh{我可以用吗？} (Wǒ kěyǐ yòng ma ?) « Je peux l'utiliser ? »
\item Ne pas confondre \zh{可以} (permission : « avoir le droit ») et \zh{会} (capacité acquise : « savoir faire »). \zh{我会说汉语。} (Wǒ huì shuō Hànyǔ.) « Je sais parler chinois. » / \zh{我可以说汉语吗？} (Wǒ kěyǐ shuō Hànyǔ ma ?) « Puis-je parler chinois ? » (ai-je le droit ?).
\item Ne pas traduire mot à mot « Est-ce que je peux... ? » par une inversion : en chinois, l'ordre des mots ne change jamais à la question ; on ajoute seulement \zh{吗}.
\item Dans la réponse brève, ne pas répéter le verbe : on répond \zh{可以。} et non \zh{可以用。} (forme lourde, à éviter à l'examen).
\item Attention au complément de temps : il se place avant le verbe, souvent avant \zh{可以} : \zh{我今天晚上可以看电视。} et non après le verbe comme en français.
\end{itemize}
\end{attention}

\begin{methode}[Répondre à une question en 可以...吗]
\begin{enumerate}
\item \textbf{Identifier} la demande de permission : repérer \zh{可以} et le verbe principal.
\item \textbf{Choisir} la réponse brève : \zh{可以。} (oui) ou \zh{不可以。} (non).
\item \textbf{Justifier} avec une phrase complète en reprenant la structure : \zh{你可以用电脑，但是不可以玩游戏。} (Nǐ kěyǐ yòng diànnǎo, dànshì bù kěyǐ wán yóuxì.) « Tu peux utiliser l'ordinateur, mais tu ne peux pas jouer. »
\item \textbf{Vérifier} l'ordre des mots : Sujet + (不)可以 + Verbe + complément.
\end{enumerate}
\end{methode}

\begin{exoresolu}[Demander et donner la permission]
Énoncé : Traduisez en chinois, puis répondez par une réponse brève négative justifiée : « Maman, je peux regarder la télévision ce soir ? »
\tcblower
Solution détaillée :
\begin{itemize}
\item Sujet : \zh{我} ; modal : \zh{可以} ; verbe : \zh{看} ; complément : \zh{电视} ; temps : \zh{今天晚上} ; particule : \zh{吗}.
\item Question : \zh{妈妈，我今天晚上可以看电视吗？} (Māma, wǒ jīntiān wǎnshang kěyǐ kàn diànshì ma ?)
\item Réponse brève : \zh{不可以。} (Bù kěyǐ.) « Non. »
\item Justification : \zh{你明天有考试，要先学习。} (Nǐ míngtiān yǒu kǎoshì, yào xiān xuéxí.) « Tu as un examen demain, tu dois d'abord étudier. »
\end{itemize}
Points de vigilance : \zh{吗} se place à la fin ; le complément de temps \zh{今天晚上} se place avant \zh{可以} ; la réponse brève ne répète pas \zh{看电视}.
\end{exoresolu}

\begin{exercice}[À vous de jouer]
\begin{enumerate}
\item Traduisez : « Les élèves peuvent-ils utiliser le téléphone portable en classe ? » — « Non. »
\item Mettez à la forme négative : \zh{他可以上网。}
\item Construisez une question avec \zh{可以...吗} pour demander la permission d'entrer dans le bureau des professeurs.
\end{enumerate}
\end{exercice}

\begin{corrige}[Exercice « À vous de jouer »]
\begin{enumerate}
\item \zh{学生上课可以用手机吗？} (Xuésheng shàngkè kěyǐ yòng shǒujī ma ?) — \zh{不可以。} (Bù kěyǐ.)
\item \zh{他不可以上网。} (Tā bù kěyǐ shàngwǎng.) « Il ne peut pas aller sur Internet. »
\item \zh{老师，我可以进办公室吗？} (Lǎoshī, wǒ kěyǐ jìn bàngōngshì ma ?) « Monsieur, je peux entrer dans le bureau ? »
\end{enumerate}
\end{corrige}

\coursec{Leçon 36 — Grammaire : antonyme 新 / 旧 ; verbe modal 可以}

\courssub{Observation guidée : dialogue au magasin de téléphones}

\begin{textebox}[Dialogue : Acheter un téléphone à Yaoundé]
\textbf{Contexte :} Dans une boutique de téléphonie au quartier Mokolo (Yaoundé), un élève cherche un téléphone.
\medskip
\noindent \zh{售货员：欢迎光临！您想买新的手机还是旧的手机？} \\
\emph{Shòuhuòyuán: Huānyíng guānglín! Nín xiǎng mǎi xīn de shǒujī háishì jiù de shǒujī?} \\
— Vendeur : Bienvenue ! Voulez-vous acheter un téléphone neuf ou d'occasion ?
\medskip
\noindent \zh{学生：我想买新的。这个可以试用吗？} \\
\emph{Xuéshēng: Wǒ xiǎng mǎi xīn de. Zhège kěyǐ shìyòng ma?} \\
— Élève : Je veux en acheter un neuf. Puis-je essayer celui-ci ?
\medskip
\noindent \zh{售货员：可以。但旧的手机不可以试用，只能看。} \\
\emph{Shòuhuòyuán: Kěyǐ. Dàn jiù de shǒujī bù kěyǐ shìyòng, zhǐ néng kàn.} \\
— Vendeur : C'est possible. Mais les téléphones d'occasion ne peuvent pas être essayés, on peut seulement les regarder.
\medskip
\noindent \zh{学生：好的。那我可以用微信付款吗？} \\
\emph{Xuéshēng: Hǎo de. Nà wǒ kěyǐ yòng Wēixìn fùkuǎn ma?} \\
— Élève : D'accord. Puis-je payer avec WeChat ?
\medskip
\noindent \zh{售货员：可以。请扫这个码。} \\
\emph{Shòuhuòyuán: Kěyǐ. Qǐng sǎo zhège mǎ.} \\
— Vendeur : Oui. Veuillez scanner ce code.
\end{textebox}

\begin{experience}[Activité orale : Jeu de rôle « Au magasin »]
En binôme : l'un joue le vendeur, l'autre le client.
\begin{enumerate}
    \item Le client demande s'il y a des téléphones neufs (\zh{新的}) et d'occasion (\zh{旧的}).
    \item Le client demande la permission d'essayer (\zh{可以试用吗？}), de regarder (\zh{可以看吗？}), de payer par Mobile Money / WeChat (\zh{可以用...付款吗？}).
    \item Le vendeur répond par \zh{可以} ou \zh{不可以} + verbe (ex: \zh{不可以试用，只能看}).
    \item Échangez les rôles. Utilisez le vocabulaire du dialogue et les structures de la leçon.
\end{enumerate}
\end{experience}

\courssub{Les antonymes 新 / 旧}

\begin{definition}[Nouveau vs Vieux / Ancien]
\cle{新} (xīn) : nouveau, neuf, récent. S'emploie pour des objets, des idées, des relations, des vêtements, etc.
\cle{旧} (jiù) : vieux, ancien, usagé, d'occasion. S'emploie pour des objets, des lieux, des habitudes, des vêtements, des livres, etc.
\medskip
\noindent \textbf{Attention :} Pour les \textbf{personnes}, on n'utilise pas \cle{旧}. Voir piège n°5 ci-dessous.
\end{definition}

\begin{propriete}[Règle de position des adjectifs qualificatifs en chinois]
En chinois, les adjectifs (verbes d'état) placés en fonction d'épithète \textbf{précèdent toujours le nom} qu'ils qualifient et sont généralement reliés à ce nom par la particule structurelle \cle{的} (de).
\medskip
\noindent Schéma : \textbf{Adjectif + 的 + Nom}
\medskip
\noindent Cela vaut pour \zh{新} (xīn) et \zh{旧} (jiù) comme pour tous les adjectifs monosyllabiques ou bisyllabiques (niveau HSK 1-3).
\end{propriete}

\begin{exemplebox}[Position correcte de 新 et 旧]
\begin{itemize}
    \item \zh{新的手机} \emph{xīn de shǒujī} — un nouveau téléphone (neuf).
    \item \zh{旧的书包} \emph{jiù de shūbāo} — un vieux cartable / un cartable d'occasion.
    \item \zh{这是我的新电脑。} \emph{Zhè shì wǒ de xīn diànnǎo.} — C'est mon nouvel ordinateur.
    \item \zh{那是他的旧自行车。} \emph{Nà shì tā de jiù zìxíngchē.} — C'est son vieux vélo.
    \item \zh{新的朋友} \emph{xīn de péngyou} — un nouvel ami (récemment connu).
    \item \zh{旧的衣服} \emph{jiù de yīfu} — de vieux vêtements / des habits usagés.
\end{itemize}
\end{exemplebox}

\begin{center}
\begin{tabularx}{\linewidth}{|X|X|X|X|}
\hline
\rowcolor{popblueL} \textbf{Caractères} & \textbf{Pinyin} & \textbf{Nature} & \textbf{Traduction} \\
\hline
新 & xīn & adj. & nouveau, neuf \\
\hline
旧 & jiù & adj. & vieux, ancien, d'occasion \\
\hline
手机 & shǒujī & n. & téléphone portable \\
\hline
电脑 & diànnǎo & n. & ordinateur \\
\hline
书包 & shūbāo & n. & cartable, sac à dos \\
\hline
自行车 & zìxíngchē & n. & vélo \\
\hline
衣服 & yīfu & n. & vêtements \\
\hline
朋友 & péngyou & n. & ami \\
\hline
\end{tabularx}
\end{center}

\begin{attention}[Piège n°1 : Oublier la particule 的]
En français, l'adjectif « nouveau » se place souvent \textbf{avant} le nom sans mot-outil (\emph{un nouveau téléphone}). L'apprenant francophone a tendance à dire : \zh{*我新手机} (\emph{*wǒ xīn shǒujī}).
\medskip
\noindent \textbf{Correction :} \zh{我的新手机} (\emph{wǒ de xīn de shǒujī}) ou \zh{我新的手机} (\emph{wǒ xīn de shǒujī}). La particule \cle{的} est obligatoire entre l'adjectif et le nom (sauf cas rares d'adjectifs monosyllabiques très lexicalisés, mais en HSK 1-3, on met toujours \cle{的} pour sécurité).
\end{attention}

\begin{attention}[Piège n°2 : Confondre 旧 et 老 pour les personnes]
\cle{旧} (jiù) s'emploie pour les \textbf{objets}, les \textbf{lieux}, les \textbf{idées}, les \textbf{vêtements} (usagé, ancien).
\cle{老} (lǎo) s'emploie pour les \textbf{personnes} (âgé, senior) ou dans des composés figés (\zh{老师} lǎoshī, \zh{老虎} lǎohǔ, \zh{老板} lǎobǎn).
\medskip
\noindent \textbf{Faux :} \zh{*我旧的老师} (\emph{*wǒ jiù de lǎoshī}) — « mon vieux professeur » (sous-entendu : usagé, dévalorisant).
\noindent \textbf{Juste :} \zh{我的老老师} (\emph{wǒ de lǎo lǎoshī}) — mon ancien professeur (que je connais depuis longtemps / qui est âgé / respectueux).
\noindent \textbf{Juste :} \zh{我的新老师} (\emph{wǒ de xīn lǎoshī}) — mon nouveau professeur (qui vient d'arriver).
\end{attention}

\courssub{Le verbe modal 可以}

\begin{definition}[Permission, possibilité, capacité]
\cle{可以} (kěyǐ) est un \textbf{verbe modal} (助动词, zhùdòngcí). Il exprime principalement la \textbf{permission} (« avoir le droit de »), la \textbf{possibilité} (« c'est possible de ») ou la \textbf{capacité} (« savoir faire / être en mesure de ») selon le contexte.
\end{definition}

\begin{propriete}[Caractéristiques grammaticales de 可以]
\begin{enumerate}
    \item \textbf{Invariable :} pas de conjugaison, pas d'accord en personne ni en nombre (pas de « je peux, tu peux, il peut »).
    \item \textbf{Place fixe :} \textbf{Sujet + 可以 + Verbe d'action (+ Complément)}. Il se place \textbf{immédiatement avant} le verbe principal.
    \item \textbf{Nécessité d'un verbe d'action :} \cle{可以} ne peut pas fonctionner seul comme verbe principal transitif direct (on ne « peut » pas un objet). Il \textbf{exige} un verbe d'action derrière lui (\zh{用} yòng, \zh{买} mǎi, \zh{吃} chī, \zh{看} kàn, \zh{去} qù, \zh{来} lái, \zh{游泳} yóuyǒng…).
    \item \textbf{Négation :} \textbf{不可以} (bù kěyǐ). \textbf{Jamais} \zh{*没可以} (\emph{*méi kěyǐ}).
    \item \textbf{Interrogation :} Ajout de la particule \cle{吗} (ma) en fin de phrase : \zh{可以...吗？}
\end{enumerate}
\end{propriete}

\begin{exemplebox}[Structure correcte avec 可以]
\begin{itemize}
    \item \zh{我可以用这个手机。} \emph{Wǒ kěyǐ yòng zhège shǒujī.} — Je peux / J'ai le droit d'utiliser ce téléphone.
    \item \zh{他可以去北京。} \emph{Tā kěyǐ qù Běijīng.} — Il peut aller à Pékin.
    \item \zh{我们不可以迟到。} \emph{Wǒmen bù kěyǐ chídào.} — Nous ne devons pas / ne pouvons pas arriver en retard.
    \item \zh{你可以吃这个苹果吗？} \emph{Nǐ kěyǐ chī zhège píngguǒ ma?} — Peux-tu / As-tu le droit de manger cette pomme ?
    \item \zh{他们可以游泳。} \emph{Tāmen kěyǐ yóuyǒng.} — Ils savent nager / Ils peuvent nager.
\end{itemize}
\end{exemplebox}

\begin{center}
\begin{tabularx}{\linewidth}{|X|X|X|X|}
\hline
\rowcolor{popblueL} \textbf{Caractères} & \textbf{Pinyin} & \textbf{Nature} & \textbf{Traduction} \\
\hline
可以 & kěyǐ & v. modal & pouvoir, être autorisé à \\
\hline
不可以 & bù kěyǐ & v. modal (nég.) & ne pas pouvoir, ne pas avoir le droit \\
\hline
用 & yòng & v. & utiliser, se servir de \\
\hline
买 & mǎi & v. & acheter \\
\hline
看 & kàn & v. & regarder, lire \\
\hline
去 & qù & v. & aller \\
\hline
来 & lái & v. & venir \\
\hline
吃 & chī & v. & manger \\
\hline
游泳 & yóuyǒng & v. & nager \\
\hline
微信 & Wēixìn & n. & WeChat \\
\hline
付款 & fùkuǎn & v. & payer \\
\hline
试用 & shìyòng & v. & essayer (un appareil) \\
\hline
苹果 & píngguǒ & n. & pomme \\
\hline
房子 & fángzi & n. & maison \\
\hline
慢 & màn & adj. & lent \\
\hline
大 & dà & adj. & grand \\
\hline
\end{tabularx}
\end{center}

\courssub{Écriture des caractères : 新, 旧, 可, 以}

\begin{propriete}[Ordre des traits et clés (radicaux)]
\begin{itemize}
    \item \textbf{新} (xīn, 13 traits) — Clé : \zh{斤} (jīn, « hache », radical n°69). Structure : haut-bas. Écrire \zh{立} (debout) puis \zh{斤}. Ordre : horizontales de \zh{立}, verticale, horizontale basse ; puis \zh{斤} : horizontale, verticale, crochet, horizontale, horizontale.
    \item \textbf{旧} (jiù, 5 traits) — Clé : \zh{丨} (gǔn, « trait vertical », radical n°2). Structure : haut-bas. Écrire le trait vertical \zh{丨} puis \zh{臼} (jiù, « mortier ») : horizontale, verticale, crochet, horizontale, horizontale.
    \item \textbf{可} (kě, 5 traits) — Clé : \zh{口} (kǒu, « bouche », radical n°30). Structure : encadrement. Écrire le cadre \zh{口} (verticale, crochet, horizontale, horizontale de fermeture) PUIS le trait horizontal intérieur (le 5e trait). \textbf{Règle :} l'intérieur s'écrit après le cadre (sauf si le cadre n'est pas fermé en bas).
    \item \textbf{以} (yǐ, 5 traits) — Clé : \zh{人} (rén, « homme », radical n°9) → forme \zh{亻} (rénzìpáng). Structure : gauche-droite. Écrire \zh{亻} (trait oblique gauche, trait oblique droit) puis la partie droite : horizontale, crochet vertical, horizontale.
\end{itemize}
\end{propriete}

\courssub{Élément socioculturel : Téléphonie et marché de l'occasion Cameroun / Chine}

\begin{savaistu}[Le téléphone portable : outil vital des deux côtés]
\begin{itemize}
    \item \textbf{Cameroun :} Le téléphone est le premier outil d'accès à Internet (Mobile Money : Orange Money, MTN MoMo). Les marchés comme \emph{Mokolo} (Yaoundé) ou \emph{Congo} (Douala) regorgent de téléphones \zh{新的} (neufs, souvent garantis) et \zh{旧的} (occasion, \emph{« occasion Belgique »} ou \emph{« occasion UK »}). Marchander (\zh{讲价} jiǎngjià) est la norme pour les \zh{旧的}.
    \item \textbf{Chine :} Pénétration du smartphone quasi totale. Paiement sans espèces roi : \zh{微信支付} (Wēixìn zhīfù) et \zh{支付宝} (Zhīfùbǎo). Marché de l'occasion (\zh{二手市场} èrshǒu shìchǎng) très organisé en ligne (\zh{闲鱼} Xiányú, « Poisson paresseux », équivalent LeBonCoin). Les \zh{旧的} y sont classés par grade (99\% neuf, 95\%, etc.).
    \item \textbf{Point commun :} Dans les deux pays, demander \zh{可以试用吗？} (Puis-je essayer ?) est essentiel avant d'acheter un \zh{旧的} téléphone.
\end{itemize}
\end{savaistu}

\courssub{Comparaison avec le français et pièges des francophones}

\begin{propriete}[Règle de position des adjectifs qualificatifs en chinois (Rappel)]
En chinois, les adjectifs (ou verbes d'état) placés en fonction d'épithète \textbf{précèdent toujours le nom} qu'ils qualifient et sont généralement reliés à ce nom par la particule structurelle \cle{的} (de).
\medskip
\noindent Schéma : \textbf{Adjectif + 的 + Nom}
\medskip
\noindent Cela vaut pour \zh{新} (xīn, nouveau) et \zh{旧} (jiù, vieux/ancien) comme pour tous les adjectifs monosyllabiques ou bisyllabiques.
\end{propriete}

\begin{exemplebox}[Position correcte de 新 et 旧]
\begin{itemize}
    \item \zh{新的手机} \emph{xīn de shǒujī} — un nouveau téléphone (neuf).
    \item \zh{旧的书包} \emph{jiù de shūbāo} — un vieux cartable.
    \item \zh{这是我的新电脑。} \emph{Zhè shì wǒ de xīn de diànnǎo.} — C'est mon nouvel ordinateur.
    \item \zh{那是他的旧自行车。} \emph{Nà shì tā de jiù de zìxíngchē.} — C'est son vieux vélo.
\end{itemize}
\end{exemplebox}

\begin{attention}[Piège n°1 : Oublier la particule 的]
En français, l'adjectif « nouveau » se place souvent \textbf{avant} le nom sans mot-outil. L'apprenant francophone a tendance à dire : \zh{*我新手机} (\emph{*wǒ xīn shǒujī}).
\medskip
\noindent \textbf{Correction :} \zh{我的新手机} (\emph{wǒ de xīn de shǒujī}) ou \zh{我新的手机} (\emph{wǒ xīn de shǒujī}). La particule \cle{的} est obligatoire entre l'adjectif et le nom (sauf cas rares d'adjectifs monosyllabiques très liés au nom, mais en HSK 1-3, on met toujours \cle{的} pour sécurité).
\end{attention}

\begin{attention}[Piège n°2 : Confondre 旧 et 老 pour les personnes]
\cle{旧} (jiù) s'emploie pour les \textbf{objets}, les \textbf{lieux}, les \textbf{idées}, les \textbf{vêtements}, etc. (usagé, ancien).
\cle{老} (lǎo) s'emploie pour les \textbf{personnes} (âgé) ou dans des composés figés (\zh{老师} lǎoshī, \zh{老虎} lǎohǔ).
\medskip
\noindent \textbf{Faux :} \zh{*我旧的老师} (\emph{*wǒ jiù de lǎoshī}) — « mon vieux professeur » (sous-entendu : usagé).
\noindent \textbf{Juste :} \zh{我的老老师} (\emph{wǒ de lǎo lǎoshī}) — mon ancien professeur (que je connais depuis longtemps / qui est âgé).
\end{attention}

\begin{center}
\begin{tabularx}{\linewidth}{|X|X|X|X|}
\hline
\rowcolor{popblueL} \textbf{Français} & \textbf{Calque fautif} & \textbf{Chinois correct} & \textbf{Analyse} \\
\hline
un nouveau téléphone & *新手机 & \zh{新的手机} \emph{xīn de shǒujī} & Adj + \cle{的} + Nom \\
\hline
un vieux livre & *旧书 & \zh{旧的书} \emph{jiù de shū} & Adj + \cle{的} + Nom \\
\hline
mon ancien camarade & *我旧同学 & \zh{我的老同学} \emph{wǒ de lǎo tóngxué} & Personne $\rightarrow$ \cle{老} + \cle{的} \\
\hline
\end{tabularx}
\end{center}

\courssub{Le verbe modal 可以 : invariabilité, place et négation}

\begin{propriete}[Nature et place du modal 可以 (Rappel)]
\cle{可以} (kěyǐ) est un \textbf{verbe modal} (助动词 zhùdòngcí). Il exprime la permission, la possibilité ou la capacité selon le contexte.
\medskip
\noindent \textbf{Caractéristiques majeures :}
\begin{enumerate}
    \item \textbf{Invariable :} pas de conjugaison, pas d'accord.
    \item \textbf{Place fixe :} \textbf{Sujet + 可以 + Verbe d'action (+ Complément)}. Il se place \textbf{avant} le verbe principal.
    \item \textbf{Nécessité d'un verbe d'action :} \cle{可以} ne peut pas fonctionner seul comme verbe principal transitif direct.
    \item \textbf{Négation :} \textbf{不可以} (bù kěyǐ). Jamais \zh{*没可以} (\emph{*méi kěyǐ}).
\end{enumerate}
\end{propriete}

\begin{exemplebox}[Structure correcte avec 可以]
\begin{itemize}
    \item \zh{我可以用这个手机。} \emph{Wǒ kěyǐ yòng zhège shǒujī.} — Je peux utiliser ce téléphone.
    \item \zh{他可以去北京。} \emph{Tā kěyǐ qù Běijīng.} — Il peut aller à Pékin.
    \item \zh{我们不可以迟到。} \emph{Wǒmen bù kěyǐ chídào.} — Nous ne devons pas / ne pouvons pas arriver en retard.
\end{itemize}
\end{exemplebox}

\begin{attention}[Piège n°3 : Placer 可以 après le verbe (calque « je téléphone peux »)]
Structure fautive : \zh{*我用可以手机} (\emph{*wǒ yòng kěyǐ shǒujī}).
\medskip
\noindent \textbf{Explication :} En français, le modal « pouvoir » se conjugue et le verbe principal reste à l'infinitif après (je peux utiliser). En chinois, le modal \textbf{précède} le verbe. Règle immuable : \textbf{Sujet + 可以 + Verbe}.
\end{attention}

\begin{attention}[Piège n°4 : Utiliser 可以 comme verbe transitif direct (calque « je peux ce téléphone »)]
Phrase fautive : \zh{*我可以这个手机} (\emph{*wǒ kěyǐ zhège shǒujī}).
\medskip
\noindent \textbf{Explication :} « Pouvoir » en français peut sembler absorber l'objet dans l'oral familier (« Je peux ce gâteau ? » pour « Je peux \emph{manger} ce gâteau ? »). En chinois, \cle{可以} est \textbf{intransitif} : il \textbf{exige} un verbe d'action derrière lui (\zh{用} yòng, \zh{买} mǎi, \zh{吃} chī, \zh{看} kàn…).
\medskip
\noindent \textbf{Correction :} \zh{我可以用/买这个手机。} (\emph{Wǒ kěyǐ yòng/mǎi zhège shǒujī.})
\end{attention}

\begin{attention}[Piège n°5 : Négation avec 没 au lieu de 不]
Phrase fautive : \zh{*我没可以去} (\emph{*wǒ méi kěyǐ qù}).
\medskip
\noindent \textbf{Règle :} Les verbes modaux (\cle{可以}, \cle{能}, \cle{会}, \cle{要}, \cle{应该}…) se nient \textbf{uniquement avec 不} (bù). \cle{没} (méi) nie les actions passées/accomplies (verbes d'action + \cle{了}/ \cle{过}) ou l'existence (\zh{没有}). Un modal n'exprime pas une action accomplie, donc pas de \cle{没}.
\medskip
\noindent \textbf{Correct :} \zh{我不可以去。} (\emph{Wǒ bù kěyǐ qù.}) — Je ne peux pas y aller / Je n'ai pas le droit d'y aller.
\end{attention}

\begin{attention}[Piège subtil : 可以 seul dans une réponse elliptique]
En contexte oral, on peut répondre \zh{可以} (D'accord / C'est possible) à une demande. Mais \zh{我不可以} seul, sans verbe suivant, sonne incomplet ou enfantin (« Je ne peux pas » — mais quoi ?).
\medskip
\noindent \textbf{Préférer :} \zh{我不可以用。} (\emph{Wǒ bù kěyǐ yòng.}) — Je ne peux pas l'utiliser. \quad ou \quad \zh{我不能去。} (\emph{Wǒ bù néng qù.}) — Je ne peux pas y aller.
\end{attention}

\courssub{Tableau comparatif synthétique : Français vs Chinois}

\begin{center}
\begin{tabularx}{\linewidth}{|X|X|X|}
\hline
\rowcolor{popblueL} \textbf{Point de grammaire} & \textbf{Français (LV1)} & \textbf{Chinois (LV2)} \\
\hline
Adjectif épithète « nouveau/vieux » & Place variable (avant/après nom), pas de mot-outil & \textbf{Toujours avant le nom} + \textbf{的} obligatoire (Adj + 的 + Nom) \\
\hline
Verbe modal « pouvoir » & Se conjugue (je peux, tu peux…), suivi de l'infinitif & \textbf{Invariable} (可以), \textbf{avant} le verbe d'action (Suj + 可以 + V) \\
\hline
Négation de « pouvoir » & « ne ... pas » autour du conjugué (je ne peux pas) & \textbf{不可以} (bù kěyǐ) placé avant le verbe ; \textbf{jamais 没可以} \\
\hline
Objet direct après « pouvoir » & Possible à l'oral familier (« Je peux ce livre ? ») & \textbf{Impossible} : il faut un verbe (看, 买, 用…) après 可以 \\
\hline
« Vieux » (personne) & « vieux / ancien / âgé » selon contexte & \textbf{老} (lǎo) pour personne ; \textbf{旧} (jiù) pour objets/lieux \\
\hline
\end{tabularx}
\end{center}

\begin{popfigure}
\begin{tikzpicture}[
    node distance=1.2cm and 1.5cm,
    box/.style={draw=popdark, thick, rounded corners, fill=popcream, text width=3.2cm, align=center, font=\small},
    arrow/.style={-Stealth, thick, popblue},
    cross/.style={draw=popred, thick, line width=1.5pt},
    check/.style={draw=popgreen, thick, line width=1.5pt}
]
% French column
\node[box] (f1) {Français : « Je peux téléphone »};
\node[box, below=of f1] (f2) {Français : « Mon nouveau téléphone »};
\node[box, below=of f2] (f3) {Français : « Je ne peux pas »};
\node[box, below=of f3] (f4) {Français : « Mon vieux prof »};

% Arrow to Chinese
\node[box, right=of f1, fill=popgreenL, align=center] (c1){Chinois : \zh{我可以用手机} \\ \emph{Wǒ kěyǐ yòng shǒujī}};
\node[box, right=of f2, fill=popgreenL, align=center] (c2){Chinois : \zh{我的新手机} \\ \emph{Wǒ de xīn de shǒujī}};
\node[box, right=of f3, fill=popgreenL, align=center] (c3){Chinois : \zh{我不可以用} \\ \emph{Wǒ bù kěyǐ yòng}};
\node[box, right=of f4, fill=popgreenL, align=center] (c4){Chinois : \zh{我的老老师} \\ \emph{Wǒ de lǎo lǎoshī}};

% Wrong calques (red crosses)
\node[box, left=of f1, fill=poppinkL, text width=3.5cm] (w1) {Calque fautif : \zh{*我可以手机} \\ \emph{*Wǒ kěyǐ shǒujī}};
\node[box, left=of f2, fill=poppinkL, text width=3.5cm] (w2) {Calque fautif : \zh{*我新手机} \\ \emph{*Wǒ xīn shǒujī}};
\node[box, left=of f3, fill=poppinkL, text width=3.5cm] (w3) {Calque fautif : \zh{*我没可以} \\ \emph{*Wǒ méi kěyǐ}};
\node[box, left=of f4, fill=poppinkL, text width=3.5cm] (w4) {Calque fautif : \zh{*我旧老师} \\ \emph{*Wǒ jiù lǎoshī}};

% Arrows and marks
\foreach \i in {1,2,3,4} {
    \draw[arrow] (f\i) -- (c\i);
    \draw[cross] (w\i.north west) -- (w\i.south east);
    \draw[cross] (w\i.north east) -- (w\i.south west);
    \draw[check, shorten >=2pt, shorten <=2pt] (c\i.north west) -- ++(0.8, -0.8);
    \draw[check, shorten >=2pt, shorten <=2pt] (c\i.south west) -- ++(0.8, 0.8);
}
\node[align=center, font=\footnotesize, text=popred] at ($(w1)!0.5!(w4) - (0, 1.8)$) {Calques francophones \\ à proscrire (×)};
\node[align=center, font=\footnotesize, text=popgreen] at ($(c1)!0.5!(c4) - (0, 1.8)$) {Structures chinoises \\ correctes (✓)};
\node[align=center, font=\footnotesize, text=popblue] at ($(f1)!0.5!(f4) - (0, 1.8)$) {Source française};
\end{tikzpicture}
\legende{Cartographie des interférences français $\rightarrow$ chinois : les calques fautifs (rose, croix rouges) vs structures cibles (vert, coches).}
\end{popfigure}

\courssub{Méthode d'auto-correction en 3 questions}

\begin{methode}[Avant de valider ta phrase, pose-toi ces 3 questions]
\begin{enumerate}
    \item \textbf{Y a-t-il un verbe d'action après 可以 ?} \\
    Si ta phrase s'arrête à \zh{可以} (ou \zh{不可以}) et qu'un objet suit directement, \textbf{ajoute le verbe} (\zh{用}, \zh{买}, \zh{看}, \zh{去}…). \\
    \textit{Exemple :} \zh{*我可以这个} $\rightarrow$ \zh{我可以买这个} (\emph{Wǒ kěyǐ mǎi zhège}).
    \item \textbf{La particule 的 est-elle présente entre 新/旧 et le nom ?} \\
    Si tu as écrit \zh{新手机} ou \zh{旧书包}, \textbf{insère 的} : \zh{新的手机}, \zh{旧的书包}. \\
    \textit{Astuce :} \cle{的} = marqueur « adjectif $\rightarrow$ nom ». Pas de 的, pas de lien clair.
    \item \textbf{La négation est-elle 不 (et non 没) ?} \\
    Devant \cle{可以}, \cle{能}, \cle{会}, \cle{要}, \cle{应该} : \textbf{toujours 不}. \\
    \textit{Règle mnémotechnique :} « Modal = Pas d'action finie = Pas de 没 ».
\end{enumerate}
\end{methode}

\courssub{Exercice résolu : correction de phrases fautives}

\begin{exoresolu}[Correction guidée de 5 phrases pièges]
\textbf{Énoncé :} Corrigez les phrases suivantes en appliquant la méthode en 3 questions. Écrivez la phrase correcte en caractères, pinyin et français.
\medskip
\noindent 1. \zh{*我旧电脑很慢。} \\
2. \zh{*他可以这本书吗？} \\
3. \zh{*我们没可以用微信。} \\
4. \zh{*这是我新朋友。} \\
5. \zh{*你可以去吗？ — 我不可以。} (contexte : réponse complète attendue)

\tcblower
\textbf{Solution détaillée :}
\medskip
\noindent \textbf{1.} \zh{*我旧电脑很慢。} \\
$\rightarrow$ Question 2 : \cle{的} manquant entre \zh{旧} et \zh{电脑}. \\
\textbf{Correct :} \zh{我的旧电脑很慢。} \emph{Wǒ de jiù diànnǎo hěn màn.} — Mon vieil ordinateur est très lent.

\medskip
\noindent \textbf{2.} \zh{*他可以这本书吗？} \\
$\rightarrow$ Question 1 : verbe d'action manquant après \zh{可以}. Objet \zh{这本书} présent. Verbe probable : \zh{看} (lire) ou \zh{买} (acheter) ou \zh{借} (emprunter). \\
\textbf{Correct :} \zh{他可以看这本书吗？} \emph{Tā kěyǐ kàn zhè běn shū ma?} — Peut-il lire ce livre ?

\medskip
\noindent \textbf{3.} \zh{*我们没可以用微信。} \\
$\rightarrow$ Question 3 : négation \zh{没} interdite devant modal. Remplacer par \zh{不}. \\
\textbf{Correct :} \zh{我们不可以用微信。} \emph{Wǒmen bù kěyǐ yòng Wēixìn.} — Nous ne pouvons pas utiliser WeChat. (Interdiction / impossibilité).

\medskip
\noindent \textbf{4.} \zh{*这是我新朋友。} \\
$\rightarrow$ Question 2 : \cle{的} manquant. De plus, « ami » = personne, mais \zh{新} (nouveau/récent) est correct pour « nouvel ami » (récemment connu). \zh{老} serait « vieux ami / ami de longue date ». Ici \zh{新} convient. \\
\textbf{Correct :} \zh{这是我的新朋友。} \emph{Zhè shì wǒ de xīn péngyou.} — C'est mon nouvel ami.

\medskip
\noindent \textbf{5.} \zh{*你可以去吗？ — 我不可以。} \\
$\rightarrow$ Question 1 (sur la réponse) : \zh{不可以} seul est elliptique, acceptable à l'oral très familier, mais incomplet à l'écrit / examen. Ajouter le verbe \zh{去}. \\
\textbf{Correct :} \zh{你可以去吗？ — 我不可以去。} \emph{Nǐ kěyǐ qù ma? — Wǒ bù kěyǐ qù.} — Peux-tu y aller ? — Je ne peux pas y aller.
\end{exoresolu}

\begin{aretenir}[Synthèse des 5 pièges à éviter absolument]
\begin{enumerate}
    \item \textbf{Oublier 的} après \cle{新}/\cle{旧} : \zh{*新手机} $\rightarrow$ \zh{新的手机}.
    \item \textbf{Placer 可以 après le verbe} : \zh{*用可以} $\rightarrow$ \zh{可以用}.
    \item \textbf{Utiliser 可以 sans verbe d'action} : \zh{*我可以这个} $\rightarrow$ \zh{我可以用这个}.
    \item \textbf{Nier 可以 avec 没} : \zh{*没可以} $\rightarrow$ \zh{不可以}.
    \item \textbf{Dire 旧 pour une personne âgée} : \zh{*旧老师} $\rightarrow$ \zh{老老师} (rappel HSK 2).
\end{enumerate}
\end{aretenir}

\courssub{Exercices d'application et de transformation}

\begin{exercice}[Entraînement autonome : transformation et correction]
\textbf{Consigne :} Réécrivez les phrases suivantes correctement (caractères + pinyin + français). Appliquez la méthode en 3 questions.
\begin{enumerate}
    \item \zh{*他们旧房子很大。}
    \item \zh{*我可以那个苹果？}
    \item \zh{*她没可以来学校。}
    \item \zh{*这是你新老师吗？}
    \item \zh{*我们可以游泳吗？ — 我们不可以。} (complétez la réponse)
\end{enumerate}
\end{exercice}

\begin{corrige}[Exercice : transformation et correction]
\begin{enumerate}
    \item \zh{他们的旧房子很大。} \emph{Tāmen de jiù fángzi hěn dà.} — Leur vieille maison est très grande.
    \item \zh{我可以吃那个苹果吗？} \emph{Wǒ kěyǐ chī nàge píngguǒ ma?} — Puis-je manger cette pomme ? (verbe \zh{吃} ajouté).
    \item \zh{她不可以来学校。} \emph{Tā bù kěyǐ lái xuéxiào.} — Elle ne peut pas venir à l'école. (\zh{没} $\rightarrow$ \zh{不}).
    \item \zh{这是你的新老师吗？} \emph{Zhè shì nǐ de xīn lǎoshī ma?} — C'est ton nouveau professeur ? (\cle{的} ajouté).
    \item \zh{我们可以游泳吗？ — 我们不可以游泳。} \emph{Wǒmen kěyǐ yóuyǒng ma? — Wǒmen bù kěyǐ yóuyǒng.} — Pouvons-nous nager ? — Nous ne pouvons pas nager. (verbe répété dans la réponse).
\end{enumerate}
\end{corrige}

\coursec{Leçon 36 — Grammaire : antonyme 新 / 旧 ; verbe modal 可以}

\courssub{Observation guidée des exemples}

\begin{textebox}[Dialogue : Au cybercafé du lycée de Yaoundé]
\zh{— 王伟，这部手机是新的吗？} \emph{— Wáng Wěi, zhè bù shǒujī shì xīn de ma?} \\
\zh{— 是的，是新买的。旧手机坏了，不能上网了。} \emph{— Shì de, shì xīn mǎi de. Jiù shǒujī huài le, bù néng shàngwǎng le.} \\
\zh{— 那我可以用你的新手机查一下成绩吗？} \emph{— Nà wǒ kěyǐ yòng nǐ de xīn shǒujī chá yíxià chéngjī ma?} \\
\zh{— 可以，但不可以玩游戏，课马上要开始了。} \emph{— Kěyǐ, dàn bù kěyǐ wán yóuxì, kè mǎshàng yào kāishǐ le.}
\end{textebox}

\begin{center}
\begin{tabularx}{\linewidth}{|X|X|X|X|}\hline
\rowcolor{popblueL} Caractères & Pinyin & Nature & Traduction \\ \hline
新 & xīn & adj. & nouveau, neuf \\ \hline
旧 & jiù & adj. & vieux, ancien, usagé \\ \hline
手机 & shǒujī & n. & téléphone portable \\ \hline
电脑 & diànnǎo & n. & ordinateur \\ \hline
可以 & kěyǐ & v. modal & pouvoir (permission, possibilité) \\ \hline
上网 & shàngwǎng & v. & aller sur Internet, se connecter \\ \hline
查 & chá & v. & consulter, vérifier, chercher \\ \hline
成绩 & chéngjī & n. & notes, résultats scolaires \\ \hline
玩游戏 & wán yóuxì & v. + n. & jouer à des jeux vidéo \\ \hline
坏 & huài & adj./v. & cassé, en panne / se casser \\ \hline
马上 & mǎshàng & adv. & tout de suite, immédiatement \\ \hline
\end{tabularx}
\end{center}

\begin{experience}[Analyse du dialogue]
En binôme, relevez dans le dialogue :\\
1. Les deux antonymes qui qualifient \zh{手机}.\\
2. La structure utilisée pour demander la permission (\zh{可以…吗？}).\\
3. La structure pour refuser/interdire (\zh{不可以}).\\
4. La place du complément de temps \zh{马上} par rapport au verbe.
\emph{Mise en commun orale : l'enseignant note au tableau les deux structures cibles : \zh{新/旧 + Nom} et \zh{Sujet + 可以 + Verbe}.}
\end{experience}

\courssub{Les antonymes 新  et 旧}

\begin{propriete}[Position des adjectifs qualificatifs monosyllabiques]
En chinois, les adjectifs monosyllabiques comme \zh{新} et \zh{旧} se placent \textbf{devant le nom} qu'ils qualifient, sans \zh{的} dans l'usage courant (sauf s'ils sont complétés par un adverbe de degré comme \zh{很}, \zh{真}, ou s'il y a un possesseur).
\begin{center}
\begin{tabularx}{\linewidth}{|X|X|X|}\hline
\rowcolor{popblueL} Structure & Exemple (caractères + pinyin) & Traduction \\ \hline
Adj + Nom & \zh{新书} \emph{xīn shū} & un livre neuf / des livres neufs \\ \hline
Adj + Nom & \zh{旧电脑} \emph{jiù diànnǎo} & un vieil ordinateur \\ \hline
Possesseur + \zh{的} + Adj + Nom & \zh{我的新手机} \emph{wǒ de xīn shǒujī} & mon nouveau téléphone \\ \hline
Possesseur + \zh{的} + Adj + Nom & \zh{老师的旧书} \emph{lǎoshī de jiù shū} & le vieux livre du prof \\ \hline
\end{tabularx}
\end{center}
\end{propriete}

\begin{attention}[Piège francophone : \zh{的} superflu et ordre possesseur/adjectif]
\begin{itemize}
\item \textbf{Pas de \zh{的} simple :} On dit \zh{新书}, pas *\zh{新的书} (sauf si \zh{新} est mis en valeur : \zh{真新的书} « un livre vraiment neuf »).
\item \textbf{Ordre possesseur – adjectif :} En français « mon vieux livre » $\rightarrow$ possesseur + adjectif + nom. En chinois : \zh{我的旧书} (Possesseur + \zh{的} + Adjectif + Nom). \textbf{Interdit :} *\zh{我旧的书} ou *\zh{旧我的书}.
\end{itemize}
\end{attention}

\begin{exemplebox}[Contexte camerounais]
\zh{我们学校买了新电脑，旧电脑送给了小学。} \emph{Wǒmen xuéxiào mǎi le xīn diànnǎo, jiù diànnǎo sòng gěi le xiǎoxué.} « Notre école a acheté de nouveaux ordinateurs, les vieux ont été donnés à l'école primaire. »\\
\zh{这是我的新书包，那个旧书包坏了。} \emph{Zhè shì wǒ de xīn shūbāo, nàge jiù shūbāo huài le.} « Voici mon nouveau cartable, celui-là (l'ancien) est cassé. »
\end{exemplebox}

\begin{exercice}[Entraînement : 新 ou 旧 ?]
Complétez avec \zh{新} ou \zh{旧} (ajoutez le possesseur + \zh{的} si nécessaire).\\
1. \zh{这部\_\_\_\_手机是我生日礼物，很好用。} (Cadeau d'anniversaire $\rightarrow$ neuf)\\
2. \zh{图书馆的\_\_\_\_报纸不能带走。} (Bibliothèque, ne pas emporter $\rightarrow$ vieux)\\
3. \zh{王老师的\_\_\_\_字典很厚。} (Professeur, dictionnaire épais $\rightarrow$ probablement vieux)\\
4. \zh{我买了\_\_\_\_自行车，骑着去上学。} (Acheté, aller en cours $\rightarrow$ neuf)
\end{exercice}

\begin{corrige}[Exercice d'observation]
1. \zh{新} (\zh{这部新手机…}) — 2. \zh{旧} (\zh{图书馆的旧报纸…}) — 3. \zh{旧} (\zh{王老师的旧字典…}) — 4. \zh{新} (\zh{我买了新自行车…})
\end{corrige}

\courssub{Le verbe modal 可以}

\begin{definition}[可以  — Permission, possibilité, capacité]
\zh{可以} exprime l'autorisation (permission donnée par une règle, une personne, les circonstances) ou la possibilité matérielle. Il se place \textbf{toujours avant le verbe principal}.
\zh{Structure : Sujet + 可以 + Verbe (+ Complément)}
\end{definition}

\begin{propriete}[Les trois formes de base]
\begin{center}
\begin{tabularx}{\linewidth}{|X|X|X|}\hline
\rowcolor{popblueL} Forme & Structure & Exemple \\ \hline
Affirmative & Sujet + \zh{可以} + V + Compl. & \zh{学生可以用电脑。} \emph{Xuésheng kěyǐ yòng diànnǎo.} \\ \hline
Négative & Sujet + \zh{不可以} + V + Compl. & \zh{学生不可以用电脑。} \emph{Xuésheng bù kěyǐ yòng diànnǎo.} \\ \hline
Interrogative & Sujet + \zh{可以} + V + Compl. + \zh{吗} ? & \zh{学生可以用电脑吗？} \emph{Xuésheng kěyǐ yòng diànnǎo ma?} \\ \hline
\end{tabularx}
\end{center}
\end{propriete}

\begin{methode}[Construire les trois formes avec 可以]
\begin{enumerate}
\item \textbf{Identifier le noyau} : Sujet + \zh{可以} + Verbe principal. Ex: \zh{我可以上网}。
\item \textbf{Négation} : Insérer \zh{不} \textbf{immédiatement devant} \zh{可以} $\rightarrow$ \zh{不可以}. Ne touchez à rien d'autre.
\item \textbf{Question} : Ajouter \zh{吗} \textbf{à la toute fin} de la phrase. Ne touchez à rien d'autre.
\item \textbf{Vérifier l'ordre} : Sujet – \zh{(不)可以} – Verbe – Compléments – \zh{吗?}. L'ordre est immuable.
\end{enumerate}
\end{methode}

\begin{exemplebox}[Variantes avec compléments de temps/lieu]
Les compléments de temps/lieu se placent \textbf{avant} \zh{可以} (sujet + Temps/Lieu + \zh{可以} + V) ou en tête de phrase.
\zh{下课后，我们可以去网吧上网。} \emph{Xiàkè hòu, wǒmen kěyǐ qù wǎngbā shàngwǎng.} « Après les cours, on peut aller au cybercafé. »\\
\zh{在图书馆不可以大声说话。} \emph{Zài túshūguǎn bù kěyǐ dàshēng shuōhuà.} « À la bibliothèque, on ne peut pas parler fort. »
\end{exemplebox}

\begin{attention}[可以 vs 能 vs 会 — Nuances HSK 3]
\begin{itemize}
\item \zh{可以} : Permission / Règle / Autorisation (« avoir le droit de »). \emph{Ex: 这里可以抽烟吗？}
\item \zh{能} : Capacité physique / Circonstances matérielles (« être capable de », « pouvoir » au sens physique). \emph{Ex: 我病了，不能去上学。}
\item \zh{会} : Savoir-faire acquis / Compétence (« savoir faire »). \emph{Ex: 他会说中文。}
\end{itemize}
\emph{Dans cette leçon (niveau 1ère A), on se concentre sur \zh{可以} (permission/règle scolaire).}
\end{attention}

\courssub{Comparaison avec le français et pièges des francophones}

\begin{propriete}[Trois différences majeures Français / Chinois]
\begin{center}
\begin{tabularx}{\linewidth}{|X|X|X|}\hline
\rowcolor{popblueL} Phénomène & Français (inversion / déplacement) & Chinois (ajout seul, ordre figé) \\ \hline
\textbf{Question} & Inversion : « Peux-tu… ? » / « Est-ce que… » & Ajout final \zh{吗} : \zh{你可以…吗？} \\ \hline
\textbf{Négation} & \zh{ne… pas} autour du verbe : « ne peux pas » & Infixe \zh{不} devant \zh{可以} : \zh{不可以} \\ \hline
\textbf{Adjectif + Nom} & Possesseur – Adj – Nom : « mon vieux téléphone » & Possesseur – \zh{的} – Adj – Nom : \zh{我的旧手机} \\ \hline
\end{tabularx}
\end{center}
\end{propriete}

\begin{attention}[Les 4 erreurs interdites en 1ère A]
\begin{enumerate}
\item \textbf{Inversion sujet/modal :} *\zh{可以你用电脑吗？} $\rightarrow$ \textbf{\zh{你可以用电脑吗？}}
\item \textbf{\zh{吗} au milieu :} *\zh{你可以吗用电脑？} $\rightarrow$ \textbf{\zh{你可以用电脑吗？}}
\item \textbf{\zh{的} entre Adj monosyllabique et Nom sans possesseur :} *\zh{新的手机} (usage simple) $\rightarrow$ \textbf{\zh{新手机}}. (Avec adverbe de degré : \zh{很新的手机} \checkmark).
\item \textbf{Ordre Possesseur/Adj :} *\zh{我旧的书} / *\zh{旧我的书} $\rightarrow$ \textbf{\zh{我的旧书}}.
\end{enumerate}
\end{attention}

\begin{exoresolu}[Analyse d'erreurs typiques — Corrigez avant de continuer]
Corrigez et expliquez l'erreur (règle violée).\\
1. *\zh{旧的电脑很慢。} (Pas de possesseur, adj monosyllabique) \\
2. *\zh{可以他用手机吗？} (Inversion interdite) \\
3. *\zh{我不可以用新电脑很贵。} (Deux prédicats \zh{用} et \zh{贵} sans liaison) \\
4. *\zh{我的新的手机。} (\zh{的} superflu sans adverbe de degré)
\tcblower
\textbf{Solutions :} \\
1. \zh{旧电脑很慢。} (Règle : Adj mono + Nom sans \zh{的}). \\
2. \zh{他可以用手机吗？} (Règle : Sujet en tête, pas d'inversion). \\
3. Deux phrases : \zh{我不可以用新电脑。} / \zh{新电脑很贵。} (Une phrase = un noyau prédicatif principal). \\
4. \zh{我的新手机}. (Règle : Poss + \zh{的} + Adj mono + Nom).
\end{exoresolu}

\courssub{Exercices d'application et de transformation}

Nous avons vu dans les sections précédentes que les francophones ont tendance à déplacer les mots comme en français (« peux-tu… ? ») ou à coller \zh{的} partout. Les exercices qui suivent vont te permettre de vérifier que tu maîtrises les deux règles d'or de cette leçon : \cle{l'ordre des mots ne change jamais} dans les transformations, et \cle{新 / 旧 se placent devant le nom} (avec ou sans \zh{的} selon le contexte). Nous progresserons par étapes, du plus guidé au plus libre.

\begin{popfigure}
\begin{tikzpicture}[node distance=0.35cm, every node/.style={font=\small}]
\node[draw=popblue, fill=popblueL, rounded corners, align=center, minimum width=2.6cm, minimum height=0.9cm] (e1) {1. Complétion\\ \zh{新 / 旧 / 可以}};
\node[draw=popgreen, fill=popgreenL, rounded corners, align=center, minimum width=2.6cm, minimum height=0.9cm, right=of e1] (e2) {2--3. Transformation\\ \zh{不} / \zh{吗}};
\node[draw=poporange, fill=poporangeL, rounded corners, align=center, minimum width=2.6cm, minimum height=0.9cm, right=of e2] (e3) {4. Correction\\ de phrases fautives};
\node[draw=poppurple, fill=poppurpleL, rounded corners, align=center, minimum width=2.6cm, minimum height=0.9cm, right=of e3] (e4) {5--6. Traduction\\ thème / version};
\node[draw=poppink, fill=poppinkL, rounded corners, align=center, minimum width=2.6cm, minimum height=0.9cm, right=of e4] (e5) {7. Production\\ libre};
\draw[-{Stealth}, thick, popdark] (e1) -- (e2);
\draw[-{Stealth}, thick, popdark] (e2) -- (e3);
\draw[-{Stealth}, thick, popdark] (e3) -- (e4);
\draw[-{Stealth}, thick, popdark] (e4) -- (e5);
\end{tikzpicture}
\legende{Frise de progression : de la complétion guidée à la production libre.}
\end{popfigure}

\courssub{Méthode : comment transformer une phrase avec 可以}

\begin{methode}[Transformer une phrase avec 可以]
\begin{enumerate}
\item \textbf{Identifier} le verbe modal \zh{可以} dans la phrase (il se place toujours juste avant le verbe principal).
\item Pour la \textbf{négation} : insérer \zh{不} immédiatement devant \zh{可以} $\rightarrow$ \zh{不可以}. Tous les autres mots restent à leur place.
\item Pour la \textbf{question} : ajouter \zh{吗} à la toute fin de la phrase et remplacer le point par un point d'interrogation. Tous les autres mots restent à leur place.
\item \textbf{Vérifier} : relire la phrase et contrôler que l'ordre Sujet + \zh{可以} + Verbe + Complément n'a pas bougé.
\end{enumerate}
\end{methode}

\begin{exemplebox}[Une phrase, trois formes]
\zh{我可以用新手机上网。} \emph{Wǒ kěyǐ yòng xīn shǒujī shàngwǎng.} « Je peux aller sur Internet avec le nouveau téléphone. »\\
Négation : \zh{我不可以用新手机上网。} \emph{Wǒ bù kěyǐ yòng xīn shǒujī shàngwǎng.} « Je ne peux pas aller sur Internet avec le nouveau téléphone. »\\
Question : \zh{我可以用新手机上网吗？} \emph{Wǒ kěyǐ yòng xīn shǒujī shàngwǎng ma?} « Est-ce que je peux aller sur Internet avec le nouveau téléphone ? »
\end{exemplebox}

\begin{attention}[L'ordre des mots ne change jamais]
En français, on inverse le sujet et le verbe (« Tu peux » $\rightarrow$ « Peux-tu ? »). En chinois, c'est \textbf{interdit} : on ne dit jamais \zh{可以你…？}. La négation se fait uniquement en ajoutant \zh{不} devant \zh{可以}, et la question uniquement en ajoutant \zh{吗} à la fin. Le reste de la phrase est figé.
\end{attention}

\begin{center}
\begin{tabularx}{\linewidth}{|X|X|X|}\hline
\rowcolor{popblueL} Type de phrase & Exemple (caractères + pinyin) & Traduction \\ \hline
Affirmative & \zh{他可以用电脑。} Tā kěyǐ yòng diànnǎo. & Il peut utiliser l'ordinateur. \\ \hline
Négative & \zh{他不可以用电脑。} Tā bù kěyǐ yòng diànnǎo. & Il ne peut pas utiliser l'ordinateur. \\ \hline
Interrogative & \zh{他可以用电脑吗？} Tā kěyǐ yòng diànnǎo ma? & Peut-il utiliser l'ordinateur ? \\ \hline
Avec \zh{新} & \zh{这是我的新手机。} Zhè shì wǒ de xīn shǒujī. & Voici mon nouveau téléphone. \\ \hline
Avec \zh{旧} & \zh{旧电脑很慢。} Jiù diànnǎo hěn màn. & Le vieil ordinateur est lent. \\ \hline
\end{tabularx}
\end{center}

\courssub{Exercice 1 — Complétion : 新, 旧, 可以 ou 不可以}

\begin{exoresolu}[Compléter avec le mot qui convient]
Énoncé. Complète les phrases suivantes avec \zh{新}, \zh{旧}, \zh{可以} ou \zh{不可以}.\\
1. 我们学校有一台\_\_\_\_电脑。 (Wǒmen xuéxiào yǒu yì tái \_\_\_\_ diànnǎo.)\\
2. 上课的时候，学生\_\_\_\_看手机。\\
3. 这本\_\_\_\_书是我爸爸的，已经很老了。\\
4. 下了课，我们\_\_\_\_去上网。\\
5. 他买了一部\_\_\_\_手机，很好。\\
6. 图书馆的旧报纸不\_\_\_\_带回家。
\tcblower
Solution détaillée.\\
1. \zh{新} : « une \textbf{nouvelle} machine » — le contexte (l'école vient d'acheter) appelle l'antonyme de \zh{旧}. \zh{我们学校有一台新电脑。} \emph{Wǒmen xuéxiào yǒu yì tái xīn diànnǎo.} « Notre école a un nouvel ordinateur. »\\
2. \zh{不可以} : en classe, c'est interdit $\rightarrow$ négation du modal. \zh{上课的时候，学生不可以看手机。} \emph{Shàngkè de shíhou, xuésheng bù kěyǐ kàn shǒujī.} « En cours, les élèves ne peuvent pas regarder leur téléphone. »\\
3. \zh{旧} : « déjà très vieux » annonce l'antonyme \zh{旧}. \zh{这本旧书是我爸爸的。} \emph{Zhè běn jiù shū shì wǒ bàba de.} « Ce vieux livre est à mon père. »\\
4. \zh{可以} : après les cours, c'est permis $\rightarrow$ forme affirmative. \zh{下了课，我们可以去上网。} \emph{Xià le kè, wǒmen kěyǐ qù shàngwǎng.} « Après les cours, nous pouvons aller sur Internet. »\\
5. \zh{新} : on achète en général du neuf ; \zh{很好} confirme. \zh{他买了一部新手机。} \emph{Tā mǎi le yí bù xīn shǒujī.} « Il a acheté un nouveau téléphone. »\\
6. \zh{可以} : la phrase contient déjà \zh{不}, il manque le modal lui-même. \zh{图书馆的旧报纸不可以带回家。} \emph{Túshūguǎn de jiù bàozhǐ bù kěyǐ dài huíjiā.} « On ne peut pas ramener à la maison les vieux journaux de la bibliothèque. »
\end{exoresolu}

\courssub{Exercice 2 — Transformation : affirmative $\rightarrow$ négative}

\begin{exercice}[Passer à la forme négative]
Transforme chaque phrase en ajoutant \zh{不} au bon endroit.\\
1. 他可以用电脑。 (Tā kěyǐ yòng diànnǎo.)\\
2. 我们可以看电视。 (Wǒmen kěyǐ kàn diànshì.)\\
3. 你可以用我的手机。 (Nǐ kěyǐ yòng wǒ de shǒujī.)\\
4. 学生可以在教室里打电话。 (Xuésheng kěyǐ zài jiàoshì lǐ dǎ diànhuà.)
\end{exercice}

\begin{corrige}[Exercice 2]
1. \zh{他不可以用电脑。} \emph{Tā bù kěyǐ yòng diànnǎo.} « Il ne peut pas utiliser l'ordinateur. »\\
2. \zh{我们不可以看电视。} \emph{Wǒmen bù kěyǐ kàn diànshì.} « Nous ne pouvons pas regarder la télévision. » (On garde \zh{看}, on ne change que \zh{可以} $\rightarrow$ \zh{不可以}).\\
3. \zh{你不可以用我的手机。} \emph{Nǐ bù kěyǐ yòng wǒ de shǒujī.} « Tu ne peux pas utiliser mon téléphone. »\\
4. \zh{学生不可以在教室里打电话。} \emph{Xuésheng bù kěyǐ zài jiàoshì lǐ dǎ diànhuà.} « Les élèves ne peuvent pas téléphoner dans la salle de classe. »\\
Dans les quatre cas, seul \zh{不} a été ajouté devant \zh{可以} ; aucun autre mot n'a bougé.
\end{corrige}

\courssub{Exercice 3 — Transformation : déclarative $\rightarrow$ interrogative avec 吗}

\begin{exercice}[Poser la question]
Transforme chaque phrase en question avec \zh{吗}.\\
1. 我可以用你的手机。 (Wǒ kěyǐ yòng nǐ de shǒujī.)\\
2. 我们可以看这个节目。 (Wǒmen kěyǐ kàn zhège jiémù.)\\
3. 你可以上网。 (Nǐ kěyǐ shàngwǎng.)\\
4. 她可以用新电脑。 (Tā kěyǐ yòng xīn diànnǎo.)
\end{exercice}

\begin{corrige}[Exercice 3]
1. \zh{我可以用你的手机吗？} \emph{Wǒ kěyǐ yòng nǐ de shǒujī ma?} « Est-ce que je peux utiliser ton téléphone ? »\\
2. \zh{我们可以看这个节目吗？} \emph{Wǒmen kěyǐ kàn zhège jiémù ma?} « Pouvons-nous regarder cette émission ? »\\
3. \zh{你可以上网吗？} \emph{Nǐ kěyǐ shàngwǎng ma?} « Peux-tu aller sur Internet ? »\\
4. \zh{她可以用新电脑吗？} \emph{Tā kěyǐ yòng xīn diànnǎo ma?} « Peut-elle utiliser le nouvel ordinateur ? »\\
Rappel : pas d'inversion comme en français ; \zh{吗} se colle simplement à la fin.
\end{corrige}

\courssub{Exercice 4 — Correction de phrases fautives}

\begin{exoresolu}[Trouve l'erreur et corrige]
Énoncé. Les phrases suivantes contiennent chacune une erreur typique des francophones. Corrige-les et justifie.\\
1. *\zh{我旧的书很好。}\\
2. *\zh{你可以去吗上网。}\\
3. *\zh{可以他用电脑吗？}\\
4. *\zh{他不可以用旧电脑很慢。}
\tcblower
Solution détaillée.\\
1. Correction : \zh{我的旧书很好。} \emph{Wǒ de jiù shū hěn hǎo.} « Mon vieux livre est très bien. » Justification : en chinois, le possesseur (\zh{我的}) précède l'adjectif (\zh{旧}) ; on ne peut pas dire \zh{我旧的}. Ordre : possesseur + \zh{的} + adjectif + nom.\\
2. Correction : \zh{你可以去上网吗？} \emph{Nǐ kěyǐ qù shàngwǎng ma?} « Peux-tu aller sur Internet ? » Justification : \zh{吗} se place \textbf{à la fin de la phrase entière}, jamais au milieu. Le français « peux-tu » au début induit en erreur.\\
3. Correction : \zh{他可以用电脑吗？} \emph{Tā kěyǐ yòng diànnǎo ma?} Justification : pas d'inversion sujet–verbe en chinois ; le sujet \zh{他} reste en tête et on ajoute seulement \zh{吗} final.\\
4. Correction : deux phrases possibles — \zh{他不可以用旧电脑。} \emph{Tā bù kěyǐ yòng jiù diànnǎo.} « Il ne peut pas utiliser le vieil ordinateur » ou \zh{旧电脑很慢。} \emph{Jiù diànnǎo hěn màn.} « Le vieil ordinateur est lent. » Justification : on ne peut pas empiler deux verbes prédicatifs (\zh{用} et \zh{很}) dans une même phrase sans structure adaptée (conjonction, subordination).
\end{exoresolu}

\begin{attention}[Les trois erreurs à bannir]
\begin{itemize}
\item Placer \zh{吗} au milieu de la phrase : *\zh{你可以吗上网…} est impossible.
\item Inverser sujet et verbe modal comme en français : *\zh{可以你…？} est impossible.
\item Mettre \zh{的} entre l'adjectif antonyme et le nom quand il n'y a pas de possesseur : on dit \zh{旧书} « vieux livre », pas *\zh{旧的书} dans un usage simple (avec un possesseur : \zh{我的旧书}).
\end{itemize}
\end{attention}

\courssub{Exercices 5 et 6 — Traduction : thème et version}

\begin{exercice}[Thème : français $\rightarrow$ chinois]
Traduis en chinois (caractères + pinyin) :\\
1. « Je peux utiliser le nouvel ordinateur, mais pas le vieux. »\\
2. « Les élèves ne peuvent pas téléphoner en classe. »\\
3. « Peux-tu regarder la télévision ce soir ? »
\end{exercice}

\begin{corrige}[Exercice 5]
1. \zh{我可以用新电脑，不可以用旧电脑。} \emph{Wǒ kěyǐ yòng xīn diànnǎo, bù kěyǐ yòng jiù diànnǎo.} — « mais pas le vieux » se rend par la répétition du modal à la négative : \zh{不可以用旧电脑}. Le chinois répète le verbe \zh{用}, contrairement au français qui sous-entend.\\
2. \zh{学生不可以在课堂上打电话。} \emph{Xuésheng bù kěyǐ zài kètáng shang dǎ diànhuà.} — \zh{在课堂上} « en classe / pendant le cours ».\\
3. \zh{你今天晚上可以看电视吗？} \emph{Nǐ jīntiān wǎnshang kěyǐ kàn diànshì ma?} — le complément de temps \zh{今天晚上} se place avant le verbe modal \zh{可以}.
\end{corrige}

\begin{exercice}[Version : chinois $\rightarrow$ français]
Traduis en français :\\
1. 学生不可以在课堂上看电视。 (Xuésheng bù kěyǐ zài kètáng shang kàn diànshì.)\\
2. 这部旧手机还可以用。 (Zhè bù jiù shǒujī hái kěyǐ yòng.)\\
3. 妈妈说不可以用新电脑玩游戏。 (Māma shuō bù kěyǐ yòng xīn diànnǎo wán yóuxì.)
\end{exercice}

\begin{corrige}[Exercice 6]
1. « Les élèves ne peuvent pas regarder la télévision en classe. » — \zh{在课堂上} = « en classe / pendant le cours ».\\
2. « Ce vieux téléphone peut encore servir / fonctionne encore. » — \zh{还} = « encore » ; il se place avant \zh{可以}. \zh{用} ici signifie « servir / fonctionner ».\\
3. « Maman dit qu'on ne peut pas jouer aux jeux vidéo avec le nouvel ordinateur. » — \zh{说} introduit une parole rapportée, sans mot-outil équivalent à « que ». \zh{玩游戏} = jouer à des jeux.
\end{corrige}

\courssub{Exercice 7 — Production guidée : « mon téléphone »}

\begin{experience}[À toi d'écrire]
Rédige \textbf{quatre phrases} sur le thème « mon téléphone » (\zh{我的手机}), en utilisant chacune des quatre formes de la leçon : une phrase avec \zh{新}, une avec \zh{旧}, une avec \zh{可以}, une avec \zh{不可以}.\\
Modèles pour t'inspirer :\\
\zh{我的新手机很漂亮。} \emph{Wǒ de xīn shǒujī hěn piàoliang.} « Mon nouveau téléphone est très beau. »\\
\zh{旧手机还可以打电话。} \emph{Jiù shǒujī hái kěyǐ dǎ diànhuà.} « Le vieux téléphone peut encore téléphoner. »\\
\zh{我可以用手机看新闻。} \emph{Wǒ kěyǐ yòng shǒujī kàn xīnwén.} « Je peux lire les informations sur mon téléphone. »\\
\zh{上课的时候，我不可以用手机。} \emph{Shàngkè de shíhou, wǒ bù kěyǐ yòng shǒujī.} « Pendant les cours, je ne peux pas utiliser mon téléphone. »\\
Consigne de relecture : vérifie que \zh{新}/\zh{旧} sont bien devant le nom, que \zh{可以} précède le verbe, et que la négation est \zh{不可以}.
\end{experience}

\begin{aretenir}[Ce que tu dois savoir faire après cette leçon]
\begin{itemize}
\item Placer \zh{新} / \zh{旧} devant le nom : \zh{新手机}, \zh{旧电脑} ; avec un possesseur : \zh{我的旧书}.
\item Construire les trois formes avec \zh{可以} : affirmative (\zh{可以}), négative (\zh{不可以}), interrogative (\zh{可以…吗}).
\item Transformer une phrase \textbf{sans jamais changer l'ordre des mots} : on ajoute seulement \zh{不} ou \zh{吗}.
\item Traduire dans les deux sens des phrases du domaine des médias et des télécommunications.
\item Distinguer \zh{可以} (permission) de \zh{能} (capacité) et \zh{会} (savoir-faire).
\end{itemize}
\end{aretenir}

\newpage

\coursec{Activité d'intégration}

\coursec{Leçon 36 — Grammaire : antonymes 新 / 旧 ; verbe modal 可以}

\courssub{Observation guidée}

\begin{textebox}[Modèle d'affiche : Règlement du cybercafé « 新新网吧 »]
新新网吧规定 \\
Xīnxīn Wǎngbā guīdìng \\
« Règlement du cybercafé Nouveau-Nouveau » \\[0.3em]
一、你可以用新电脑。 \\
Yī, nǐ kěyǐ yòng xīn diànnǎo. \\
« 1. Tu peux utiliser les nouveaux ordinateurs. » \\[0.3em]
二、旧电脑也可以，但是很慢。 \\
Èr, jiù diànnǎo yě kěyǐ, dànshì hěn màn. \\
« 2. Les vieux ordinateurs aussi sont permis, mais ils sont très lents. » \\[0.3em]
三、你不可以在电脑旁边吃饭。 \\
Sān, nǐ bù kěyǐ zài diànnǎo pángbiān chīfàn. \\
« 3. Tu ne peux pas manger à côté des ordinateurs. »
\end{textebox}

\begin{experience}[Activité d'observation (5 min)]
En binôme, observez l'affiche et répondez :
\begin{enumerate}
\item Soulignez toutes les occurrences de \zh{可以} et \zh{不可以}. Entourez \zh{新} et \zh{旧}.
\item Quelle est la position de \zh{可以} / \zh{不可以} par rapport au sujet \zh{你} et au verbe (\zh{用}, \zh{吃}) ?
\item \zh{新} et \zh{旧} sont-ils placés avant ou après le nom \zh{电脑} ? Y a-t-il \zh{的} entre l'adjectif et le nom ?
\item Traduisez les trois règles en français.
\end{enumerate}
\end{experience}

\courssub{Les antonymes 新 (xīn) et 旧 (jiù)}

\begin{propriete}[Emploi des adjectifs attributs 新 et 旧]
En chinois, les adjectifs monosyllabiques comme \zh{新} (nouveau) et \zh{旧} (vieux, ancien, usagé) fonctionnent comme des \textbf{adjectifs attributs} : ils se placent \textbf{directement devant le nom} qu'ils qualifient, sans verbe \zh{是} (être) et généralement \textbf{sans \zh{的}} s'ils sont monosyllabiques et que le nom est courant.
\begin{center}
\begin{tabularx}{\linewidth}{|X|X|X|}
\hline
\rowcolor{popblueL} Structure & Exemple (Caractères + Pinyin) & Traduction \\
\hline
Adj (monosyll.) + Nom & \zh{新电脑} (xīn diànnǎo) & le nouvel ordinateur / les nouveaux ordinateurs \\
\hline
Adj (monosyll.) + Nom & \zh{旧书} (jiù shū) & le vieux livre / les vieux livres \\
\hline
Adj + \zh{的} + Nom (insistance / adj. bisyll.) & \zh{新的手机} (xīn de shǒujī) & le téléphone qui est nouveau (insistance sur la nouveauté) \\
\hline
\end{tabularx}
\end{center}
\end{propriete}

\begin{exemplebox}[Exemples contextuels]
\begin{itemize}
\item \zh{这是我的新书包。} (Zhè shì wǒ de xīn shūbāo.) — « Voici mon nouveau cartable. »
\item \zh{他不喜欢旧衣服。} (Tā bù xǐhuan jiù yīfu.) — « Il n'aime pas les vieux vêtements. »
\item \zh{新旧社会} (xīnjiù shèhuì) — « l'ancienne et la nouvelle société » (composé figé).
\item \zh{旧朋友} (jiù péngyou) — « vieux ami / ami de longue date » (ne signifie pas « ami âgé »).
\end{itemize}
\end{exemplebox}

\begin{attention}[Pièges francophones : 新 / 旧]
\begin{itemize}
\item \textbf{Interdiction du \zh{是} attributif} : On ne dit jamais \zh{*这是是新电脑} pour « C'est un nouvel ordinateur ». \zh{是} n'est utilisé que pour identifier le sujet (\zh{这是新电脑} = « Ceci EST un nouvel ordi »).
\item \textbf{Oubli du \zh{的} avec adjectif bisyllabique ou phrase adjectivale} : \zh{全新的电脑} (quánxīn de diànnǎo, « ordinateur tout neuf ») nécessite \zh{的}. \zh{新电脑} (monosyllabique) ne le nécessite pas.
\item \textbf{Sens de \zh{旧}} : \zh{旧} = « usagé, ancien, d'autrefois ». Pour « ancien » (ex-ancien président), on utilise \zh{前} (qián) ou \zh{原} (yuán) : \zh{前总统} (ancien président).
\end{itemize}
\end{attention}

\courssub{Le verbe modal 可以 (kěyǐ) : permission et possibilité}

\begin{propriete}[Structure centrale : Sujet + 可以 + Verbe (+ Complément)]
\zh{可以} (kěyǐ) exprime la \textbf{permission} (« avoir le droit de », « pouvoir » au sens déontique) ou la \textbf{possibilité favorable} (« c'est possible / ça marche »).
\medskip
\textbf{1. Affirmatif (Permission / Autorisation)}
\[
\text{Sujet} + \textbf{可以} + \text{Verbe} + (\text{Complément})
\]
Ex. \zh{你可以上网。} (Nǐ kěyǐ shàngwǎng.) « Tu peux / Tu as le droit d'aller sur Internet. »

\medskip
\textbf{2. Négatif (Interdiction) : \zh{不可以} (bù kěyǐ)}
\[
\text{Sujet} + \textbf{不可以} + \text{Verbe}
\]
Ex. \zh{你不可以打电话。} (Nǐ bù kěyǐ dǎ diànhuà.) « Tu ne peux pas / Tu n'as pas le droit de téléphoner. »
\emph{Attention :} La négation est \zh{不} (bù), jamais \zh{没} (méi). \zh{*没可以} n'existe pas.

\medskip
\textbf{3. Interrogatif (Demande de permission) : \zh{可以……吗？}}
\[
\text{Sujet} + \textbf{可以} + \text{Verbe} + \textbf{吗}？
\]
Ex. \zh{我可以用这台电脑吗？} (Wǒ kěyǐ yòng zhè tái diànnǎo ma?) « Puis-je utiliser cet ordinateur ? »
Réponse courte : \zh{可以。} (Oui, tu peux.) / \zh{不可以。} (Non, tu ne peux pas.)
\end{propriete}

\begin{exemplebox}[Variantes et nuances]
\begin{itemize}
\item \zh{也可以} (yě kěyǐ) : « peut aussi / est aussi permis ». Ex. \zh{旧电脑也可以。} « Les vieux ordis vont aussi bien. »
\item \zh{可以不……} : « Peut-on ne pas… ? » (demande de dispense). Ex. \zh{我可以不去吗？} « Est-ce que je peux ne pas y aller ? »
\item \zh{可以} seul en réponse : \zh{—— 我可以进来吗？ —— 可以。} (« — Puis-je entrer ? — Oui / Entrez. »)
\end{itemize}
\end{exemplebox}

\begin{propriete}[Règle de changement de ton (Sandhi) pour 可以]
\zh{可以} s'écrit avec deux 3\textsuperscript{e} tons (kě + yǐ). \textbf{À l'oral}, deux 3\textsuperscript{e} tons consécutifs déclenchent la règle du sandhi : le premier 3\textsuperscript{e} ton devient un \textbf{2\textsuperscript{e} ton} (montant).
\begin{center}
\begin{tabularx}{\linewidth}{|X|X|}
\hline
\rowcolor{poporangeL} Écriture (Citation) & Prononciation réelle (Contexte) \\
\hline
kěyǐ (3-3) & \textbf{kéyǐ} (2-3) \\
\hline
\end{tabularx}
\end{center}
\emph{Exercice :} Répétez : \zh{可以} (kéyǐ), \zh{不可以} (bú kěyǐ — \zh{不} devient 2\textsuperscript{e} ton devant 3\textsuperscript{e} ton), \zh{可以吗} (kéyǐ ma).
\end{propriete}

\begin{popfigure}
\begin{tikzpicture}[
  node distance=1.2cm and 1.5cm,
  box/.style={draw=popblue, thick, rounded corners, fill=popblueL, text width=3.2cm, align=center, font=\small},
  arrow/.style={-Stealth, thick, popdark}
]
\node[box, align=center] (suj){Sujet\\ (你, 我, 他...)};
\node[box, right=of suj, align=center] (modal){\textbf{可以 / 不可以}\\ (kéyǐ / bú kěyǐ)};
\node[box, right=of modal, align=center] (verb){Verbe d'action\\ (用, 吃, 去, 看...)};
\node[box, right=of verb, align=center] (comp){Complément\\ (objet, lieu, manière)};
\draw[arrow] (suj) -- (modal);
\draw[arrow] (modal) -- (verb);
\draw[arrow] (verb) -- (comp);
\node[below=0.3cm of modal, text=poporange, font=\bfseries\small] {Position fixe : après le sujet, avant le verbe};
\end{tikzpicture}
\legende{Schéma de la structure modale avec 可以}
\end{popfigure}

\courssub{Comparaison avec le français et pièges des francophones}

\begin{attention}[Erreurs fréquentes : 可以 vs français « pouvoir »]
\begin{itemize}
\item \textbf{Ordre des mots} : FR « Peux-tu venir ? » (inversion) $\neq$ CH \zh{*可以你来？}. CH : \zh{你可以来吗？} (Sujet - Modal - Verbe - Ma). \textbf{Jamais d'inversion sujet-verbe en chinois.}
\item \textbf{Négation} : FR « Tu ne peux pas » (ne...pas autour du verbe). CH : \zh{你不可以来} (Négation \zh{不} \textbf{devant} le modal \zh{可以}, pas devant le verbe principal). \zh{*你可以不来} = « Tu peux ne pas venir » (permission de ne pas venir), ce n'est pas une interdiction.
\item \textbf{可以 vs 能 vs 会} :
\begin{center}
\begin{tabularx}{\linewidth}{|X|X|X|}
\hline
\rowcolor{poppurpleL} Verbe modal & Sens principal & Exemple \\
\hline
\zh{可以} (kěyǐ) & Permission / Règle / Autorisation & \zh{这里可以抽烟。} « On a le droit de fumer ici. » \\
\hline
\zh{能} (néng) & Capacité physique / Circonstance / Possibilité & \zh{我会游泳，但今天生病了，不能去。} « Je sais nager, mais je suis malade, je ne peux pas y aller. » \\
\hline
\zh{会} (huì) & Savoir-faire acquis / Compétence & \zh{他会说中文。} « Il sait parler chinois. » \\
\hline
\end{tabularx}
\end{center}
\item \textbf{Traduction de « Tu peux »} : Selon le contexte, choisir \zh{可以} (permission), \zh{能} (capacité/possibilité), \zh{会} (savoir). Ne pas traduire systématiquement par \zh{可以}.
\end{itemize}
\end{attention}

\courssub{Exercices d'application et de transformation : Situation « Cybercafé à Douala »}

\begin{textebox}[Contexte]
Tu es membre d'une équipe de trois élèves de Première A. Votre oncle vient d'ouvrir un cybercafé à Douala, quartier Akwa, près du lycée. Le cybercafé s'appelle « \zh{新新网吧} » (Xīnxīn Wǎngbā, « le cybercafé Nouveau-Nouveau »). Il a acheté dix \cle{nouveaux} ordinateurs, mais il a gardé quelques \cle{vieux} ordinateurs pour les jeux. Beaucoup de clients sont des élèves : certains mangent devant les écrans, d'autres téléphonent très fort, d'autres encore veulent regarder des matchs de football pendant que leurs camarades font leurs recherches scolaires. Votre oncle vous demande de rédiger, en chinois, le \cle{règlement} qui sera affiché à l'entrée, puis de le présenter aux clients (la classe).
\end{textebox}

\begin{center}
\begin{tabularx}{\linewidth}{|X|X|X|X|}
\hline
\rowcolor{popblueL} Caractères & Pinyin & Nature & Traduction \\
\hline
网吧 & wǎngbā & n. & cybercafé \\
\hline
规定 & guīdìng & n. & règlement, règle \\
\hline
电脑 & diànnǎo & n. & ordinateur \\
\hline
台 & tái & m. & classificateur pour machines/ordinateurs \\
\hline
用 & yòng & v. & utiliser \\
\hline
上网 & shàngwǎng & v. & aller sur Internet, surfer \\
\hline
打电话 & dǎ diànhuà & v. & téléphoner \\
\hline
看电视 & kàn diànshì & v. & regarder la télévision \\
\hline
听音乐 & tīng yīnyuè & v. & écouter de la musique \\
\hline
玩游戏 & wán yóuxì & v. & jouer à des jeux (vidéo) \\
\hline
吃饭 & chīfàn & v. & manger (un repas) \\
\hline
喝水 & hē shuǐ & v. & boire de l'eau \\
\hline
说话 & shuōhuà & v. & parler, discuter \\
\hline
大声 & dàshēng & adv. & fort (voix), à haute voix \\
\hline
旁边 & pángbiān & n. (loc.) & à côté (de) \\
\hline
慢 & màn & adj. & lent, lente \\
\hline
欢迎 & huānyíng & v. & accueillir, souhaiter la bienvenue \\
\hline
新 & xīn & adj. & nouveau, neuf \\
\hline
旧 & jiù & adj. & vieux, ancien, usagé \\
\hline
\end{tabularx}
\end{center}

\courssub{Tâche 1 — Compréhension et analyse (10 min)}
\begin{exercice}[Analyse du modèle]
\begin{enumerate}
\item Relevez dans le modèle (Section 1) les 3 structures : \zh{可以 + V}, \zh{也可以}, \zh{不可以 + V}.
\item Pourquoi \zh{新电脑} et \zh{旧电脑} s'écrivent-ils sans \zh{的} ?
\item Dans la règle 3, identifiez le complément de lieu (\zh{在...旁边}) et sa place par rapport au verbe.
\end{enumerate}
\end{exercice}

\courssub{Tâche 2 — Production écrite : Rédigez le règlement complet (25 min)}
\begin{exercice}[Consignes de rédaction]
Rédigez l'affiche complète sur une feuille A4 (ou dans le cahier) en respectant \textbf{strictement} les contraintes suivantes :
\begin{enumerate}
\item \textbf{Titre} : \zh{新新网吧规定} (Xīnxīn Wǎngbā guīdìng).
\item \textbf{Cinq règles numérotées} : \zh{一、} \zh{二、} \zh{三、} \zh{四、} \zh{五、}.
\item \textbf{Contraintes grammaticales obligatoires} :
      \begin{itemize}
      \item Au moins \textbf{deux règles avec \zh{可以} ou \zh{也可以}} (permissions).
      \item Au moins \textbf{deux règles avec \zh{不可以}} (interdictions).
      \item Les mots \zh{新} et \zh{旧} doivent apparaître \textbf{au moins deux fois chacun} dans les cinq règles.
      \item Structure de chaque règle : \textbf{Sujet + 可以 / 不可以 + Verbe (+ Complément)}. Sujet recommandé : \zh{你} (nǐ).
      \item Ponctuation chinoise : virgule \zh{、} pour énumérer dans une règle, point final \zh{。}.
      \end{itemize}
\item \textbf{Phrase de bienvenue finale} : \zh{欢迎光临！} (Huānyíng guānglín !) ou \zh{欢迎您！}.
\end{enumerate}
\end{exercice}

\courssub{Tâche 3 — Présentation orale et interaction (15 min par groupe)}
\begin{experience}[Déroulement oral]
\begin{enumerate}
\item \textbf{Lecture} : Un élève lit le titre et la bienvenue ; les deux autres lisent les 5 règles (2-2-1 ou 3-2). \textbf{Focus prononciation} : \zh{可以} $\rightarrow$ \emph{kéyǐ} (2-3) ; \zh{不可以} $\rightarrow$ \emph{bú kěyǐ} (2-3) ; \zh{新} (xīn, 1), \zh{旧} (jiù, 4).
\item \textbf{Jeu de rôle « Clients / Gérants »} : Les autres groupes posent \textbf{au moins 2 questions} avec \zh{可以……吗？}.
      \begin{itemize}
      \item Ex. \zh{我可以用旧电脑玩游戏吗？} (Puis-je jouer sur les vieux ordis ?)
      \item Ex. \zh{我可以在新电脑旁边喝水吗？} (Puis-je boire de l'eau à côté des nouveaux ordis ?)
      \end{itemize}
\item \textbf{Réponse attendue} : \zh{可以。} / \zh{不可以。} + justification courte (ex. \zh{不可以。规定第三条：不可以吃饭。}).
\end{enumerate}
\end{experience}

\courssub{Proposition de production modèle et corrigé commenté}

\begin{exoresolu}[Affiche modèle respectant toutes les contraintes]
Énoncé : Rédiger le règlement complet (5 règles, titre, bienvenue) avec $\ge$ 2 \zh{可以}, $\ge$ 2 \zh{不可以}, $\ge$ 2 \zh{新}, $\ge$ 2 \zh{旧}.
\tcblower
\textbf{Production écrite :}\\[0.3em]
新新网吧规定 \\
Xīnxīn Wǎngbā guīdìng \\[0.3em]
一、你可以用新电脑学习。 \\
Yī, nǐ kěyǐ yòng xīn diànnǎo xuéxí. \\
« 1. Tu peux utiliser les nouveaux ordinateurs pour étudier. » \\[0.3em]
二、旧电脑也可以，但是很慢。 \\
Èr, jiù diànnǎo yě kěyǐ, dànshì hěn màn. \\
« 2. Les vieux ordinateurs sont aussi autorisés, mais ils sont très lents. » \\[0.3em]
三、你不可以在新电脑旁边吃饭。 \\
Sān, nǐ bù kěyǐ zài xīn diànnǎo pángbiān chīfàn. \\
« 3. Tu ne peux pas manger à côté des nouveaux ordinateurs. » \\[0.3em]
四、你不可以在旧电脑上看电视。 \\
Sì, nǐ bù kěyǐ zài jiù diànnǎo shàng kàn diànshì. \\
« 4. Tu ne peux pas regarder la télévision sur les vieux ordinateurs. » \\[0.3em]
五、你可以听音乐，也可以喝水。 \\
Wǔ, nǐ kěyǐ tīng yīnyuè, yě kěyǐ hē shuǐ. \\
« 5. Tu peux écouter de la musique, et tu peux aussi boire de l'eau. » \\[0.3em]
欢迎光临！ \\
Huānyíng guānglín! \\
« Bienvenue ! » \\[0.5em]
\textbf{Vérification des contraintes :}
\begin{center}
\begin{tabularx}{\linewidth}{|X|X|X|X|X|X|}
\hline
\rowcolor{popgreenL} Règle & Structure & \zh{可以} / \zh{不可以} & \zh{新} & \zh{旧} & Validation \\
\hline
一 & 可以 + 用 & \zh{可以} (Perm) & 1 & 0 & OK \\
\hline
二 & 也可以 (elliptique) & \zh{也可以} (Perm) & 0 & 1 & OK \\
\hline
三 & 不可以 + 吃 & \zh{不可以} (Interd) & 1 & 0 & OK \\
\hline
四 & 不可以 + 看 & \zh{不可以} (Interd) & 0 & 1 & OK \\
\hline
五 & 可以 + 听 / 也可以 + 喝 & \zh{可以}, \zh{也可以} (Perm) & 0 & 0 & OK \\
\hline
\textbf{Total} & & \textbf{3 Perm / 2 Interd} & \textbf{2} & \textbf{2} & \textbf{TOUTES CONTRAINTES RESPECTÉES} \\
\hline
\end{tabularx}
\end{center}

\textbf{Commentaire linguistique détaillé :}
\begin{itemize}
\item \textbf{Antonymes} : \zh{新} et \zh{旧} sont utilisés 2 fois chacun comme adjectifs attributs monosyllabiques devant \zh{电脑} (R1, R3 pour \zh{新} ; R2, R4 pour \zh{旧}). Pas de \zh{的} nécessaire. Naturel et correct.
\item \textbf{Modal 可以} : Position invariante après le sujet \zh{你} (R1, R3, R4, R5) ou après le sujet \zh{旧电脑} (R2). Négation correcte \zh{不可以} (R3, R4). Coordination \zh{也可以} (R2, R5) maîtrisée.
\item \textbf{Compléments de lieu} : R3 \zh{在新电脑旁边} (à côté des nouveaux ordis) ; R4 \zh{在旧电脑上} (sur les vieux ordis). Structure \zh{在 + Lieu + Verbe} respectée.
\item \textbf{Vocabulaire module} : \zh{学习} (étudier), \zh{看电视}, \zh{听音乐}, \zh{喝水}, \zh{上网} (implicite dans \zh{用电脑}) réinvestis.
\end{itemize}
\end{exoresolu}

\courssub{Exercices de consolidation individuels (Devoirs / Travail autonome)}

\begin{exercice}[Transformation : Affirmatif $\rightarrow$ Négatif / Interrogatif]
Transformez les phrases suivantes selon le modèle. Écrivez en caractères, pinyin et français.
\begin{enumerate}
\item \zh{你可以用这台新电脑。} $\rightarrow$ (Négatif) $\rightarrow$ (Interrogatif \zh{吗})
\item \zh{他们可以在旧电脑上玩游戏。} $\rightarrow$ (Négatif) $\rightarrow$ (Interrogatif \zh{可以不……})
\item \zh{我们可以喝水。} $\rightarrow$ (Négatif) $\rightarrow$ (Interrogatif \zh{吗})
\end{enumerate}
\end{exercice}

\begin{exercice}[Traduction thème : Français $\rightarrow$ Chinois]
Traduisez en chinois (caractères + pinyin) en utilisant \zh{新}, \zh{旧}, \zh{可以}, \zh{不可以}.
\begin{enumerate}
\item Ce vieil ordinateur est très lent, mais on peut l'utiliser. (\textit{Indice : sujet = 这台旧电脑})
\item Les nouveaux ordinateurs sont interdits aux jeux. (\textit{Indice : 玩游戏 wán yóuxì ; structure : 新电脑不可以……})
\item Puis-je utiliser un nouvel ordinateur pour aller sur Internet ?
\item Vous ne pouvez pas parler fort à côté de l'ordinateur.
\item Les vieux téléphones portables ne peuvent pas prendre de photos. (\textit{Indice : 手机 shǒujī, 照相 zhàoxiàng})
\end{enumerate}
\end{exercice}

\begin{corrige}[Exercice Transformation]
\begin{enumerate}
\item Nég. : \zh{你不可以用这台新电脑。} (Nǐ bù kěyǐ yòng zhè tái xīn diànnǎo.) « Tu ne peux pas utiliser ce nouvel ordi. » \\
      Int. : \zh{你可以用这台新电脑吗？} (Nǐ kěyǐ yòng zhè tái xīn diànnǎo ma?) « Peux-tu utiliser ce nouvel ordi ? »
\item Nég. : \zh{他们不可以在旧电脑上玩游戏。} (Tāmen bù kěyǐ zài jiù diànnǎo shàng wán yóuxì.) \\
      Int. (dispense) : \zh{他们可以不在旧电脑上玩游戏吗？} (Tāmen kěyǐ bù zài jiù diànnǎo shàng wán yóuxì ma?) « Peuvent-ils ne pas jouer sur les vieux ordis ? »
\item Nég. : \zh{我们不可以喝水。} (Wǒmen bù kěyǐ hē shuǐ.) \\
      Int. : \zh{我们可以喝水吗？} (Wǒmen kěyǐ hē shuǐ ma?)
\end{enumerate}
\end{corrige}

\begin{corrige}[Exercice Traduction]
\begin{enumerate}
\item \zh{这台旧电脑很慢，但是可以用。} (Zhè tái jiù diànnǎo hěn màn, dànshì kěyǐ yòng.)
\item \zh{新电脑不可以玩游戏。} (Xīn diànnǎo bù kěyǐ wán yóuxì.)
\item \zh{我可以用新电脑上网吗？} (Wǒ kěyǐ yòng xīn diànnǎo shàngwǎng ma?)
\item \zh{你不可以在电脑旁边大声说话。} (Nǐ bù kěyǐ zài diànnǎo pángbiān dàshēng shuōhuà.)
\item \zh{旧手机不可以照相。} (Jiù shǒujī bù kěyǐ zhàoxiàng.)
\end{enumerate}
\end{corrige}

\courssub{Bilan et auto-évaluation}

\begin{aretenir}[Synthèse grammaticale de la leçon 36]
\begin{itemize}
\item \textbf{Antonymes} : \zh{新} (xīn, nouveau) $\leftrightarrow$ \zh{旧} (jiù, vieux/ancien). Adjectifs attributs monosyllabiques $\rightarrow$ placés \textbf{devant le nom sans \zh{的}} (\zh{新电脑}, \zh{旧书}).
\item \textbf{Verbe modal 可以} : Permission / Autorisation.
      \begin{itemize}
      \item Structure : \textbf{S + 可以 + V (+ Compl)}.
      \item Négation : \textbf{S + 不可以 + V} (Interdiction). Pas de \zh{没可以}.
      \item Question : \textbf{S + 可以 + V + 吗？}
      \item Prononciation : \zh{可以} $\rightarrow$ \emph{kéyǐ} (2-3) ; \zh{不可以} $\rightarrow$ \emph{bú kěyǐ} (2-3).
      \end{itemize}
\item \textbf{Distinction cruciale} : \zh{可以} (permission/règle) $\neq$ \zh{能} (capacité/circonstance) $\neq$ \zh{会} (savoir-faire).
\end{itemize}
\end{aretenir}

\begin{experience}[Grille d'auto-évaluation finale (à coller dans le cahier)]
Cochez \ding{51} si acquis, \ding{55} si à réviser.
\begin{itemize}
\item \ding{51} \ding{55} Je sais placer \zh{新} et \zh{旧} correctement devant un nom (sans \zh{的} pour monosyllabiques).
\item \ding{51} \ding{55} Je construis correctement : Sujet + \zh{可以} + Verbe (permission).
\item \ding{51} \ding{55} Je construis correctement : Sujet + \zh{不可以} + Verbe (interdiction).
\item \ding{51} \ding{55} Je forme la question de permission : Sujet + \zh{可以} + Verbe + \zh{吗} ?
\item \ding{51} \ding{55} Je prononce \zh{可以} \emph{kéyǐ} (ton 2 + ton 3) et \zh{不可以} \emph{bú kěyǐ}.
\item \ding{51} \ding{55} Je distingue \zh{可以} (permission) de \zh{能} (capacité) et \zh{会} (savoir).
\item \ding{51} \ding{55} J'ai rédigé un règlement de 5 règles respectant toutes les contraintes (vocabulaire, grammaire, quantificateurs).
\item \ding{51} \ding{55} J'ai participé à l'oral (lecture + questions-réponses \zh{可以……吗？}).
\end{itemize}
\end{experience}

\begin{savaistu}[Le cybercafé en Chine : 网吧 (wǎngbā)]
Dans les années 2000-2010, les \zh{网吧} étaient le lieu principal d'accès à Internet pour la jeunesse chinoise (jeux en ligne, messagerie QQ). Aujourd'hui, avec la généralisation du smartphone (\zh{智能手机} zhìnéng shǒujī) et de la 4G/5G, les cybercafés classiques ont décliné. Ils se transforment en \zh{电竞馆} (diànjìnguǎn, « arènes e-sport ») haut de gamme avec des chaises ergonomiques, des écrans 144Hz/240Hz, et service de restauration. En zone rurale ou pour les étudiants sans PC personnel, le \zh{网吧} reste un espace social abordable (souvent 2-5 CNY/heure, soit 150-400 FCFA). La réglementation chinoise impose l'interdiction aux mineurs de moins de 18 ans (\zh{未成年人禁止入内} wèichéngniánrén jìnzhǐ rùnèi) et la fermeture nocturne (souvent 00h-8h). Au Cameroun, les cybercafés (quartiers Akwa, Messa, Université) restent vitaux pour l'impression, la formation bureautique et les démarches administratives en ligne.
\end{savaistu}

\newpage

\coursec{Méthodes et exercices résolus}

\coursec{Leçon 36 — Grammaire : antonyme 新 / 旧 ; verbe modal 可以}

\courssub{Vocabulaire clé de la leçon}
\begin{definition}[Mots et expressions à maîtriser (HSK 2-3)]
\begin{center}
\begin{tabularx}{\linewidth}{|X|X|X|X|}
\hline
\rowcolor{popblueL} \textbf{Caractères} & \textbf{Pinyin} & \textbf{Nature} & \textbf{Traduction} \\
\hline
新 & xīn & adj. & nouveau, neuf \\
\hline
旧 & jiù & adj. & ancien, vieux, usagé (choses) \\
\hline
老 & lǎo & adj. & vieux, âgé (personnes) ; ancien (titre de respect) \\
\hline
可以 & kěyǐ & v. modal & pouvoir (permission), pouvoir (possibilité) \\
\hline
不可以 & bù kěyǐ & v. modal & ne pas avoir le droit de, interdiction \\
\hline
会 & huì & v. modal & savoir (compétence apprise) \\
\hline
能 & néng & v. modal & pouvoir (capacité physique, circonstancielle) \\
\hline
很 & hěn & adv. & très (liaison sujet-adjectif) \\
\hline
的 & de & part. & particule d'épithète / possession \\
\hline
辆 & liàng & class. & classificateur pour véhicules (vélo, voiture) \\
\hline
台 & tái & class. & classificateur pour machines, appareils (TV, ordi, téléphone) \\
\hline
部 & bù & class. & classificateur pour téléphone, voiture, film \\
\hline
\end{tabularx}
\end{center}
\end{definition}

\courssub{Fiche méthode 1 : Employer correctement l'antonyme 新 / 旧}

\begin{methode}[Opposer 新 (nouveau) et 旧 (ancien, usagé)]
\begin{enumerate}
\item \textbf{Identifier la nature du mot} : \zh{新} (xīn, « nouveau, neuf ») et \zh{旧} (jiù, « ancien, vieux, usagé ») sont des \cle{adjectifs} qui qualifient des \emph{objets, des lieux ou des concepts abstraits} (vêtements, machines, bâtiments, journaux, idées). En chinois, un adjectif seul ne peut pas être prédicat direct : il nécessite un adverbe de degré (\zh{很} hěn, \zh{真} zhēn, \zh{太} tài…) ou une négation (\zh{不} bù).
\item \textbf{Construire le prédicat (adjectif attribut)} : Structure : \textbf{Sujet + 很 + 新 / 旧}. L'adverbe \zh{很} (hěn) perd souvent son sens de « très » pour faire simple lien syntaxique.
\begin{itemize}
\item \zh{我的手机很新。} (Wǒ de shǒujī hěn xīn.) « Mon téléphone est neuf. »
\item \zh{这栋楼很旧。} (Zhè dòng lóu hěn jiù.) « Cet immeuble est vieux. »
\end{itemize}
\item \textbf{Nier correctement} : Structure : \textbf{Sujet + 不 + 新 / 旧}. \zh{不} (bù) précède directement l'adjectif.
\begin{itemize}
\item \zh{这本书不旧。} (Zhè běn shū bù jiù.) « Ce livre n'est pas vieux. »
\item \zh{他的想法不新。} (Tā de xiǎngfǎ bù xīn.) « Ses idées ne sont pas nouvelles. »
\end{itemize}
\item \textbf{Placer l'adjectif en épithète (devant un nom)} :
\begin{itemize}
\item \textbf{Adjectif monosyllabique} (新, 旧) : souvent \textbf{sans 的} directement devant le nom. Exemples : \zh{新手机} (xīn shǒujī, « nouveau téléphone »), \zh{旧电脑} (jiù diànnǎo, « vieil ordinateur »), \zh{新衣服} (xīn yīfu, « nouveaux vêtements »).
\item \textbf{Adjectif bisyllabique ou phrase adjectivale} : avec \textbf{的}. Exemple : \zh{崭新的书包} (zhǎnxīn de shūbāo, « cartable flambant neuf »).
\end{itemize}
\item \textbf{Choisir le bon antonyme selon l'animacité du nom} : \zh{旧} ne s'emploie \textbf{jamais} pour une personne vivante.
\begin{itemize}
\item \textbf{Chose / Objet} : opposition \zh{新} / \zh{旧}. \zh{新鞋} (chaussures neuves) vs \zh{旧鞋} (vieilles chaussures).
\item \textbf{Personne âgée} : on utilise \zh{老} (lǎo). \zh{老人} (personne âgée), \zh{老老师} (professeur âgé / titre de respect).
\item \textbf{Piège} : \zh{老师} (lǎoshī, « professeur ») est un mot composé figé où \zh{老} est un préfixe honorifique, non l'adjectif « vieux ». On ne dit jamais \zh{新师} ni \zh{旧师}.
\end{itemize}
\end{enumerate}
\textbf{Réflexe} : avant d'écrire, se poser deux questions : « Est-ce une chose ou une personne ? » et « L'adjectif est-il prédicat (avec 很/不) ou épithète (souvent sans 的 si monosyllabique) ? »
\end{methode}

\begin{exoresolu}[Choix entre 新, 旧 et 老]
\textbf{Énoncé} — Complétez avec \zh{新}, \zh{旧} ou \zh{老}, puis traduisez.\\
1. 我买了一辆\underline{\hspace{2cm}}自行车。\\
2. 他的电脑很\underline{\hspace{2cm}}，不能用了。\\
3. 我的\underline{\hspace{2cm}}师姓王。\\
4. 这是\underline{\hspace{2cm}}手机，昨天买的。
\tcblower
\textbf{Raisonnement} : on applique la règle de sélection. Les phrases 1, 2 et 4 qualifient des \emph{choses} (vélo, ordinateur, téléphone) : opposition \zh{新} / \zh{旧}. La phrase 3 qualifie une \emph{personne} (le professeur) : seul \zh{老} convient, d'autant que \zh{老师} (lǎoshī, « professeur ») est un mot figé.\\
\textbf{Solutions} :\\
1. \zh{我买了一辆新自行车。} (Wǒ mǎi le yí liàng xīn zìxíngchē.) « J'ai acheté un vélo neuf. » — épithète monosyllabique sans \zh{的} ; classificateur \zh{辆} (liàng) pour les véhicules.\\
2. \zh{他的电脑很旧，不能用了。} (Tā de diànnǎo hěn jiù, bù néng yòng le.) « Son ordinateur est très vieux, il ne marche plus. » — prédicat : présence obligatoire de \zh{很} ; \zh{不能} (bù néng) = incapacité physique/circonstancielle.\\
3. \zh{我的老师姓王。} (Wǒ de lǎoshī xìng Wáng.) « Mon professeur s'appelle Wang. » — \zh{旧师} n'existe pas : piège classique. \zh{老} ici est un préfixe honorifique dans le mot \zh{老师}.\\
4. \zh{这是新手机，昨天买的。} (Zhè shì xīn shǒujī, zuótiān mǎi de.) « C'est un nouveau téléphone, acheté hier. » — épithète \zh{新} sans \zh{的} ; structure \zh{是……的} pour insister sur le temps passé (\zh{昨天买的}).\\
\textbf{Conclusion} : vérifier toujours la nature du nom qualifié (animé / inanimé) avant de choisir l'antonyme.
\end{exoresolu}

\courssub{Fiche méthode 2 : Construire une phrase avec le verbe modal 可以}

\begin{methode}[La structure Sujet + 可以 + verbe]
\begin{enumerate}
\item \textbf{Mémoriser le schéma} : \textbf{Sujet + 可以 (kěyǐ) + Verbe (+ Complément)}. \zh{可以} exprime la \cle{permission} (« avoir le droit de », autorisation donnée par une règle, une personne, la loi) ou la \cle{possibilité} objective (« c'est possible de », conditions réunies).
\item \textbf{Placer 可以 AVANT le verbe principal}, jamais après.
\begin{itemize}
\item \zh{我可以进去吗？} (Wǒ kěyǐ jìnqù ma?) « Puis-je entrer ? » (Demande de permission)
\item \zh{这里可以停车。} (Zhèlǐ kěyǐ tíngchē.) « On peut se garer ici. » (Possibilité / Règle)
\end{itemize}
\item \textbf{Former la négation (interdiction)} : \textbf{Sujet + 不可以 + Verbe}. \zh{不} (bù) se place \textbf{devant} le modal \zh{可以}, jamais devant le verbe principal.
\begin{itemize}
\item \zh{教室里不可以打电话。} (Jiàoshì lǐ bù kěyǐ dǎ diànhuà.) « Il est interdit de téléphoner dans la salle de classe. »
\item \zh{未成年人不可以喝酒。} (Wèichéngniánrén bù kěyǐ hējiǔ.) « Les mineurs n'ont pas le droit de boire de l'alcool. »
\end{itemize}
\item \textbf{Former la question} : trois possibilités courantes.
\begin{itemize}
\item Particule \zh{吗} (ma) en fin de phrase : \zh{我可以用你的笔吗？} (Wǒ kěyǐ yòng nǐ de bǐ ma?)
\item Alternance \zh{可以不可以} + V (insistance sur le choix) : \zh{这件事可以不可以告诉他？}
\item Forme contractée \zh{可不可以} + V (plus courant à l'oral, plus poli) : \zh{可不可以借我一下？} (Kě bù kěyǐ jiè wǒ yíxià?)
\end{itemize}
\item \textbf{Répondre brièvement} : on reprend le modal seul, le verbe principal est omis.
\begin{itemize}
\item Affirmatif : \zh{可以。} (Kěyǐ.) « Oui, c'est permis / possible. »
\item Négatif : \zh{不可以。} (Bù kěyǐ.) « Non, c'est interdit / impossible. »
\end{itemize}
\item \textbf{Distinguer 可以 / 会 / 能} (point crucial pour le Bac) :
\begin{center}
\begin{tabularx}{\linewidth}{|X|X|X|}
\hline
\rowcolor{popblueL} \textbf{Modal} & \textbf{Valeur principale} & \textbf{Exemple} \\
\hline
\zh{可以} (kěyǐ) & Permission, autorisation, possibilité externe (règle, météo, feu vert) & \zh{雨停了，我们可以出去了。} (La pluie s'est arrêtée, on peut sortir.) \\
\hline
\zh{会} (huì) & Compétence apprise, savoir-faire, habileté & \zh{我会说中文。} (Je sais parler chinois.) \\
\hline
\zh{能} (néng) & Capacité physique, aptitude innée, possibilité circonstancielle (temps, place, santé) & \zh{他病了，不能来。} (Il est malade, il ne peut pas venir.) \\
\hline
\end{tabularx}
\end{center}
\end{enumerate}
\end{methode}

\begin{exoresolu}[Traduire avec le bon modal : 可以, 能, 会]
\textbf{Énoncé} — Traduisez en chinois avec \zh{可以}, \zh{能} ou \zh{会}.\\
1. Est-ce que je peux utiliser ton dictionnaire ? (demande de permission)\\
2. Mon petit frère sait nager. (compétence)\\
3. Aujourd'hui il ne pleut pas, on peut jouer au football. (circonstance météo)\\
4. Il est malade, il ne peut pas venir au lycée. (incapacité physique)
\tcblower
\textbf{Raisonnement} : on identifie la valeur du modal français « pouvoir / savoir ».\\
1. Permission demandée $\rightarrow$ \zh{可以} : \zh{我可以用你的词典吗？} (Wǒ kěyǐ yòng nǐ de cídiǎn ma?) — structure question avec \zh{吗} final.\\
2. Compétence apprise $\rightarrow$ \zh{会} : \zh{我弟弟会游泳。} (Wǒ dìdi huì yóuyǒng.) — \zh{可以} serait une erreur : nager est un savoir-faire, pas une permission.\\
3. Possibilité due aux circonstances (météo favorable) $\rightarrow$ \zh{可以} (ou \zh{能}) : \zh{今天不下雨，我们可以踢足球。} (Jīntiān bù xià yǔ, wǒmen kěyǐ tī zúqiú.) — \zh{可以} insiste sur la faisabilité externe.\\
4. Capacité physique empêchée par la maladie $\rightarrow$ \zh{能} : \zh{他病了，不能来学校。} (Tā bìng le, bù néng lái xuéxiào.) — \zh{不可以} signifierait « il n'a pas le droit », ce qui change le sens.\\
\textbf{Conclusion} : le choix du modal dépend du \emph{sens profond}, pas du mot français « pouvoir ». Traduire mentalement : « avoir le droit de » $\rightarrow$ \zh{可以} ; « savoir faire » $\rightarrow$ \zh{会} ; « être physiquement/circonstanciellement capable de » $\rightarrow$ \zh{能}。
\end{exoresolu}

\courssub{Fiche méthode 3 : Transformer des phrases (affirmation, négation, question, épithète)}

\begin{methode}[Transformations avec 可以 et les adjectifs 新 / 旧]
\begin{enumerate}
\item \textbf{Affirmation $\rightarrow$ Négation avec 可以} : insérer \zh{不} \textbf{devant} \zh{可以} (jamais devant le verbe principal seul).
\begin{itemize}
\item \zh{他可以上网。} $\rightarrow$ \zh{他不可以上网。} (Tā bù kěyǐ shàngwǎng.) « Il n'a pas le droit d'aller sur Internet. »
\item \textbf{Attention} : \zh{他可以不上网} (Tā kěyǐ bù shàngwǎng) signifie « Il a le droit de NE PAS aller sur Internet » (permission de s'abstenir), ce qui est l'opposé de l'interdiction.
\end{itemize}
\item \textbf{Affirmation $\rightarrow$ Question (particule 吗)} : ajouter \zh{吗} en fin de phrase, sans changer l'ordre des mots.
\begin{itemize}
\item \zh{你可以进来。} $\rightarrow$ \zh{你可以进来吗？} (Nǐ kěyǐ jìnlái ma?) « Peux-tu entrer ? »
\end{itemize}
\item \textbf{Question $\rightarrow$ Réponse courte} : reprendre le modal seul.
\begin{itemize}
\item \zh{—— 我可以用这个吗？} \zh{—— 可以。} / \zh{—— 不可以。}
\end{itemize}
\item \textbf{Adjectif prédicat $\rightarrow$ Épithète (phrase avec 是……的 ou simple NP)} : ajouter classificateur + (的) + nom.
\begin{itemize}
\item \zh{这台电视机很旧。} $\rightarrow$ \zh{这是一台旧电视机。} (Zhè shì yì tái jiù diànshìjī.) « C'est un vieux téléviseur. »
\item \textbf{Règle du classificateur} : \zh{一} (yī) + \textbf{Classificateur} + \textbf{Adjectif} + \textbf{Nom}. Pour les appareils (TV, ordi, téléphone, machine à laver) : \zh{台} (tái) ou \zh{部} (bù, pour téléphone/voiture). L'adjectif monosyllabique (\zh{旧}, \zh{新}) se passe généralement de \zh{的}.
\end{itemize}
\item \textbf{Vérification finale} : relire en contrôlant l'ordre \textbf{Sujet + (不) + 可以 + Verbe}, la place des compléments (Lieu/Temps avant le verbe), et la ponctuation chinoise (。 ？).
\end{enumerate}
\end{methode}

\begin{exoresolu}[Atelier de transformations]
\textbf{Énoncé} — Transformez selon la consigne.\\
1. \zh{学生可以在图书馆看书。} $\rightarrow$ mettez à la forme négative (interdiction).\\
2. \zh{我可以用你的手机。} $\rightarrow$ transformez en question avec \zh{吗} et donnez la réponse affirmative courte.\\
3. \zh{他的电脑很旧。} $\rightarrow$ transformez en phrase nominale avec \zh{是……的} : « C'est un vieil ordinateur. »
\tcblower
\textbf{Solutions détaillées} :\\
1. Négation : \zh{不} se place devant le modal \zh{可以}, pas devant \zh{看} : \zh{学生不可以在图书馆看书。} (Xuésheng bù kěyǐ zài túshūguǎn kàn shū.) « Les élèves n'ont pas le droit de lire à la bibliothèque. » Le complément de lieu \zh{在图书馆} reste avant le verbe \zh{看}.\\
2. Question : \zh{我可以用你的手机吗？} (Wǒ kěyǐ yòng nǐ de shǒujī ma?) « Puis-je utiliser ton téléphone ? » Réponse courte : \zh{可以。} (Kěyǐ.) « Oui, tu peux. » — on ne répète pas le verbe \zh{用}.\\
3. Épithète / Phrase nominale : \zh{这是一台旧电脑。} (Zhè shì yì tái jiù diànnǎo.) « C'est un vieil ordinateur. » — classificateur \zh{台} (tái) obligatoire entre \zh{一} et le nom ; l'adjectif monosyllabique \zh{旧} se passe de \zh{的}.\\
\textbf{Conclusion} : chaque transformation ne touche qu'un seul point de la phrase (négation, interrogation, nominalisation) ; le reste de l'ordre des mots (Sujet - Temps - Lieu - Modal - Verbe - Objet) est conservé.
\end{exoresolu}

\courssub{Fiche méthode 4 : Rédiger un court texte avec 新 / 旧 et 可以}

\begin{methode}[Produire un paragraphe cohérent de 5-6 phrases (Module : Mass médias)]
\begin{enumerate}
\item \textbf{Choisir un sujet concret} lié aux médias / télécoms : un smartphone, un ordinateur, la connexion Internet au lycée, la radio, une application.
\item \textbf{Planifier en trois temps logiques} :
\begin{enumerate}
\item \textbf{Description} : présentation de l'objet (neuf ou vieux ?) $\rightarrow$ \zh{很新} / \zh{很旧} / \zh{新……} / \zh{旧……}。
\item \textbf{Fonctions / Permissions} : ce qu'on \emph{peut} faire avec (\zh{可以} + verbes d'action : \zh{上网}, \zh{看新闻}, \zh{听音乐}, \zh{发微信}, \zh{拍照})。
\item \textbf{Contrainte / Contraste / Conclusion} : interdiction (\zh{不可以}), comparaison avec l'ancien (\zh{旧} vs \zh{新}), sentiment final.
\end{enumerate}
\item \textbf{Varier les structures obligatoires} (critères de notation) :
\begin{itemize}
\item Au moins une phrase avec adjectif prédicat (\zh{很新} / \zh{很旧}).
\item Au moins une phrase avec \zh{可以} affirmatif.
\item Au moins une phrase avec \zh{不可以} ou question \zh{可以吗}。
\item Utilisation correcte des classificateurs : \zh{一部/一台手机} (yí bù / yí tái shǒujī), \zh{一台电脑} (yí tái diànnǎo), \zh{一辆车} (yí liàng chē).
\end{itemize}
\item \textbf{Relire (check-list)} : tons du pinyin (surtout 3e ton \zh{可以} kěyǐ, \zh{旧} jiù), ordre Sujet + Modal + Verbe, place du complément de temps/lieu \textbf{avant} le modal ou devant le sujet, ponctuation pleine chasse (。 ， ？ ！).
\end{enumerate}
\end{methode}

\begin{exoresolu}[Composition guidée : Mon nouveau téléphone]
\textbf{Énoncé} — Rédigez 5 phrases en chinois sur le thème « Mon nouveau téléphone » en utilisant au moins deux fois \zh{可以} et une fois \zh{新} ou \zh{旧}.
\tcblower
\textbf{Démarche} : présentation (phrase 1) $\rightarrow$ usages permis (phrases 2-3) $\rightarrow$ interdiction contexte scolaire (phrase 4) $\rightarrow$ conclusion comparative (phrase 5).\\
\textbf{Production modèle} :\\
1. \zh{爸爸给我买了一部新手机。} (Bàba gěi wǒ mǎi le yí bù xīn shǒujī.) « Papa m'a acheté un nouveau téléphone. » — \zh{给} (gěi) marque le bénéficiaire ; classificateur \zh{部} (bù) pour téléphone ; \zh{新} épithète sans \zh{的}。\\
2. \zh{我可以用它打电话，也可以上网。} (Wǒ kěyǐ yòng tā dǎ diànhuà, yě kěyǐ shàngwǎng.) « Je peux l'utiliser pour téléphoner, et aussi pour aller sur Internet. » — \zh{可以} répété pour deux verbes différents ; \zh{也} (yě) « aussi ».\\
3. \zh{我还可以看新闻、听音乐。} (Wǒ hái kěyǐ kàn xīnwén, tīng yīnyuè.) « Je peux aussi lire les informations et écouter de la musique. » — \zh{还} (hái) « en plus » ; énumération de verbes \zh{看}, \zh{听} partageant \zh{可以}。\\
4. \zh{上课的时候，我们不可以用手机。} (Shàngkè de shíhou, wǒmen bù kěyǐ yòng shǒujī.) « Pendant les cours, nous n'avons pas le droit d'utiliser le téléphone. » — Complément de temps \zh{上课的时候} en tête de phrase (sujet \zh{我们}) ; \zh{不可以} = interdiction.\\
5. \zh{我的旧手机太慢了，新手机很快，我很高兴。} (Wǒ de jiù shǒujī tài màn le, xīn shǒujī hěn kuài, wǒ hěn gāoxìng.) « Mon vieux téléphone était trop lent, le nouveau est rapide, je suis content. » — Opposition \zh{旧} / \zh{新} ; \zh{太……了} (tài…le) « trop… » ; adjectifs prédicats \zh{慢} (avec \zh{太}), \zh{快} (avec \zh{很}).\\
\textbf{Justifications} : classificateur \zh{部} avec \zh{手机} ; modal \zh{可以} placé avant chaque verbe (\zh{打}, \zh{上}, \zh{看}, \zh{听}) ; complément de temps \zh{上课的时候} en tête de phrase ; opposition finale \zh{旧} / \zh{新} pour le contraste.\\
\textbf{Conclusion} : un texte court mais varié, avec les deux structures de la leçon employées dans des fonctions différentes (description, permission, interdiction, comparaison).
\end{exoresolu}

\courssub{Bilan et points de vigilance pour l'évaluation}
\begin{attention}[Les 5 erreurs qui coûtent des points au Bac / Contrôle]
\begin{enumerate}
\item \textbf{Dire \zh{旧} pour une personne} : \zh{旧} ne qualifie que les choses inanimées. Pour une personne âgée, c'est \zh{老} (lǎo) ou \zh{年纪大} (niánjì dà).
\begin{itemize}
\item \textbf{Faux} : \zh{我的旧老师} (mon « vieux » prof — sens : prof usagé).
\item \textbf{Juste} : \zh{我的老师年纪大了} (mon prof est âgé) ou \zh{老老师} (professeur expérimenté / titre de respect).
\end{itemize}
\item \textbf{Mal placer la négation avec \zh{可以}} : le calque du français « il peut ne pas venir » (ambiguïté permission/possibilité) pousse à écrire \zh{他可以不回来}.
\begin{itemize}
\item \textbf{Interdiction} (négation du modal) : \zh{他不可以回来。} (Tā bù kěyǐ huílái.) « Il n'a pas le droit de revenir. »
\item \textbf{Permission de ne pas faire} (négation du verbe) : \zh{他可以不回来。} (Tā kěyǐ bù huílái.) « Il a le droit de ne pas revenir. »
\item \textbf{Règle} : pour interdire, \zh{不} \textbf{précède} \zh{可以}.
\end{itemize}
\item \textbf{Oublier \zh{很} (ou autre adverbe de degré) devant l'adjectif prédicat} : \zh{我的手机新} est syntaxiquement incomplet (sonne comme une ellipse de comparaison : « Mon téléphone [est] nouveau [comparé à l'autre] »).
\begin{itemize}
\item \textbf{Correct} : \zh{我的手机很新。} (Wǒ de shǒujī hěn xīn.)
\item \textbf{Correct} : \zh{我的手机真新！} (Zhēn xīn !) « Mon téléphone est vraiment neuf ! »
\end{itemize}
\item \textbf{Confondre \zh{可以}, \zh{会} et \zh{能}} : « Je sais nager » se traduit \zh{我会游泳}, jamais \zh{我可以游泳} (qui signifie « j'ai le droit de nager / la piscine est ouverte »).
\begin{itemize}
\item \zh{可以} = permission / possibilité externe (règle, condition météo, autorisation parentale).
\item \zh{会} = compétence apprise (savoir-faire : nager, conduire, parler une langue, coder).
\item \zh{能} = capacité physique / circonstancielle (force, santé, temps, place, argent).
\end{itemize}
\item \textbf{Oublier le classificateur dans la transformation « Prédicat $\rightarrow$ Épithète »} : « C'est un vieil ordinateur » exige \zh{这是一台旧电脑} avec \zh{台}.
\begin{itemize}
\item \textbf{Faute} : \zh{这是一个旧电脑} (\zh{个} ne convient pas aux appareils électroniques).
\item \textbf{Faute} : \zh{这是旧电脑} (phrase nominale incomplète sans démonstratif + numéral + classificateur).
\item \textbf{Correct} : \zh{这是一台旧电脑。} / \zh{这是一部旧电脑。} (\zh{部} possible pour ordi portable).
\end{itemize}
\end{enumerate}
\end{attention}

\newpage

\coursec{Exercices}

\courssub{Je vérifie mes connaissances}

\begin{exercice}[QCM : la bonne phrase]
Pour chaque question, entoure la phrase correcte (une seule réponse).

\textbf{1.} « Je peux utiliser l'ordinateur. »
\begin{itemize}
\item[A.] 我用可以电脑。(Wǒ yòng kěyǐ diànnǎo.)
\item[B.] 我可以用电脑。(Wǒ kěyǐ yòng diànnǎo.)
\item[C.] 我可以电脑用。(Wǒ kěyǐ diànnǎo yòng.)
\item[D.] 可以我用电脑。(Kěyǐ wǒ yòng diànnǎo.)
\end{itemize}

\textbf{2.} « C'est mon ancien téléphone. »
\begin{itemize}
\item[A.] 这是我旧手机。(Zhè shì wǒ jiù shǒujī.)
\item[B.] 这是我的旧手机。(Zhè shì wǒ de jiù shǒujī.)
\item[C.] 这是我旧的。(Zhè shì wǒ jiù de.)
\item[D.] 这是旧我的手机。(Zhè shì jiù wǒ de shǒujī.)
\end{itemize}

\textbf{3.} « On ne peut pas utiliser le téléphone en classe. »
\begin{itemize}
\item[A.] 在课堂上没可以用手机。(Zài kètáng shang méi kěyǐ yòng shǒujī.)
\item[B.] 在课堂上不可以手机用。(Zài kètáng shang bù kěyǐ shǒujī yòng.)
\item[C.] 在课堂上不可以用手机。(Zài kètáng shang bù kěyǐ yòng shǒujī.)
\item[D.] 在课堂上可以用不手机。(Zài kètáng shang kěyǐ yòng bù shǒujī.)
\end{itemize}

\textbf{4.} « Ma nouvelle école est grande. »
\begin{itemize}
\item[A.] 我的新学校很大。(Wǒ de xīn xuéxiào hěn dà.)
\item[B.] 我新学校的很大。(Wǒ xīn xuéxiào de hěn dà.)
\item[C.] 我的学校新很大。(Wǒ de xuéxiào xīn hěn dà.)
\item[D.] 新我的学校很大。(Xīn wǒ de xuéxiào hěn dà.)
\end{itemize}
\end{exercice}

\begin{exercice}[Vrai ou faux ? Justifie ta réponse]
Dis si chaque phrase est correcte (V) ou incorrecte (F), puis justifie en français en citant la règle de grammaire concernée.

\begin{enumerate}
\item \zh{这本书很旧。}(Zhè běn shū hěn jiù.) — « Ce livre est très vieux. »
\item \zh{他不可以说汉语。}(Tā bù kěyǐ shuō Hànyǔ.) — « Il ne peut pas parler chinois. » (sens : interdiction)
\item \zh{我可以去你家吗？}(Wǒ kěyǐ qù nǐ jiā ma ?) — « Est-ce que je peux aller chez toi ? »
\item \zh{我没可以看电视。}(Wǒ méi kěyǐ kàn diànshì.) — « Je ne peux pas regarder la télévision. »
\item \zh{老师的新电脑可以上网。}(Lǎoshī de xīn diànnǎo kěyǐ shàngwǎng.) — « Le nouvel ordinateur du professeur permet d'aller sur Internet. »
\end{enumerate}
\end{exercice}

\begin{exercice}[Vocabulaire : complète le tableau]
Complète le tableau suivant en écrivant les éléments manquants (caractères, pinyin ou traduction).

\begin{center}\begin{tabularx}{\linewidth}{|X|X|X|X|}\hline
\rowcolor{popblueL} Caractères & Pinyin & Nature & Traduction \\ \hline
新 & \ldots\ldots\ldots & adjectif & \ldots\ldots\ldots \\ \hline
\ldots\ldots\ldots & jiù & adjectif & vieux, ancien (chose) \\ \hline
可以 & \ldots\ldots\ldots & verbe modal & \ldots\ldots\ldots \\ \hline
\ldots\ldots\ldots & shǒujī & nom & téléphone portable \\ \hline
用 & \ldots\ldots\ldots & verbe & \ldots\ldots\ldots \\ \hline
\ldots\ldots\ldots & shàngwǎng & verbe & aller sur Internet \\ \hline
旧书 & jiù shū & locution nominale & \ldots\ldots\ldots \\ \hline
\ldots\ldots\ldots & xīn tóngxué & locution nominale & nouveau camarade de classe \\ \hline
\end{tabularx}\end{center}
\end{exercice}

\begin{exercice}[Associe le pinyin aux caractères]
Relie chaque mot en caractères à son pinyin correct.

\begin{center}
\begin{tabular}{ll}
\textbf{Caractères} & \textbf{Pinyin} \\
1. 新手机 & a. kěyǐ \\
2. 旧电脑 & b. bù kěyǐ \\
3. 可以 & c. xīn shǒujī \\
4. 不可以 & d. kè táng \\
5. 课堂 & e. jiù diànnǎo \\
\end{tabular}
\end{center}

Puis traduis chaque mot en français.
\end{exercice}

\courssub{Je m'entraîne}

\begin{exercice}[Traduis en français]
Traduis les phrases suivantes en français.

\begin{enumerate}
\item \zh{这是我的旧手机，那是我的新手机。}(Zhè shì wǒ de jiù shǒujī, nà shì wǒ de xīn shǒujī.)
\item \zh{我可以用电脑上网。}(Wǒ kěyǐ yòng diànnǎo shàngwǎng.)
\item \zh{在课堂上，学生不可以用手机。}(Zài kètáng shang, xuésheng bù kěyǐ yòng shǒujī.)
\item \zh{他的旧自行车还可以用。}(Tā de jiù zìxíngchē hái kěyǐ yòng.)
\item \zh{妈妈给我买了一个新书包。}(Māma gěi wǒ mǎi le yí ge xīn shūbāo.)
\end{enumerate}
\end{exercice}

\begin{exercice}[Transformation : négation et question]
Pour chaque phrase : a) mets-la à la forme négative ; b) mets-la à la forme interrogative avec 吗. Écris les caractères, le pinyin et la traduction des deux phrases obtenues.

Exemple : 我可以用电脑。→ négation : 我不可以用电脑。→ question : 你可以用电脑吗？

\begin{enumerate}
\item \zh{我可以用电脑。}(Wǒ kěyǐ yòng diànnǎo.)
\item \zh{学生可以在课堂上吃东西。}(Xuésheng kěyǐ zài kètáng shang chī dōngxi.)
\item \zh{我们可以看电视。}(Wǒmen kěyǐ kàn diànshì.)
\item \zh{你可以用我的新手机。}(Nǐ kěyǐ yòng wǒ de xīn shǒujī.)
\item \zh{他可以今天来学校。}(Tā kěyǐ jīntiān lái xuéxiào.)
\end{enumerate}
\end{exercice}

\begin{exercice}[Texte à trous]
Complète le texte suivant avec les mots de la liste : \zh{新}、\zh{旧}、\zh{可以}、\zh{不可以}、\zh{的}、\zh{用}.

\begin{textebox}[Chez Ali, à Yaoundé]
我叫阿里，我住在雅温得。我有一台\_\_\_\_电脑，它很\_\_\_\_，已经五年了，但是还\_\_\_\_用。昨天爸爸给我买了一台\_\_\_\_电脑，我很高兴！在学校，我们\_\_\_\_在课上\_\_\_\_手机，这是学校\_\_\_\_规定。\\[4pt]
Wǒ jiào Ālǐ, wǒ zhù zài Yǎwēndé. Wǒ yǒu yì tái \_\_\_\_ diànnǎo, tā hěn \_\_\_\_, yǐjīng wǔ nián le, dànshì hái \_\_\_\_ yòng. Zuótiān bàba gěi wǒ mǎi le yì tái \_\_\_\_ diànnǎo, wǒ hěn gāoxìng ! Zài xuéxiào, wǒmen \_\_\_\_ zài kè shang \_\_\_\_ shǒujī, zhè shì xuéxiào \_\_\_\_ guīdìng.\\[4pt]
« Je m'appelle Ali, j'habite à Yaoundé. J'ai un ordinateur, il est très \ldots, il a déjà cinq ans, mais on peut encore l'utiliser. Hier, papa m'a acheté un ordinateur \ldots, je suis très content ! À l'école, nous \ldots utiliser le téléphone en classe, c'est le règlement de l'école. »
\end{textebox}
\end{exercice}

\begin{exercice}[Remets les mots dans l'ordre]
Remets les mots dans le bon ordre pour faire une phrase correcte, puis traduis-la.

\begin{enumerate}
\item 可以 / 我 / 上网 / 用 / 新手机
\item 旧 / 这 / 的 / 是 / 书 / 我
\item 不 / 课堂上 / 在 / 可以 / 手机 / 用 / 学生
\item 新 / 很 / 学校 / 我们 / 的 / 大
\item 吗 / 你 / 可以 / 去 / 我家 / 今天
\end{enumerate}
\end{exercice}

\courssub{Je me prépare au Bac}

\begin{exercice}[Type Bac — Compréhension de l'écrit et questions de langue]
Lis le texte suivant, puis réponds aux questions.

\begin{textebox}[Le téléphone de Marie]
我叫玛丽，我是雅温得一所中学的学生。去年，妈妈给我买了一个新手机，我的旧手机给了弟弟。新手机很好，我可以用它上网、看新闻，也可以学习汉语。但是在学校，我们不可以在课堂上用手机。老师说：“手机可以在家里用，不可以在课上用。”我觉得老师说得对。\\[4pt]
Wǒ jiào Mǎlì, wǒ shì Yǎwēndé yì suǒ zhōngxué de xuésheng. Qùnián, māma gěi wǒ mǎi le yí ge xīn shǒujī, wǒ de jiù shǒujī gěi le dìdi. Xīn shǒujī hěn hǎo, wǒ kěyǐ yòng tā shàngwǎng, kàn xīnwén, yě kěyǐ xuéxí Hànyǔ. Dànshì zài xuéxiào, wǒmen bù kěyǐ zài kètáng shang yòng shǒujī. Lǎoshī shuō : “Shǒujī kěyǐ zài jiā lǐ yòng, bù kěyǐ zài kè shang yòng.” Wǒ juéde lǎoshī shuō de duì.\\[4pt]
« Je m'appelle Marie, je suis élève dans un collège de Yaoundé. L'année dernière, maman m'a acheté un nouveau téléphone, et j'ai donné mon ancien téléphone à mon petit frère. Le nouveau téléphone est très bien, je peux l'utiliser pour aller sur Internet, lire les informations, et aussi pour apprendre le chinois. Mais à l'école, nous ne pouvons pas utiliser le téléphone en classe. Le professeur dit : "On peut utiliser le téléphone à la maison, on ne peut pas l'utiliser en cours." Je trouve que le professeur a raison. »
\end{textebox}

\textbf{A. Questions de compréhension} (réponds en français)
\begin{enumerate}
\item Qui a acheté le nouveau téléphone de Marie ?
\item Qu'a fait Marie de son ancien téléphone ?
\item Cite deux choses que Marie peut faire avec son nouveau téléphone.
\item Quelle est la règle de l'école concernant les téléphones ?
\end{enumerate}

\textbf{B. Questions de langue}
\begin{enumerate}
\item Relève dans le texte deux phrases contenant \zh{可以} et une phrase contenant \zh{不可以}, puis traduis-les.
\item Dans la phrase \zh{我可以用它上网}, quel est le verbe principal ? Où se place le verbe modal \zh{可以} ?
\item Trouve dans le texte un exemple de \zh{新} et un exemple de \zh{旧}, et explique la différence de sens.
\item Mets à la forme interrogative avec \zh{吗} : \zh{我可以用它上网。}
\end{enumerate}
\end{exercice}

\begin{exercice}[Type Bac — Thème et version]
\textbf{A. Thème.} Traduis les phrases suivantes en chinois (caractères obligatoires ; le pinyin est apprécié).

\begin{enumerate}
\item Les élèves ne peuvent pas utiliser leurs nouveaux téléphones en classe.
\item Mon ancien ordinateur est très vieux, mais je peux encore l'utiliser.
\item Est-ce que je peux aller sur Internet avec ton nouveau téléphone ?
\item C'est le nouveau livre de mon professeur de chinois.
\item À la maison, nous pouvons regarder la télévision ; à l'école, nous ne pouvons pas.
\end{enumerate}

\textbf{B. Version.} Traduis le texte suivant en français.

\begin{textebox}[À traduire]
这是我的旧手机，我可以用新手机上网。新手机是妈妈给我买的。在学校里，我们不可以用手机，但是在家可以。\\[4pt]
Zhè shì wǒ de jiù shǒujī, wǒ kěyǐ yòng xīn shǒujī shàngwǎng. Xīn shǒujī shì māma gěi wǒ mǎi de. Zài xuéxiào lǐ, wǒmen bù kěyǐ yòng shǒujī, dànshì zài jiā kěyǐ.
\end{textebox}
\end{exercice}

\begin{exercice}[Type Bac — Correction de texte et expression écrite]
\textbf{A. Correction de texte.} Le paragraphe suivant contient \textbf{5 erreurs} (ordre des mots, mot-outil oublié, négation incorrecte, mauvais caractère). Relève chaque erreur, écris la forme corrigée et justifie la correction en français.

\begin{textebox}[Texte à corriger]
我用可以新手机。这是我的旧的书。在课堂上，学生没可以用手机。我旧电脑还可以上网。老师的新很大学校。\\[4pt]
Wǒ yòng kěyǐ xīn shǒujī. Zhè shì wǒ de jiù de shū. Zài kètáng shang, xuésheng méi kěyǐ yòng shǒujī. Wǒ jiù diànnǎo hái kěyǐ shàngwǎng. Lǎoshī de xīn hěn dà xuéxiào.
\end{textebox}

\textbf{B. Expression écrite guidée.} Rédige un court texte en chinois de \textbf{60 à 80 caractères} sur le sujet suivant : « Mon vieux téléphone et mon nouveau téléphone » (我的旧手机和新手机).

Consignes obligatoires :
\begin{itemize}
\item utilise au moins \textbf{trois fois} \zh{可以} ou \zh{不可以} ;
\item utilise au moins \textbf{deux fois} les antonymes \zh{新} et \zh{旧} ;
\item utilise au moins une fois la structure \zh{用 + objet + verbe} (ex. \zh{用手机上网}) ;
\item écris en caractères, avec le pinyin en dessous de chaque phrase.
\end{itemize}

Piste de départ : \zh{我有两个手机。}(Wǒ yǒu liǎng ge shǒujī.)
\end{exercice}

\newpage

\coursec{Corrigés des exercices}

\begin{corrige}[Exercice 1 — QCM : la bonne phrase]
\textbf{1. Réponse B.} \zh{我可以用电脑。}(Wǒ kěyǐ yòng diànnǎo.) — « Je peux utiliser l'ordinateur. »
Justification : le verbe modal \zh{可以} se place \textbf{avant} le verbe principal, jamais après : Sujet + 可以 + verbe. Les phrases A, C et D violent l'ordre des mots chinois.

\textbf{2. Réponse B.} \zh{这是我的旧手机。}(Zhè shì wǒ de jiù shǒujī.) — « C'est mon ancien téléphone. »
Justification : entre le possessif \zh{我} et le nom, il faut le mot-outil \zh{的} ; l'adjectif \zh{旧} se place entre \zh{的} et le nom : 我 + 的 + 旧 + 手机. En A, \zh{的} manque ; en D, \zh{旧} est mal placé.

\textbf{3. Réponse C.} \zh{在课堂上不可以用手机。}(Zài kètáng shang bù kěyǐ yòng shǒujī.) — « On ne peut pas utiliser le téléphone en classe. »
Justification : la négation de \zh{可以} est \zh{不可以} (jamais \zh{没可以}) ; le complément de lieu \zh{在课堂上} se place en tête de phrase.

\textbf{4. Réponse A.} \zh{我的新学校很大。}(Wǒ de xīn xuéxiào hěn dà.) — « Ma nouvelle école est grande. »
Justification : schéma 我 + 的 + 新 + 学校 ; l'adjectif \zh{新} précède le nom qu'il qualifie, et le prédicat adjectival \zh{很大} se place après le nom (Sujet + 很 + adjectif).
\end{corrige}

\begin{corrige}[Exercice 2 — Vrai ou faux ?]
\begin{enumerate}
\item \textbf{V.} \zh{这本书很旧。} — Correct : adjectif attribut du prédicat précédé de \zh{很} (Sujet + 很 + adjectif). \zh{旧} = vieux, ancien (pour une chose).
\item \textbf{V.} \zh{他不可以说汉语。} — Correct : la négation de \zh{可以} est \zh{不 + 可以}, placée avant le verbe modal ; sens : interdiction (« il n'a pas le droit de parler chinois »).
\item \textbf{V.} \zh{我可以去你家吗？} — Correct : question en \zh{吗} ajouté en fin de phrase affirmative, sans changer l'ordre des mots.
\item \textbf{F.} \zh{我没可以看电视。} — Incorrect : \zh{可以} ne se nie jamais avec \zh{没} ; il faut \zh{不}. Correction : \zh{我不可以看电视。}(Wǒ bù kěyǐ kàn diànshì.) — « Je ne peux pas regarder la télévision. »
\item \textbf{V.} \zh{老师的新电脑可以上网。} — Correct : 老师 + 的 + 新 + 电脑 (possessif + 的 + adjectif + nom), puis 可以 + verbe \zh{上网}.
\end{enumerate}
\end{corrige}

\begin{corrige}[Exercice 3 — Vocabulaire : complète le tableau]
\begin{center}\begin{tabularx}{\linewidth}{|X|X|X|X|}\hline
\rowcolor{popblueL} Caractères & Pinyin & Nature & Traduction \\ \hline
新 & xīn & adjectif & nouveau, neuf \\ \hline
旧 & jiù & adjectif & vieux, ancien (chose) \\ \hline
可以 & kěyǐ & verbe modal & pouvoir, avoir le droit de \\ \hline
手机 & shǒujī & nom & téléphone portable \\ \hline
用 & yòng & verbe & utiliser \\ \hline
上网 & shàngwǎng & verbe & aller sur Internet \\ \hline
旧书 & jiù shū & locution nominale & vieux livre, livre d'occasion \\ \hline
新同学 & xīn tóngxué & locution nominale & nouveau camarade de classe \\ \hline
\end{tabularx}\end{center}
Point de vigilance : \zh{旧} (jiù, quatrième ton) ne se confond pas avec \zh{就} (jiù, « alors, aussitôt »), homophone au sens totalement différent. Autre piège : \zh{旧} est l'antonyme de \zh{新} uniquement pour les \textbf{choses} ; pour une personne âgée, on dit \zh{老} (lǎo), jamais \zh{旧}.
\end{corrige}

\begin{corrige}[Exercice 4 — Associe le pinyin aux caractères]
Correspondances : \textbf{1-c, 2-e, 3-a, 4-b, 5-d.}
\begin{itemize}
\item 1. \zh{新手机} — xīn shǒujī — « nouveau téléphone (portable) »
\item 2. \zh{旧电脑} — jiù diànnǎo — « vieil ordinateur »
\item 3. \zh{可以} — kěyǐ — « pouvoir, avoir le droit de »
\item 4. \zh{不可以} — bù kěyǐ — « ne pas pouvoir, ne pas avoir le droit de »
\item 5. \zh{课堂} — kètáng — « classe, cours (salle de classe) »
\end{itemize}
\end{corrige}

\begin{corrige}[Exercice 5 — Traduis en français]
\begin{enumerate}
\item \zh{这是我的旧手机，那是我的新手机。} — « Ceci est mon ancien téléphone, cela est mon nouveau téléphone. »
\item \zh{我可以用电脑上网。} — « Je peux utiliser l'ordinateur pour aller sur Internet. »
\item \zh{在课堂上，学生不可以用手机。} — « En classe, les élèves ne peuvent pas utiliser le téléphone. »
\item \zh{他的旧自行车还可以用。} — « Son vieux vélo peut encore servir (on peut encore l'utiliser). »
\item \zh{妈妈给我买了一个新书包。} — « Maman m'a acheté un nouveau cartable. » (注意 : \zh{给} + personne + verbe = « faire quelque chose pour quelqu'un ».)
\end{enumerate}
\end{corrige}

\begin{corrige}[Exercice 6 — Transformation : négation et question]
Rappel : négation = \zh{不} devant \zh{可以} ; question = phrase affirmative + \zh{吗} (souvent avec passage 我 → 你).

\begin{enumerate}
\item Négation : \zh{我不可以用电脑。}(Wǒ bù kěyǐ yòng diànnǎo.) — « Je ne peux pas utiliser l'ordinateur. »\\
Question : \zh{你可以用电脑吗？}(Nǐ kěyǐ yòng diànnǎo ma ?) — « Peux-tu utiliser l'ordinateur ? »
\item Négation : \zh{学生不可以在课堂上吃东西。}(Xuésheng bù kěyǐ zài kètáng shang chī dōngxi.) — « Les élèves ne peuvent pas manger en classe. »\\
Question : \zh{学生可以在课堂上吃东西吗？}(Xuésheng kěyǐ zài kètáng shang chī dōngxi ma ?) — « Les élèves peuvent-ils manger en classe ? »
\item Négation : \zh{我们不可以看电视。}(Wǒmen bù kěyǐ kàn diànshì.) — « Nous ne pouvons pas regarder la télévision. »\\
Question : \zh{你们可以看电视吗？}(Nǐmen kěyǐ kàn diànshì ma ?) — « Pouvez-vous regarder la télévision ? »
\item Négation : \zh{你不可以用我的新手机。}(Nǐ bù kěyǐ yòng wǒ de xīn shǒujī.) — « Tu ne peux pas utiliser mon nouveau téléphone. »\\
Question : \zh{我可以用你的新手机吗？}(Wǒ kěyǐ yòng nǐ de xīn shǒujī ma ?) — « Est-ce que je peux utiliser ton nouveau téléphone ? »
\item Négation : \zh{他今天不可以来学校。}(Tā jīntiān bù kěyǐ lái xuéxiào.) — « Il ne peut pas venir à l'école aujourd'hui. »\\
Question : \zh{他今天可以来学校吗？}(Tā jīntiān kěyǐ lái xuéxiào ma ?) — « Peut-il venir à l'école aujourd'hui ? »
\end{enumerate}
Point de vigilance : ne jamais écrire \zh{没可以} ; ne pas déplacer \zh{可以} après le verbe ; le complément de temps \zh{今天} se place juste après le sujet, avant \zh{可以}.
\end{corrige}

\begin{corrige}[Exercice 7 — Texte à trous]
Texte complété :\\
\zh{我叫阿里，我住在雅温得。我有一台\textbf{旧}电脑，它很\textbf{旧}，已经五年了，但是还\textbf{可以}用。昨天爸爸给我买了一台\textbf{新}电脑，我很高兴！在学校，我们\textbf{不可以}在课上\textbf{用}手机，这是学校\textbf{的}规定。}\\[4pt]
Wǒ jiào Ālǐ, wǒ zhù zài Yǎwēndé. Wǒ yǒu yì tái \textbf{jiù} diànnǎo, tā hěn \textbf{jiù}, yǐjīng wǔ nián le, dànshì hái \textbf{kěyǐ} yòng. Zuótiān bàba gěi wǒ mǎi le yì tái \textbf{xīn} diànnǎo, wǒ hěn gāoxìng ! Zài xuéxiào, wǒmen \textbf{bù kěyǐ} zài kè shang \textbf{yòng} shǒujī, zhè shì xuéxiào \textbf{de} guīdìng.\\[4pt]
Traduction : « Je m'appelle Ali, j'habite à Yaoundé. J'ai un vieil ordinateur, il est très vieux, il a déjà cinq ans, mais on peut encore l'utiliser. Hier, papa m'a acheté un nouvel ordinateur, je suis très content ! À l'école, nous ne pouvons pas utiliser le téléphone en classe, c'est le règlement de l'école. »

Justifications : trous 1 et 2 : antonyme \zh{旧} (contexte « cinq ans ») ; trou 3 : \zh{可以} + verbe \zh{用} ; trou 4 : \zh{新} (opposé logique, « acheté hier ») ; trou 5 : \zh{不可以} (interdiction scolaire) ; trou 6 : verbe \zh{用} ; trou 7 : mot-outil \zh{的} entre le nom \zh{学校} et \zh{规定}. Rappel : le classificateur de \zh{电脑} est \zh{台} (一台电脑).
\end{corrige}

\begin{corrige}[Exercice 8 — Remets les mots dans l'ordre]
\begin{enumerate}
\item \zh{我可以用新手机上网。}(Wǒ kěyǐ yòng xīn shǒujī shàngwǎng.) — « Je peux utiliser le nouveau téléphone pour aller sur Internet. » (Sujet + 可以 + 用 + objet + verbe)
\item \zh{这是我的旧书。}(Zhè shì wǒ de jiù shū.) — « C'est mon vieux livre. » (这 + 是 + 我 + 的 + 旧 + 书)
\item \zh{在课堂上学生不可以用手机。}(Zài kètáng shang xuésheng bù kěyǐ yòng shǒujī.) — « En classe, les élèves ne peuvent pas utiliser le téléphone. » (complément de lieu en tête, puis Sujet + 不可以 + verbe)
\item \zh{我们的新学校很大。}(Wǒmen de xīn xuéxiào hěn dà.) — « Notre nouvelle école est grande. »
\item \zh{你今天可以去我家吗？}(Nǐ jīntiān kěyǐ qù wǒ jiā ma ?) — « Peux-tu aller chez moi aujourd'hui ? » (complément de temps \zh{今天} avant \zh{可以} ; \zh{吗} en fin de phrase)
\end{enumerate}
\end{corrige}

\begin{corrige}[Exercice 9 — Type Bac : compréhension de l'écrit et questions de langue]
\textbf{Barème indicatif (sur 20)} : A. compréhension 8 pts (2 pts par question) ; B. langue 12 pts (3 pts par question).

\textbf{A. Questions de compréhension}
\begin{enumerate}
\item C'est sa maman (sa mère) qui a acheté le nouveau téléphone de Marie (\zh{妈妈给我买了一个新手机}).
\item Elle a donné son ancien téléphone à son petit frère (\zh{我的旧手机给了弟弟}).
\item Avec son nouveau téléphone, Marie peut aller sur Internet, lire les informations et apprendre le chinois (deux réponses suffisent : \zh{上网、看新闻、学习汉语}).
\item La règle de l'école : on ne peut pas utiliser le téléphone en classe (\zh{不可以在课堂上用手机}) ; on peut l'utiliser à la maison.
\end{enumerate}

\textbf{B. Questions de langue}
\begin{enumerate}
\item Avec \zh{可以} : \zh{我可以用它上网。} — « Je peux l'utiliser pour aller sur Internet » ; \zh{也可以学习汉语} — « je peux aussi apprendre le chinois » (ou \zh{手机可以在家里用} — « on peut utiliser le téléphone à la maison »). Avec \zh{不可以} : \zh{我们不可以在课堂上用手机。} — « Nous ne pouvons pas utiliser le téléphone en classe. »
\item Dans \zh{我可以用它上网}, le verbe principal est \zh{上网} (avec \zh{用} « utiliser » comme premier verbe de la série). Le verbe modal \zh{可以} se place \textbf{avant} le (premier) verbe, juste après le sujet : Sujet + 可以 + verbe.
\item \zh{新} : \zh{新手机} « nouveau téléphone » ; \zh{旧} : \zh{旧手机} « ancien téléphone ». Ce sont des antonymes : \zh{新} qualifie ce qui est récent, neuf ; \zh{旧} qualifie ce qui est ancien, usagé (pour les choses).
\item \zh{我可以用它上网吗？}(Wǒ kěyǐ yòng tā shàngwǎng ma ?) — « Est-ce que je peux l'utiliser pour aller sur Internet ? » (on ajoute simplement \zh{吗} en fin de phrase).
\end{enumerate}
Points de vigilance : en B2, ne pas dire que \zh{可以} est le verbe principal ; en B4, ne pas inverser l'ordre des mots comme en français.
\end{corrige}

\begin{corrige}[Exercice 10 — Type Bac : thème et version]
\textbf{Barème indicatif (sur 20)} : thème 12 pts (2 à 3 pts par phrase : caractères 1 pt, grammaire 1 pt, pinyin 0,5 pt bonus) ; version 8 pts.

\textbf{A. Thème — corrigé}
\begin{enumerate}
\item \zh{学生不可以在课堂上用他们的新手机。}(Xuésheng bù kěyǐ zài kètáng shang yòng tāmen de xīn shǒujī.)
\item \zh{我的旧电脑很旧，但是我还可以用它。}(Wǒ de jiù diànnǎo hěn jiù, dànshì wǒ hái kěyǐ yòng tā.)
\item \zh{我可以用你的新手机上网吗？}(Wǒ kěyǐ yòng nǐ de xīn shǒujī shàngwǎng ma ?)
\item \zh{这是我的汉语老师的新书。}(Zhè shì wǒ de Hànyǔ lǎoshī de xīn shū.)
\item \zh{在家里，我们可以看电视；在学校，我们不可以。}(Zài jiā lǐ, wǒmen kěyǐ kàn diànshì ; zài xuéxiào, wǒmen bù kěyǐ.)
\end{enumerate}
Points de vigilance : phrase 1, le lieu \zh{在课堂上} se place avant le verbe ; phrase 2, reprendre l'objet par \zh{它} ; phrase 5, \zh{不可以} peut s'employer seul en réponse elliptique, le verbe étant sous-entendu.

\textbf{B. Version — traduction}
« Ceci est mon ancien téléphone ; je peux utiliser le nouveau téléphone pour aller sur Internet. Le nouveau téléphone, c'est maman qui me l'a acheté. À l'école, nous ne pouvons pas utiliser le téléphone, mais à la maison, nous le pouvons. »

Points de vigilance : \zh{是妈妈给我买的} = tournure d'identification « c'est maman qui me l'a acheté » ; \zh{在家可以} = « à la maison, on peut (le faire) », le verbe est sous-entendu.
\end{corrige}

\begin{corrige}[Exercice 11 — Type Bac : correction de texte et expression écrite]
\textbf{Barème indicatif (sur 20)} : A. correction 10 pts (2 pts par erreur : 1 pt identification, 1 pt justification) ; B. expression écrite 10 pts (respect des consignes 4 pts, grammaire 3 pts, vocabulaire et caractères 2 pts, cohérence 1 pt).

\textbf{A. Correction de texte — les 5 erreurs}
\begin{enumerate}
\item \zh{我用可以新手机。} → \zh{我可以用新手机。} Erreur : ordre des mots ; le modal \zh{可以} précède le verbe \zh{用} (Sujet + 可以 + verbe).
\item \zh{这是我的旧的书。} → \zh{这是我的旧书。} Erreur : double \zh{的} ; après un adjectif monosyllabique comme \zh{旧}, on n'ajoute pas \zh{的} devant le nom : 我的 + 旧书.
\item \zh{学生没可以用手机。} → \zh{学生不可以用手机。} Erreur : négation incorrecte ; \zh{可以} se nie avec \zh{不}, jamais avec \zh{没}.
\item \zh{我旧电脑还可以上网。} → \zh{我的旧电脑还可以上网。} Erreur : mot-outil \zh{的} oublié entre le possessif \zh{我} et le nom.
\item \zh{老师的新很大学校。} → \zh{老师的新学校很大。} Erreur : ordre des mots ; l'adjectif \zh{新} qualifie \zh{学校} (新学校), et le prédicat \zh{很大} se place après le nom.
\end{enumerate}

\textbf{B. Expression écrite — proposition de rédaction modèle}

\zh{我有两个手机。旧手机是黑色的，它很旧，但是还可以用。新手机是妈妈给我买的，我可以用它上网、看新闻，也可以学习汉语。在学校，我们不可以用手机。我很喜欢我的新手机！}\\[4pt]
Wǒ yǒu liǎng ge shǒujī. Jiù shǒujī shì hēisè de, tā hěn jiù, dànshì hái kěyǐ yòng. Xīn shǒujī shì māma gěi wǒ mǎi de, wǒ kěyǐ yòng tā shàngwǎng, kàn xīnwén, yě kěyǐ xuéxí Hànyǔ. Zài xuéxiào, wǒmen bù kěyǐ yòng shǒujī. Wǒ hěn xǐhuan wǒ de xīn shǒujī !\\[4pt]
« J'ai deux téléphones. L'ancien est noir, il est très vieux, mais on peut encore l'utiliser. Le nouveau, c'est maman qui me l'a acheté ; je peux l'utiliser pour aller sur Internet, lire les informations, et aussi apprendre le chinois. À l'école, nous ne pouvons pas utiliser le téléphone. J'aime beaucoup mon nouveau téléphone ! »

Vérification des consignes : \zh{可以}/\zh{不可以} employé 4 fois ; \zh{旧} et \zh{新} employés plusieurs fois ; structure \zh{用 + objet + verbe} présente (\zh{用它上网}) ; environ 75 caractères.

Points de vigilance : bien placer \zh{可以} avant le verbe ; négation \zh{不可以} ; ne pas oublier \zh{的} après le possessif ; compter les caractères pour rester entre 60 et 80.
\end{corrige}

\newpage

\coursec{Fiche bilan}

\begin{popfigure}
\begin{tikzpicture}[mindmap, grow cyclic, every node/.style={concept, font=\small}, concept color=popblue!40, level 1/.append style={level distance=3.2cm, sibling angle=60}, level 2/.append style={level distance=2.2cm}]
\node[concept color=popblue!60, font=\small\bfseries, align=center]{Leçon 36\\ 新 / 旧 et 可以}
  child[concept color=popgreen!50] { node[align=center]{Antonymes\\ 新 xīn / 旧 jiù}
    child { node {nouveau / vieux} }
    child { node[align=center]{adj. ou préfixe\\ 新书 / 旧手机} }
  }
  child[concept color=poporange!50] { node[align=center]{Verbe modal\\ 可以 kěyǐ}
    child { node[align=center]{permission\\ « pouvoir »} }
    child { node[align=center]{possibilité\\ « il est possible »} }
  }
  child[concept color=poppurple!50] { node[align=center]{Structure\\ S + 可以 + V}
    child { node[align=center]{Négation :\\ 不可以 / 不能} }
    child { node[align=center]{Question :\\ 可以 … 吗 ?} }
  }
  child[concept color=poppink!50] { node[align=center]{Comparaison\\ 会 / 能 / 可以}
    child { node {会 = savoir-faire} }
    child { node {能 = capacité} }
  }
  child[concept color=popgold!60] { node[align=center]{Pièges\\ des francophones}
    child { node {可以 avant le verbe} }
    child { node {pas de 是 + adj.} }
  }
  child[concept color=popblue!30] { node[align=center]{Médias\\ et téléphone}
    child { node {新手机 / 旧电脑} }
  };
\end{tikzpicture}
\legende{Carte mentale de la leçon : les antonymes 新 / 旧 et le verbe modal 可以.}
\end{popfigure}

\begin{bilan}
\textbf{1. Lexique clé à connaître par cœur}
\begin{center}
\begin{tabularx}{\linewidth}{|X|X|X|X|}
\hline
\rowcolor{popblueL} Caractères & Pinyin & Nature & Traduction \\
\hline
新 & xīn & adjectif & nouveau, neuf \\
\hline
旧 & jiù & adjectif & vieux, ancien (objet) \\
\hline
可以 & kěyǐ & verbe modal & pouvoir (permission, possibilité) \\
\hline
新书 & xīn shū & nom & livre neuf \\
\hline
旧手机 & jiù shǒujī & nom & vieux téléphone portable \\
\hline
新电脑 & xīn diànnǎo & nom & ordinateur neuf \\
\hline
旧报纸 & jiù bàozhǐ & nom & vieux journal \\
\hline
新同学 & xīn tóngxué & nom & nouveau camarade \\
\hline
\end{tabularx}
\end{center}

\textbf{2. Règles de grammaire à connaître par cœur}
\begin{center}
\begin{tabularx}{\linewidth}{|X|X|X|}
\hline
\rowcolor{popblueL} Structure & Exemple (caractères + pinyin) & Traduction \\
\hline
新 / 旧 + nom (sans 的) & \zh{旧手机} jiù shǒujī & vieux téléphone \\
\hline
Sujet + 可以 + verbe & \zh{我可以用手机。} Wǒ kěyǐ yòng shǒujī. & Je peux utiliser le téléphone. \\
\hline
Sujet + 不可以 + verbe & \zh{上课不可以打电话。} Shàngkè bù kěyǐ dǎ diànhuà. & En cours, on ne peut pas téléphoner. \\
\hline
可以 + verbe + 吗 ? & \zh{我可以进来吗？} Wǒ kěyǐ jìnlái ma ? & Est-ce que je peux entrer ? \\
\hline
Réponse courte & \zh{可以。/ 不可以。} Kěyǐ. / Bù kěyǐ. & Oui, tu peux. / Non. \\
\hline
\end{tabularx}
\end{center}

\textbf{3. Distinction des modaux (à ne pas confondre)}
\begin{itemize}
\item \cle{会} huì : savoir-faire appris. \zh{我会说汉语。} Wǒ huì shuō Hànyǔ. « Je sais parler chinois. »
\item \cle{能} néng : capacité physique ou circonstancielle. \zh{我今天能来。} Wǒ jīntiān néng lái. « Je peux venir aujourd'hui. »
\item \cle{可以} kěyǐ : permission ou autorisation. \zh{你可以看我的新手机。} Nǐ kěyǐ kàn wǒ de xīn shǒujī. « Tu peux regarder mon nouveau téléphone. »
\end{itemize}

\textbf{4. Méthodes express}
\begin{itemize}
\item Pour demander la permission : \cle{我可以 + verbe + 吗 ?} ; répondre par \zh{可以。} ou \zh{不可以。} (jamais par \zh{是}).
\item Pour opposer deux objets : \cle{新} (neuf) contre \cle{旧} (usagé) ; on les place directement devant le nom : \zh{新电脑} / \zh{旧电脑}.
\item Pour interdire : \cle{不可以 + verbe} ou \cle{不能 + verbe} : \zh{这里不能上网。} Zhèlǐ bù néng shàngwǎng. « Ici, on ne peut pas aller sur Internet. »
\item Attention : \zh{旧} s'emploie pour les objets ; pour les personnes âgées, on dit \zh{老} lǎo.
\end{itemize}
\end{bilan}

\begin{aretenir}[Checklist avant le Bac]
\begin{itemize}
\item[$\square$] Je connais les antonymes \zh{新} xīn « nouveau » et \zh{旧} jiù « vieux » et leur place devant le nom.
\item[$\square$] Je sais que \zh{新} et \zh{旧} se placent directement devant le nom, sans \zh{的} dans les composés courants (\zh{新书}, \zh{旧手机}).
\item[$\square$] Je sais construire la structure affirmative : Sujet + \zh{可以} + verbe (+ complément).
\item[$\square$] Je sais construire la négation : Sujet + \zh{不可以} / \zh{不能} + verbe.
\item[$\square$] Je sais poser une question de permission avec \zh{可以 … 吗？} et y répondre correctement.
\item[$\square$] Je distingue \zh{会} (savoir-faire), \zh{能} (capacité) et \zh{可以} (permission).
\item[$\square$] Je place toujours le verbe modal \zh{可以} AVANT le verbe principal, jamais après.
\item[$\square$] Je n'utilise pas \zh{旧} pour une personne âgée : je dis \zh{老}.
\item[$\square$] Je sais écrire les caractères \zh{新} et \zh{旧} dans le bon ordre des traits (gauche $\rightarrow$ droite).
\item[$\square$] Je sais employer ces structures dans le thème des médias : \zh{新手机}, \zh{旧电脑}, \zh{可以上网}.
\item[$\square$] Je réponds à une question avec \zh{可以} par \zh{可以} ou \zh{不可以}, jamais par \zh{是} ou \zh{对}.
\item[$\square$] Je fais attention aux tons : xīn (1\ier{} ton), jiù (4\ieme{} ton), kěyǐ (3\ieme{} + 3\ieme{} ton).
\end{itemize}
\end{aretenir}

\findecours

\end{document}
